\part{The Eastern Realms}
% After \part{The Eastern Realms}
\begin{quote}
  \small
  \textit{``The maps of the empire draw the East as a single sweep -- Sihai, Nihori, Ayokha -- all pressed together like tiles on a game board. The roads know better. Sihai measures time in dynasties; Nihori counts it in storms. Ayokha negotiates between them with the patience of the monsoon. And beneath them all, the older peoples keep their own thresholds, their own debts, their own ways of not being seen.''}
  \noindent --- Dusana of the Raven Road, at the edge of the Inland Sea
  \end{quote}
% ------------------------------
% CHAPTER: SIHAI
% ------------------------------
\chapter{Sihai}
\emph{The Central Kingdom.}

\noindent
{\Large
\textit{The Mandate of Heaven is not a right. It is a loan. The emperor borrows legitimacy from the ancestors, from the land, from the Hollow itself. When the loan is due, the empire falls. I do not serve the emperor. I serve the ledger that tracks the interest.}

\vspace{0.5em}

I am a Provincial Censor. My brush is heavier than any sword. A single stroke can erase a name, a family, a dynasty. The emperor reads my reports and trembles. He knows that I have written nine reports this year. The ninth is about him.

\vspace{0.5em}

Sihai is the central kingdom, the seat of the Emperor and the Mandate of Heaven. Its canals carry grain from the south; its examination halls produce bureaucrats who govern by the brush. Power here is measured in ink, not iron. The Censorate watches everyone; the Salt Monopoly controls trade; the Lethai-thora keep the secrets that the Mandate prefers to forget. But here's the secret the Emperor doesn't want you to know: the Mandate is not a divine right. It is a loan. The ancestors are the creditors. The Hollow collects the interest. Every dynasty that has ever fallen did so because the interest came due and the emperor could not pay.

\vspace{0.5em}

Beneath the mortal bureaucracy runs a deeper one: the Celestial Bureaucracy of gods, ghosts, and ancestors, where every official has a heavenly counterpart and every erased name becomes a hungry ghost. The Lethai-thora---the ancient archivists who taught the first emperors to write---still keep the hidden ledgers, the ones that record what the Censorate burns. A well-written couplet can appease a ghost; a forged seal can summon a demon. The Hollow's ninth debt is the unpaid interest on the Mandate itself.

\vspace{0.5em}

Never speak of succession before naming the Emperor. The Censorate has ears in every tea house. A misplaced word can erase your name from the rolls.

\vspace{1em}

\noindent\textit{--- Quiet Chancellor Wei, who has never smiled and whose seal is a laughing mouth}
}

\vspace{1em}

\begin{almanac}
\textbf{What do people eat for breakfast?} Congee with preserved vegetables---the congee is rice boiled until it is memory, the vegetables are salted and aged in clay jars. The jars are sealed with the emperor's own wax. The wax is a promise. The promise is that the vegetables will not poison you. The promise is sometimes broken.

\textbf{What does a child learn before they can walk?} To write their name. Sihai children are taught the eight strokes of the character ``person''---the ninth stroke is the one that turns a name into a signature. A signature is a contract. A contract is a debt. A debt is a life.

\textbf{What is the first thing you notice about a stranger?} Their chop---the seal they carry to sign documents. A stranger whose chop is jade is a noble. A stranger whose chop is brass is a clerk. A stranger whose chop is wood is a merchant. A stranger with no chop is a ghost. Ghosts cannot sign contracts. Ghosts cannot be trusted.

\textbf{What is the most common crime?} Examination fraud---but fraud is not a crime, it is a family business. The examinations determine who governs. The questions are leaked in advance. The answers are memorized. The examiners are bribed. The system is a sieve. The sieve is called ``merit.''

\textbf{What does no one talk about but everyone knows?} The Emperor is a figurehead. The real power is the Censorate---nine ministers who have never been seen together. Their names are not recorded. Their faces are not painted. Their orders are whispered. The Emperor signs what they tell him to sign. The Emperor has not refused in forty years.

\textbf{Where do people go when they die?} To the ancestral tablets. The dead are not buried. Their names are written on wooden tablets, placed in the family shrine. The tablets are fed with rice and tea. The dead eat the steam. The dead are never full.

\textbf{How do you apologize for a serious wrong?} You offer a calligraphy---a scroll written in your own blood. The scroll is presented to the wronged party's ancestral shrine. If the ancestors accept, you are forgiven. If the scroll bursts into flame, you are not. If the scroll writes itself, the ancestors are angry. The ancestors' anger is a famine.

\textbf{What is the most important festival?} The Festival of the Ancestors (autumn). The tablets are cleaned and polished. The dead are invited to feast. The feast is laid out on the family altar. The food is removed after an hour. The dead have eaten the essence. The essence is memory.

\textbf{What is the secret that could destroy a powerful person?} The Censorate's ninth minister is a woman---the emperor's own mother, who was supposed to have died in childbirth. She faked her death to escape the politics of the court. She has been running the Censorate from a hidden villa in the western hills. The emperor does not know. The emperor's mother signs his death warrants. He signs them too.
\end{almanac}

\newpage
\section{Sihai --- ``The Central Kingdom''}
\label{chap:sihai}

\begin{quote}
  \textbf{Quiet Chancellor (Elite):}  
  ``The Mandate of Heaven is not a right. It is a loan. The emperor borrows legitimacy from the ancestors, from the land, from the Hollow itself. When the loan is due, the empire falls. I do not serve the emperor. I serve the ledger that tracks the interest.''
\end{quote}

\begin{quote}
  \textbf{Provincial Censor (Commoner):}  
  ``My brush is heavier than any sword. A single stroke can erase a name, a family, a dynasty. The emperor reads my reports and trembles. He knows that I have written nine reports this year. The ninth is about him.''
\end{quote}

\begin{tcolorbox}[colback=black!3,colframe=blue!60!black,title={Theme \& Atmosphere}]
Sihai is the central kingdom, the seat of the Emperor and the Mandate of Heaven. Its canals carry grain from the south; its examination halls produce bureaucrats who govern by the brush. Power here is measured in ink, not iron. The Censorate watches everyone; the Salt Monopoly controls trade; the Lethai-thora keep the secrets that the Mandate prefers to forget. Sihai is not a land of warriors — it is a land of scholars, clerks, and spies. A well-written couplet can end a feud; a forged seal can start a war. The Hollow's ninth debt is the unpaid interest on the Mandate itself. Every dynasty borrows legitimacy; every silence is a debt.

\vspace{2mm}
\textbf{What you notice first:} The smell of ink and canal water. The way officials stamp every document with a brass seal. The constant ticking of water-clocks in public squares.

\textbf{The one rule that kills newcomers:} Never speak of succession before naming the Emperor. The Censorate has ears in every tea house. A misplaced word can erase your name from the rolls.
\end{tcolorbox}

\subsection*{Regional Mood: The Weight of the Brush}

Power here is measured in examination scores, salt licenses, and the thickness of a family's archive. The Emperor rules, but the Censorate audits, the Salt Monopoly finances, and the Lethai-thora remember what the Censorate burns. A provincial governor's authority rests on his examination rank; a merchant's on his salt license. The Hollow watches from the missing ninth step, the unpoured ninth cup, the unsaid ninth name.

\textbf{Starting Location:} An examination hall where ink-stained desks face a silent proctor. A Provincial Censor holds up a scroll: ``This candidate's family name has been erased. He does not know. You will tell him — or you will take his place in the ledger.''

\subsection*{Prominent NPCs of Sihai}
\label{sec:sihai-npcs}

\begin{longtable}{|p{0.2\textwidth}|p{0.25\textwidth}|p{0.3\textwidth}|p{0.15\textwidth}|}
\hline
\textbf{Name} & \textbf{Role/Title} & \textbf{Motivation} & \textbf{Rumoured Quirk} \\
\hline
Quiet Chancellor Wei & Imperial censor & Root out corruption before the Mandate fails & He has never smiled. His seal is a laughing mouth. \\
\hline
Salt Commissioner Li & Head of the Salt Monopoly & Control the price of salt and the flow of silver & She tastes every batch of salt before it is stamped. Her tongue is white. \\
\hline
Provincial Censor Chen & Overseer of the eastern prefectures & Erase a rebel lineage from the rolls & He writes with his left hand. His right hand is for swearing oaths. \\
\hline
Lethai-thora Keeper Lian & Hidden archivist & Protect a scroll that proves the emperor's claim is false & She has not spoken aloud in twenty years. She writes on water. \\
\hline
The Hollow's Debt-Walker & Spirit of unpaid Mandate & Collect the ninth name, the ninth step, the ninth cup & It appears as a scribe with too many fingers. It never blinks. \\
\hline
\end{longtable}

\subsection*{Names of Sihai}
\label{sec:sihai-names}

Inspired by the central kingdom: short, tonal, with honorifics. Surnames often denote a clan or a prefecture.

\begin{longtable}{|p{0.25\textwidth}|p{0.25\textwidth}|p{0.2\textwidth}|p{0.2\textwidth}|}
\hline
\textbf{First Names (Male)} & \textbf{First Names (Female)} & \textbf{Clan/House} & \textbf{Epithets} \\
\hline
Wei & Lian & Chen & \textit{the Quiet} \\
\hline
Li & Mei & Zhou & \textit{Salt-Tongue} \\
\hline
Chen & Huan & Zhao & \textit{the Unerased} \\
\hline
Zhang & Ning & Wu & \textit{Ninth Seal} \\
\hline
Liu & Ping & Tang & \textit{the Brush} \\
\hline
\end{longtable}

\paragraph*{(Hall/Canal/Archive)} Examination hall with ink-stained desks; Grand Canal lock with barges queueing; Lethai-thora hidden archive behind a waterfall.

\section*{Spades — Places}
\label{sec:sihai-places}
\begin{enumerate}
\setcounter{enumi}{1}
\item \textbf{Examination Hall} --- A long building with ink-stained desks, silent proctors, and a single bell.
  \hfill \textit{The bell rings when a candidate cheats. It rings once a day on average.}
\item \textbf{Grand Canal Lock} --- A stone gate where barges queue for inspection. Officials stamp manifests.
  \hfill \textit{The stamp leaves a watermark that glows under moonlight. Forgeries are obvious.}
\item \textbf{Lethai-thora Hidden Archive} --- A pavilion behind a waterfall, filled with forbidden scrolls.
  \hfill \textit{The scrolls are sealed with silk cords. The cords have nine knots. The ninth is a lie.}
\item \textbf{Salt Commissioner's Counting House} --- A fortress of ledgers, guarded by armed clerks.
  \hfill \textit{The ledgers are written in disappearing ink. Only the Commissioner knows the reagent.}
\item \textbf{Provincial Censorate Office} --- A grey building with a brass scale above the door.
  \hfill \textit{The scale weighs truth and lies. It has never balanced.}
\item \textbf{Emperor's Forbidden Garden} --- A walled garden where no one may enter without permission.
  \hfill \textit{The garden contains a single tree with nine branches. The ninth branch is dead.}
\item \textbf{Well of Erased Names} --- A stone well where the Censorate throws scrolls of erased lineages.
  \hfill \textit{The scrolls burn as they fall. The smoke smells of ink and old blood.}
\item \textbf{Tea House of the Ninth Cup} --- A popular establishment where the ninth cup is never poured.
  \hfill \textit{The owner has a bell that rings when someone pours the ninth cup. It rings often.}
\item \textbf{Canal Mother's Court} --- A shaded bench where water rights are assigned.
  \hfill \textit{The bench is carved with the names of every canal mother for nine generations.}
\item[J] \textbf{The Emperor's Tomb (Unfinished)} --- A subterranean complex where the current emperor's tomb is being built.
  \hfill \textit{The workers are slaves. They will be buried with him. They know.}
\item[Q] \textbf{The Hollow's Gate (Subterranean)} --- A crack in the floor of the Forbidden Garden.
  \hfill \textit{The crack is warm. It whispers the names of emperors who died without heirs.}
\item[K] \textbf{The Ninth Ledger's Vault} --- A vault beneath the Counting House, accessible only by nine keys.
  \hfill \textit{The ninth key is held by the Hollow. The Hollow will lend it — for a price.}
\item[A] \textbf{The First Examination Hall (Myth)} --- Where the first emperor tested his sons.
  \hfill \textit{The hall is invisible. Only those who have passed nine examinations can see it.}
\end{enumerate}

\paragraph*{(Censor/Merchant/Keeper)} Provincial censor with a list of forbidden words; salt merchant with honest scales except when they aren't; Lethai-thora memory-keeper with a silence-token.

\section*{Hearts — People & Factions}
\label{sec:sihai-people}
\begin{enumerate}
\setcounter{enumi}{1}
\item \textbf{Provincial Censor} --- Grey-robed, a brass seal on a chain. He carries a list of forbidden words.
  \hfill \textit{``Speak the ninth word,'' he says, ``and your name will join the list.''}
\item \textbf{Salt Merchant} --- A thin woman with a brass scale. Her weights are hollow.
  \hfill \textit{``The scale is honest,'' she says. ``My hand is not.''}
\item \textbf{Lethai-thora Keeper} --- Silent, carrying a brass lamp. She writes on water with a reed brush.
  \hfill \textit{``The words vanish,'' she says. ``That is how I know they are true.''}
\item \textbf{Canal Mother} --- A matriarch in blue robes, a key to the sluice gate on her belt.
  \hfill \textit{``Water is memory,'' she says. ``The canal remembers every bribe.''}
\item \textbf{Examination Candidate} --- A young man with ink-stained fingers. He has failed nine times.
  \hfill \textit{``The tenth will be different,'' he says. He has said this nine times.}
\item \textbf{Quiet Chancellor's Spy} --- A woman disguised as a tea-seller. She listens at every table.
  \hfill \textit{She writes reports on her own skin. The ink is invisible.}
\item \textbf{Salt Worker (Slave)} --- A man with cracked hands, harvesting salt from pans.
  \hfill \textit{``The salt tastes of tears,'' he says. ``Whose tears, I do not know.''}
\item \textbf{Hollow's Debt-Walker} --- A tall figure in grey, with too many fingers. It carries a ledger.
  \hfill \textit{It asks: ``Where is the ninth name?'' It will not leave without an answer.}
\item \textbf{Emperor's Dreamer} --- A woman locked in a tower. She dreams the emperor's nightmares.
  \hfill \textit{She wakes screaming. The dreams are always about the Hollow.}
\item[J] \textbf{The First Censor (Ghost)} --- A skeleton in grey robes, sitting in the Well of Erased Names.
  \hfill \textit{He asks: ``Did I erase the right name?'' The answer is always a lie.}
\item[Q] \textbf{Canal Mother's Ghost} --- A translucent figure walking the canal at midnight.
  \hfill \textit{She checks the water level. When it drops, a debt is overdue.}
\item[K] \textbf{The Ninth Emperor's Ghost} --- A crowned skeleton in the unfinished tomb.
  \hfill \textit{He asks: ``Is the Mandate renewed?'' If you say no, he weeps.}
\item[A] \textbf{Kuva's Shadow (Hollow)} --- The vast darkness that covers the Forbidden Garden at midnight.
  \hfill \textit{It asks: ``Who holds the ninth key?'' The answer is a name no one knows.}
\end{enumerate}

\paragraph*{(Forgery/Monsoon/Ancestor)} A scroll is found to be a forgery — the party is blamed; monsoon breaks early — the canal closes; an ancestor demands a public accounting of a forgotten debt.

\section*{Clubs — Complications/Threats}
\label{sec:sihai-complications}
\begin{enumerate}
\setcounter{enumi}{1}
\item \textbf{Forged Scroll} --- A tax ledger is discovered to be a fake. The forgery is perfect.
  \hfill \textit{The forger is a Censorate clerk. The scroll is not a forgery — it is the original.}
\item \textbf{Monsoon Breaks Early} --- The canal closes for inspection. Delay 2 segments.
  \hfill \textit{The closure is a cover. The Censorate is searching for a specific barge.}
\item \textbf{Ancestor's Demand} --- A ghost appears at the family shrine. It demands a public accounting.
  \hfill \textit{The debt is from nine generations ago. The ancestor will not leave until it is paid.}
\item \textbf{Censorate Audit} --- Officials arrive with a warrant. They have a list of forbidden words.
  \hfill \textit{The list includes a word you spoke yesterday. The penalty is erasure.}
\item \textbf{Salt Monopoly Price Hike} --- The Commissioner doubles the salt tax. Riots begin.
  \hfill \textit{The hike is a bribe. The bribe is to the Quiet Chancellor. He accepted.}
\item \textbf{Canal Lock Sabotage} --- Someone drops a barge at the lock. The canal backs up.
  \hfill \textit{The saboteur is a Canal Mother's son. He wants to flood the Censorate office.}
\item \textbf{Examination Cheating Scandal} --- A candidate is caught with a cheat sheet. The sheet has your name.
  \hfill \textit{The cheat sheet was planted. The real target is your patron.}
\item \textbf{Lethai-thora Archive Fire} --- A lamp falls in the hidden pavilion. Scrolls burn.
  \hfill \textit{The fire was set by the Censorate. They want to erase a specific history.}
\item \textbf{Hollow's Debt Called} --- A Debt-Walker appears at the tea house. It points at a party member.
  \hfill \textit{The member owes a debt. The debt is the ninth cup. Pay in story or memory.}
\item[J] \textbf{Emperor's Health Fails} --- The Mandate wavers. Succession rumours spike.
  \hfill \textit{The emperor has no heir. The Hollow is watching. It will collect the ninth name.}
\item[Q] \textbf{Forbidden Garden Breached} --- Someone enters the garden without permission. The Hollow follows.
  \hfill \textit{The intruder is a child. The Hollow asks: ``Why are you here?'' The child does not know.}
\item[K] \textbf{Ninth Ledger Stolen} --- The vault is cracked. The ninth ledger is missing.
  \hfill \textit{The ledger contains the names of everyone the Censorate has erased. The thief is a ghost.}
\item[A] \textbf{Mandate Falls} --- The emperor dies without an heir. The Hollow claims the throne.
  \hfill \textit{The Hollow sits on the dragon throne. It asks: ``Who will rule?'' Silence is the only answer.}
\end{enumerate}

\paragraph*{(Pass/Seal/Token)} Stamped pass from a provincial governor (one use); copy of the Celestial Ledger (reveals one hidden tax); Lethai-thora silence-token (one question the keeper must answer).

\section*{Diamonds — Rewards/Leverage}
\label{sec:sihai-rewards}
\begin{enumerate}
\setcounter{enumi}{1}
\item \textbf{Stamped Pass (Governor's Seal)} --- A wax-sealed document. Grants passage through any provincial checkpoint.
  \hfill \textit{The seal is a tiger. It glows when a Censorate agent is nearby.}
\item \textbf{Celestial Ledger (Copy)} --- A scroll containing one hidden tax. Use it to blackmail an official.
  \hfill \textit{The ledger is written in disappearing ink. Read it aloud before it fades.}
\item \textbf{Lethai-thora Silence-Token} --- A brass charm. Present it to demand a keeper answer one question.
  \hfill \textit{The keeper cannot lie. The question must be about a forgotten secret.}
\item \textbf{Junk Seal (Official)} --- A brass seal. Impersonate a provincial official for one scene.
  \hfill \textit{The seal is warm. It glows when someone tests its authenticity.}
\item \textbf{Canal Mother's Key} --- A bronze key. Open any sluice gate in the Grand Canal.
  \hfill \textit{The key is cold. It warms when a water right is being violated.}
\item \textbf{Salt Commissioner's Waiver} --- A stamped document. Exempts you from the salt tax for one year.
  \hfill \textit{The waiver is written on silk. The silk was stolen from the Commissioner's own robes.}
\item \textbf{Examination Answer Key (Forged)} --- A cheat sheet that actually works. Pass one examination automatically.
  \hfill \textit{The answers are written in invisible ink. The reagent is moonlight.}
\item \textbf{Quiet Chancellor's Seal (Jade)} --- A jade seal. Commands the Censorate to ignore one crime.
  \hfill \textit{The seal is a forgery. The Chancellor made it himself. It still works.}
\item \textbf{Forbidden Garden Key (Iron)} --- A heavy iron key. Opens the gate to the Emperor's garden.
  \hfill \textit{The key is cursed. Each use ages the wielder one month.}
\item[J] \textbf{Hollow's Ninth Ledger Page} --- A single page from the ninth ledger. Contains nine erased names.
  \hfill \textit{Speak a name, and that person's deeds become known. The Hollow will remember too.}
\item[Q] \textbf{First Censor's Brush} --- A brush made from a phoenix feather. Write any name, and it is erased from all records.
  \hfill \textit{The brush works once. Then it crumbles. The Hollow will know.}
\item[K] \textbf{Canal Mother's Ghostly Dipper} --- A clay dipper used by a dead Canal Mother. Use it to control any sluice gate.
  \hfill \textit{The dipper is cold. It weeps water when a debt is overdue.}
\item[A] \textbf{The Mandate Itself (Scroll)} --- A scroll that proves the current emperor has no legitimate claim.
  \hfill \textit{Read it aloud, and the Mandate falls. The Hollow will take the throne.}
\end{enumerate}

\section*{Quick use notes}
\label{sec:sihai-quick-use}
\begin{itemize}
\item Draw until you have all four suits: Spade = place, Heart = actor, Club = pressure, Diamond = leverage. Highest rank sets the main Clock (2–5 → 4, 6–10 → 6, J/Q/K → 8, A → 10).
\item Diamonds are codified outcomes (passes, seals, tokens) that change position rather than call for a roll.
\item If any A appears, echo \textbf{ink-omens}—scrolls that rustle in still air, a lamp that burns green, a forgotten name whispered from a shelf.
\end{itemize}

\subsection*{Additional Features}
\label{sec:sihai-features}

\begin{itemize}
\item \textbf{The Mandate Clock [8]:} Replaces the usual Mandate/Crisis dials in Sihai courts. When it fills, the current dynasty loses the Mandate; chaos spreads, but the Hollow's hunger is briefly sated. The clock ticks when the emperor's health fails, when a succession dispute arises, or when a major corruption scandal is exposed.
\item \textbf{Censorate Audits:} The Censorate can audit any official, any merchant, any traveller. An audit is a social combat (Presence + Sway vs. Censor's Lore). Lose, and your name is added to a watchlist — or erased.
\item \textbf{Water Rights as Currency:} In the canal system, water rights are traded like silver. A Canal Mother's key is worth more than a chest of coins. Forging a water right is a capital crime.
\end{itemize}

\subsection*{Lore Echoes (Cross-Expansion Hooks)}
\begin{itemize}
  \item \textbf{The Ninth Ledger and the Kalrantha Amulet (Dhahara):} The ninth ledger is said to contain the name of the prince who stole the amulet. The Veiled Sisters want to erase it. The Breath-less want to find it.
  \item \textbf{Silence-Tokens and the Covenant (Mykkiel):} The Covenant recognises Lethai-thora silence-tokens as valid requests for closed hearings. A silence token can pause any court proceeding for the duration of a single breath.
  \item \textbf{Canal Mothers and the Fhara (Way of Silk):} Fhara well-judges and Sihai canal mothers use similar water-measuring systems. A canal mother's key can be used to negotiate a well contract.
  \item \textbf{The Mandate and the Hollow (Eastern Realms):} The Hollow is the unpaid interest on the Mandate. Aelaerem hedge-keepers say that when a dynasty falls, the Hollow eats the ninth name.
\end{itemize}

\subsection*{Sihai Superstitions (burn for +1 die once)}
\begin{itemize}
  \item \textbf{Bread to the Canal:} Crumble a crust into the water. Ignore the first \emph{Forged Scroll} penalty.
  \item \textbf{Knot the Scroll:} Tie a single knot in a silk cord. Cancel \emph{Censorate Audit} or take –1 die on your next forgery detection roll (your choice).
  \item \textbf{Name to the Well:} Whisper a dead censor's name into the Well of Erased Names. Gain +1 Effect when negotiating with the Censorate this scene, but mark \emph{Obligation} to the Hollow.
\end{itemize}

\begin{tcolorbox}[colback=black!3,colframe=black!40!white,title={Patronage \& Power}]
In Sihai, power flows through ink, salt, and the weight of the Mandate. The Emperor rules, but the Censorate audits, the Salt Monopoly finances, and the Lethai-thora remember. A well-written couplet can end a feud; a forged seal can start a war.

\textbf{For the GM:}  
Power in Sihai is measured in examination scores, salt licenses, and the thickness of a family's archive. Patronage often takes the form of stamped passes, ledger copies, or silence-tokens — each tied to a specific office, canal, or debt. To emphasize this:
\begin{itemize}
\item Use the \textbf{Mandate Clock} as a visible campaign front. The party's actions can tick or reset it.
\item Enforce \textbf{Censorate Audits} as a recurring threat — players must maintain clean papers or face erasure.
\item Let \textbf{water rights} be a form of currency that can be stolen, forged, or contested.
\end{itemize}
In Sihai, the brush forgets nothing — and the Hollow remembers every erased name.
\end{tcolorbox}

\newpage

% ------------------------------
% CHAPTER: NIHORI
% ------------------------------
\chapter{Nihori}
\emph{The Isles of the Dawn Spirit.}

\noindent
{\Large
\textit{I ruled the waves before the Hollow learned to walk. My empire sank because I spoke a name I should not have. Do not speak my name. Do not ask what it was. The sea remembers. That is enough.}

\vspace{0.5em}

I am a Wanderer. I carry a broken sword and a grudge. The sword was my father's. The grudge is mine. The man who killed him died nine years ago. I still walk the isles, looking for his ghost. The ghost is looking for me too.

\vspace{0.5em}

Nihori is a chain of volcanic islands, shrouded in mist and lashed by storms. Here, the old empire sank beneath the waves, and the Lethai-al guard the drowned kingdom's memory. The people are wanderers, duelists, and shrine maidens. They worship the storm spirits and fear the ninth wave. The Hollow breathes in the salt foam; the Oath-Keepers ride the tide. But here's the secret the Shogun doesn't want you to know: the old empire did not sink. It was pulled under. And whatever pulled it is still down there, waiting for the ninth wave to break the right way.

\vspace{0.5em}

The Shogun sits in a half-ruined castle, his authority challenged by three rival daimyo who sharpen their blades and their treaties. Oni cults whisper in the ash wastes, offering forgotten bargains. The Storm Emperor's ghost does not rule---but his silence haunts every island. The ninth taboo is absolute: never speak the name of the drowned capital. The sea will hear you, and it will ask why you called. And when the sea asks a question, it expects an answer---in blood, in memory, or in a name you did not mean to give.

\vspace{0.5em}

Never speak the name of the drowned capital. The sea will hear you, and it will ask why you called.

\vspace{1em}

\noindent\textit{--- The Storm Emperor's Ghost, who appears as a column of steam over the water}
}

\vspace{1em}

\begin{almanac}
\textbf{What do people eat for breakfast?} Grilled fish and pickled vegetables---the fish is caught in the caldera lakes, the vegetables are fermented in clay pots buried in volcanic ash. The ash is warm. The warmth is the mountain's breath. The mountain is sleeping. The mountain dreams of the drowned empire.

\textbf{What does a child learn before they can walk?} To read the tide. Nihori children are taught the eight wave-heights: Calm, Ripple, Swell, Storm, Tsunami, and the ninth---the wave that has no height because it is the sea itself rising. The ninth wave is the drowning. The drowning is the empire's memory.

\textbf{What is the first thing you notice about a stranger?} Their sword. Nihori warriors carry their ancestors' blades, passed down for generations. A stranger whose sword is rusted has forgotten their lineage. A stranger whose sword is sharp has remembered. A stranger with no sword is a coward. Cowards are not fed.

\textbf{What is the most common crime?} Dueling---but dueling is not a crime, it is a conversation. The dueling codes are written in verse. The verses are interpreted by the Shogun's court. The interpretations are contradictory. The contradictions are settled by dueling. The circle is beautiful.

\textbf{What does no one talk about but everyone knows?} The drowned capital is not underwater. It is beneath---beneath the caldera, beneath the bedrock, beneath the world. The sea is a lid. The lid is cracked. Something is leaking through. The leaks are the oni. The oni are the empire's regrets.

\textbf{Where do people go when they die?} To the sea cliffs. The dead are placed in caves overlooking the water. The waves sing to them. The songs are the empire's lullabies. The dead listen. The dead do not sleep.

\textbf{How do you apologize for a serious wrong?} You offer a duel---a matched contest with the wronged party's choice of weapon. If you lose, you are forgiven. If you win, you are not. If you refuse, you are exiled. Exile is the open sea. The sea is patient. The sea will wait.

\textbf{What is the most important festival?} The Calling of the Tides (spring). The shrine maidens walk into the water, singing the names of the drowned. The waves answer. The waves speak the names of the living who will die in the coming year. The names are never forgotten. The names are never repeated.

\textbf{What is the secret that could destroy a powerful person?} The Shogun is a ghost. He died in the last oni war, but his body was never found. The daimyo rule in his name because they are afraid of what will happen if they admit he is gone. The oni know. The oni are waiting for the admission. The admission will be the war.
\end{almanac}

\newpage
\section{Nihori --- ``The Isles of the Dawn Spirit''}
\label{chap:nihori}

\begin{quote}
  \textbf{Storm Emperor's Ghost (Elite):}  
  ``I ruled the waves before the Hollow learned to walk. My empire sank because I spoke a name I should not have. Do not speak my name. Do not ask what it was. The sea remembers. That is enough.''
\end{quote}

\begin{quote}
  \textbf{Wanderer (Commoner):}  
  ``I carry a broken sword and a grudge. The sword was my father's. The grudge is mine. The man who killed him died nine years ago. I still walk the isles, looking for his ghost. The ghost is looking for me too.''
\end{quote}

\begin{tcolorbox}[colback=black!3,colframe=blue!60!black,title={Theme \& Atmosphere}]
Nihori is a chain of volcanic islands, shrouded in mist and lashed by storms. Here, the old empire sank beneath the waves, and the Lethai-al guard the drowned kingdom's memory. The people are wanderers, duelists, and shrine maidens. They worship the storm spirits and fear the ninth wave. The Hollow breathes in the salt foam; the Debt-Walkers ride the tide. Nihori has no central government — only clans, pirate fleets, and the memory of a Storm Emperor who spoke a forbidden name and drowned his own capital.

\vspace{2mm}
\textbf{What you notice first:} The smell of salt and volcanic ash. The way every traveller carries a sword, even the tea-sellers. The constant sound of bells on shrine gates — ringing in the wind.

\textbf{The one rule that kills newcomers:} Never speak the name of the drowned capital. The sea will hear you, and it will ask why you called.
\end{tcolorbox}

\subsection*{Regional Mood: The Weight of the Wave}

Power here is measured in duels won, storms survived, and the length of a feud. The Storm Emperor's ghost is said to rule from a sunken throne, but the living follow clan chiefs, pirate admirals, and shrine maidens who can read the wind. A duel is not a fight — it is a negotiation. First blood decides the truth. The loser owes a story. The winner owes a cup of sake.

\textbf{Starting Location:} A bamboo grove on a clifftop, where two duelists face each other at dawn. A Shrine Maiden rings a copper bell: ``The dispute is over a tide-charm. One says it was stolen. The other says it was inherited. Watch the duel. The first blood tells the truth. Then you will be the witness.''

\subsection*{Prominent NPCs of Nihori}
\label{sec:nihori-npcs}

\begin{longtable}{|p{0.2\textwidth}|p{0.25\textwidth}|p{0.3\textwidth}|p{0.15\textwidth}|}
\hline
\textbf{Name} & \textbf{Role/Title} & \textbf{Motivation} & \textbf{Rumoured Quirk} \\
\hline
Storm Emperor's Ghost & Sunken ruler of the drowned kingdom & Prevent anyone from speaking the forbidden name & He appears as a column of steam over the water. \\
\hline
Shrine Maiden Hana & Keeper of the storm shrine & Pacify the spirit of the ninth wave & She rings a bell that can only be heard underwater. \\
\hline
Wanderer Ronin & Masterless swordsman & Find the ghost of his father's killer & He has never lost a duel. He has also never drawn first blood. \\
\hline
Pirate Captain Red-Hook & Kalingan raider & Control the strait between the isles & His left hand is a brass hook. The hook is also a key. \\
\hline
The Drowned Scholar (Ghost) & A Lethai-al who died before completing his research & Ask the living to retrieve a scroll from the sunken capital & He appears at low tide, holding a skeletal oar. \\
\hline
\end{longtable}

\subsection*{Names of Nihori}
\label{sec:nihori-names}

Inspired by the northern isles: short, with vowel endings and occasional glottal stops. Surnames denote a clan, a wave, or a drowned ancestor.

\begin{longtable}{|p{0.25\textwidth}|p{0.25\textwidth}|p{0.2\textwidth}|p{0.2\textwidth}|}
\hline
\textbf{First Names (Male)} & \textbf{First Names (Female)} & \textbf{Clan/Epithet} & \textbf{Wave-Name} \\
\hline
Kenji & Hana & Wave-walker & \textit{Ninth Breath} \\
\hline
Ronan & Yuki & Storm-kin & \textit{Salt-Tongue} \\
\hline
Kael & Mira & Drowned-line & \textit{Unforgotten} \\
\hline
Takeshi & Sora & Foam-born & \textit{Tide-Charm} \\
\hline
Jiro & Nami & Ancestor-salt & \textit{the Unshadowed} \\
\hline
\end{longtable}

\paragraph*{(Forge/Grove/Shrine)} A swordsmith's forge with a waterfall cooling tank; a bamboo grove where duels are fought at dawn; a storm-shrine on a cliff with ropes for climbing.

\section*{Spades — Places}
\label{sec:nihori-places}
\begin{enumerate}
\setcounter{enumi}{1}
\item \textbf{Swordsmith's Forge} --- A waterfall cooling tank; prayer strips on the bellows.
  \hfill \textit{The swordsmith prays before each strike. The prayer is a verse from the old epic.}
\item \textbf{Bamboo Grove (Duel Ground)} --- A clearing where the stalks are scarred by blade cuts.
  \hfill \textit{The bamboo bleeds sap when a duel is dishonourable. The sap is red.}
\item \textbf{Storm-Shrine (Cliff)} --- A wooden platform on a sea cliff, with ropes for climbing.
  \hfill \textit{The ropes are braided with human hair. The hair is from sailors who survived shipwrecks.}
\item \textbf{Drowned Palace (Ruins)} --- Half-submerged columns visible at low tide.
  \hfill \textit{The columns are carved with names. The names are the drowned kingdom's royalty.}
\item \textbf{Volcano Shrine (Dormant)} --- A crater where the forbidden name was spoken.
  \hfill \textit{The crater is warm. Speaking any name aloud here causes the ground to tremble.}
\item \textbf{Sake House of the Ninth Cup} --- A tavern where duels are negotiated before they are fought.
  \hfill \textit{The ninth cup is never poured. It is for the dead.}
\item \textbf{Pirate Cove (Hidden)} --- A sheltered bay behind a reef, accessible only by a narrow channel.
  \hfill \textit{The channel is marked by a single red buoy. The buoy is a skull.}
\item \textbf{Tide-Charm Vault (Sea Cave)} --- A cave where Lethai-al store their ancestors' charms.
  \hfill \textit{The charms are bones, shells, and pieces of shipwreck. Each hums with a different memory.}
\item \textbf{Ancestor's Salt Pan} --- A flat where seawater is evaporated to make salt.
  \hfill \textit{The salt is mixed with ash from funeral pyres. It tastes of grief.}
\item[J] \textbf{The Sunken Gate} --- An underwater archway that leads to the drowned capital.
  \hfill \textit{It can only be reached by holding your breath for nine heartbeats.}
\item[Q] \textbf{Storm-Reader's Lookout} --- A wooden tower on the highest cliff.
  \hfill \textit{From here, the Storm-Reader can see three storms coming. He has never been wrong.}
\item[K] \textbf{The Ninth Wave's Rock} --- A black stone that the ninth wave strikes every full moon.
  \hfill \textit{The wave leaves behind a single white shell. The shell contains a memory.}
\item[A] \textbf{The Drowned Capital (Myth)} --- A sunken city said to rise once every century.
  \hfill \textit{Those who enter can speak to the dead. The price is a breath.}
\end{enumerate}

\paragraph*{(Wanderer/Maiden/Pirate)} Wanderer with a broken sword; Shrine Maiden with a copper bell; Pirate Captain with a brass hook.

\section*{Hearts — People & Factions}
\label{sec:nihori-people}
\begin{enumerate}
\setcounter{enumi}{1}
\item \textbf{Wanderer (Ronin)} --- A man with no master, a scarred blade, and a cup of sake.
  \hfill \textit{He offers a cup to every opponent. ``Drink,'' he says. ``The duel will be fair.''}
\item \textbf{Shrine Maiden} --- A woman with a copper bell on a rope. She rings it at dawn.
  \hfill \textit{The bell's tone tells the fishing fleet whether to sail. It has never been wrong.}
\item \textbf{Pirate Captain} --- A woman with a brass hook for a hand. She lost the original to a sea monster.
  \hfill \textit{The hook is also a key. It opens a chest of treasure maps.}
\item \textbf{Storm-Reader} --- Grey-haired, eyes the colour of storm clouds. He can feel the pressure change in his bones.
  \hfill \textit{``The storm will arrive in nine hours,'' he says. ``Bring your offerings.''}
\item \textbf{Swordsmith} --- A woman with calloused hands and a prayer strip on her forehead.
  \hfill \textit{She names each sword after a dead warrior. The sword remembers.}
\item \textbf{Duelist (Clan Champion)} --- A young man with a fresh scar. He has won nine duels.
  \hfill \textit{``The tenth will be my last,'' he says. He plans to retire after.}
\item \textbf{Tide-Hunter} --- A woman who hunts sea monsters. She has a harpoon made from a ship's keel.
  \hfill \textit{She has never killed a monster. She has ridden nine of them.}
\item \textbf{Volcano Priest} --- A man who tends the dormant crater. He speaks only in whispers.
  \hfill \textit{``The forbidden name is still here,'' he says. ``Do not ask what it is.''}
\item \textbf{Drowned Scholar's Ghost} --- A translucent figure who haunts the Sunken Gate.
  \hfill \textit{He asks: ``Has anyone retrieved my scroll?'' He will not say which scroll.}
\item[J] \textbf{Storm Emperor's Ghost} --- A column of steam over the water, shaped like a crowned man.
  \hfill \textit{He asks: ``Has anyone spoken the name?'' If you say yes, he weeps.}
\item[Q] \textbf{Ninth Wave's Herald} --- A child who appears before every storm. She points at the horizon.
  \hfill \textit{She never speaks. She only points. The direction tells where the storm will hit.}
\item[K] \textbf{The First Drowned (Legend)} --- The first Lethai-al who died in the sinking. She sits at the bottom of the sea.
  \hfill \textit{If you bring her a memory, she will teach you to breathe water.}
\item[A] \textbf{Kuva's Shadow (Hollow)} --- A vast darkness that covers the sea at midnight.
  \hfill \textit{It asks: ``Who called the ninth wave?'' The answer is a name no one knows.}
\end{enumerate}

\paragraph*{(Ashfall/Duel/Pirate)} Ashfall from the volcano — all actions Desperate; a duel challenge is issued — refusal loses face; a pirate flotilla claims the strait — pay toll or fight.

\section*{Clubs — Complications/Threats}
\label{sec:nihori-complications}
\begin{enumerate}
\setcounter{enumi}{1}
\item \textbf{Ashfall (Volcano)} --- The dormant crater emits ash. Visibility drops. All actions are Desperate.
  \hfill \textit{The ash is warm. It smells of old names. Speaking a name makes the ash thicken.}
\item \textbf{Duel Challenge (Issued)} --- A Lethai-al challenges a party member to a duel. Refusal loses face.
  \hfill \textit{The duel is to first blood. The loser owes a story. The story must be true.}
\item \textbf{Pirate Flotilla} --- Kalingan pirates block the strait. They demand a toll of one tide-charm.
  \hfill \textit{The charm must be from a drowned ancestor. Refuse, and they attack.}
\item \textbf{Storm-Spirit Offended} --- A traveller insulted the wind. The storm-spirit demands compensation.
  \hfill \textit{The compensation is a promise. The promise must be about the sea.}
\item \textbf{Tide-Charm Stolen} --- Someone steals a pouch of charms. The ancestors are angry.
  \hfill \textit{The thief is a foreign merchant. The charms are in a cargo hold.}
\item \textbf{Drowned Capital Rising} --- The sunken city emerges at low tide. Spirits walk the streets.
  \hfill \textit{They ask for news of the surface. Lies will anger them. Truths will sadden them.}
\item \textbf{Volcano Name Whispered} --- Someone speaks the forbidden name near the crater. The ground shakes.
  \hfill \textit{The speaker is a child. The child did not know. The volcano remembers.}
\item \textbf{Ninth Wave Warning} --- The full moon is tomorrow. The ninth wave will be larger than any in a century.
  \hfill \textit{The fishing fleet cannot be recalled. Someone must warn them.}
\item \textbf{Sake House Brawl} --- A dispute over a duel's fairness turns violent. The ninth cup is poured.
  \hfill \textit{The ninth cup is for the dead. Someone has just died. The brawl is now a blood feud.}
\item[J] \textbf{Sunken Gate Opens} --- The underwater archway becomes passable. Something swims out.
  \hfill \textit{It is a drowned spirit. It wants its name back. It will trade a secret for it.}
\item[Q] \textbf{Storm-Reader's Silence} --- The old man stops reading the storm. The fishermen do not know where to sail.
  \hfill \textit{He is grieving. His son died at sea. He will not speak until the body is found.}
\item[K] \textbf{Duel to the Death (Betrayal)} --- A duel that was meant to be to first blood becomes a killing.
  \hfill \textit{The killer is a spy from Ashaan. The duel was an assassination. The party is blamed.}
\item[A] \textbf{The Ninth Wave Takes} --- The massive wave crashes over the cliff shrine. The bells are silent.
  \hfill \textit{The shrine maiden is gone. The storm-spirit is free. The Hollow walks the shore.}
\end{enumerate}

\paragraph*{(Token/Bell/Charm)} Swordsmith's token (one reroll on a failed melee attack); shrine bell (reveals hidden spirits for one scene); wanderer's oath-bracelet (one use to compel a service).

\section*{Diamonds — Rewards/Leverage}
\label{sec:nihori-rewards}
\begin{enumerate}
\setcounter{enumi}{1}
\item \textbf{Swordsmith's Token (Iron)} --- A small iron disc. Reroll a failed melee attack once.
  \hfill \textit{The token is stamped with a prayer. The prayer is for a clean cut.}
\item \textbf{Shrine Bell (Copper)} --- A small copper bell. Ring it to reveal hidden spirits for one scene.
  \hfill \textit{The bell's tone is inaudible to humans. The spirits hear it clearly.}
\item \textbf{Wanderer's Oath-Bracelet (Leather)} --- A braided leather bracelet. Demand a service from any wanderer once.
  \hfill \textit{The bracelet is stained with sake. The wanderer cannot refuse.}
\item \textbf{Storm-Charm (Bone)} --- A carved bone. Ignore one weather penalty on a journey leg.
  \hfill \textit{The bone is from a sailor who died in a storm. His spirit protects you.}
\item \textbf{Tide-Charm (Shell)} --- A white shell. Whisper a drowned ancestor's name to calm a storm for one hour.
  \hfill \textit{The shell is warm. It glows when the ancestor is listening.}
\item \textbf{Duelist's Sake Cup} --- A ceramic cup used in a fair duel. Demand a duel to first blood.
  \hfill \textit{The cup is chipped. The chip is from a blade that drew first blood.}
\item \textbf{Pirate Captain's Hook (Brass)} --- A brass hook prosthetic. Use it as a key to open any chest in the cove.
  \hfill \textit{The hook is also a weapon. It can cut rope and bone.}
\item \textbf{Volcano Priest's Whisper} --- A single word: the forbidden name. Speak it to cause a tremor.
  \hfill \textit{The word burns the tongue. Use once. The volcano will remember.}
\item \textbf{Drowned Scholar's Scroll (Waterproof)} --- A scroll retrieved from the sunken capital. Contains a forgotten secret.
  \hfill \textit{The scroll is wet. It will never dry. The secret is about the Storm Emperor.}
\item[J] \textbf{Ninth Wave's Shell} --- A white shell from the Ninth Wave's Rock. Break it to experience a memory.
  \hfill \textit{The memory is always a death. The death is always a drowning.}
\item[Q] \textbf{Storm-Reader's Flute (Bone)} --- A bone flute. Play it to predict the weather for the next 24 hours.
  \hfill \textit{The tune is sad. It always predicts a storm. It is always right.}
\item[K] \textbf{Sunken Gate Key (Coral)} --- A piece of coral shaped like a key. Opens the Sunken Gate.
  \hfill \textit{The key is alive. It grows. It will only fit the gate once.}
\item[A] \textbf{First Drowned's Breath (Glass Vial)} --- Seawater from the bottom of the sunken capital. Drink it to breathe water for one hour.
  \hfill \textit{The water tastes of salt and sorrow. The lungs will ache.}
\end{enumerate}

\section*{Quick use notes}
\label{sec:nihori-quick-use}
\begin{itemize}
\item Draw until you have all four suits. Highest rank sets main Clock.
\item Diamonds are codified outcomes (tokens, bells, charms) that change position rather than call for a roll.
\item If any A appears, echo \textbf{sea-omens}—a wave that breaks twice, a bell that rings underwater, a name that echoes from the foam.
\end{itemize}

\subsection*{Additional Features}
\label{sec:nihori-features}

\begin{itemize}
\item \textbf{Duel Law:} Disputes in Nihori are settled by duel to first blood. The loser owes a story; the winner owes a cup of sake. A duel can be witnessed by anyone; the witness's word is law. Refusing a duel is an admission of guilt.
\item \textbf{Storm Season Clock [8]:} Tracks the approach of the monsoon. When it fills, a typhoon strikes — all travel becomes Desperate for 2 legs. The clock ticks each day during storm season.
\item \textbf{Tide-Charms:} Drowned ancestors become tide-charms — bones, shells, or shipwreck pieces that can calm storms. Each charm holds a memory. Using it may attract the Hollow.
\end{itemize}

\subsection*{Lore Echoes (Cross-Expansion Hooks)}
\begin{itemize}
  \item \textbf{The Drowned Capital and the Kalrantha Amulet (Dhahara):} The sunken city is said to contain a temple that holds a fragment of the demon Varash. The Veiled Sisters want it.
  \item \textbf{Tide-Charms and the Tulkani (Way of Silk):} Tulkani story-tellers sometimes carry Lethai-al tide-charms as proof of a witnessed drowning.
  \item \textbf{Storm-Spirits and the Veiled Sisters (Ashaan):} The Veiled Sisters have tried to bind storm-spirits for use in their fleet. The Lethai-al have foiled nine attempts.
\end{itemize}

\subsection*{Nihori Superstitions (burn for +1 die once)}
\begin{itemize}
  \item \textbf{Bread to the Foam:} Crumble a crust at the water's edge. Ignore the first \emph{Pirate Flotilla} penalty.
  \item \textbf{Knot the Rope:} Tie a single knot in a shrine climbing rope. Cancel \emph{Storm-Spirit Offended} or take –1 die on your next sailing roll.
  \item \textbf{Name to the Wave:} Whisper a drowned ancestor's name into the ninth wave. Gain +1 Effect when navigating a storm, but mark \emph{Obligation} to the Hollow.
\end{itemize}

\begin{tcolorbox}[colback=black!3,colframe=black!40!white,title={Patronage \& Power}]
In Nihori, power flows through duels, storms, and the memory of the drowned. A clan chief's authority rests on how many duels he has won; a shrine maiden's on how many storms she has predicted.

\textbf{For the GM:}  
Use the \textbf{Storm Season Clock} as a visible threat. Let \textbf{Duel Law} be a way to resolve conflicts without combat. Make \textbf{Tide-Charms} valuable but dangerous — each use may attract the Hollow.
\end{tcolorbox}

\newpage

% ------------------------------
% CHAPTER: AYOKHA
% ------------------------------
\chapter{Ayokha}
\emph{The Monsoon Throne.}

\noindent
{\Large
\textit{The monsoon does not come because we pray. It comes because the stars align. I read the stars. I align them. Without me, the rains would fail, the rice would burn, and the kingdom would fall. I am not a servant of the God-King. I am his monsoon.}

\vspace{0.5em}

I am a River Merchant. My boat has more secrets than cargo. The lower deck is for silk. The upper deck is for lies. The tiller is for bribes. I have carried nine spies this year. The ninth was the God-King's own daughter.

\vspace{0.5em}

Ayokha is the land of the Monsoon Throne, where the God-King rules by the rhythm of the rains. Rivers flood, rice terraces glow green, and the wharves are crowded with ships from Dhahara, Sihai, and distant isles. The Lethai-ar witness every contract; the spirit-mediums speak for the nats of the ford. Power here is measured in monsoon schedules, silk vouchers, and the favor of the God-King's astrologer. But here's the secret the astrologer doesn't want you to know: the monsoon does not come because he aligns the stars. The monsoon comes because the river-spirits allow it. And the river-spirits are angry.

\vspace{0.5em}

The nats---hungry spirits who live in eddies and under bridges---demand offerings of betel leaf, candle wax, and forgotten promises. A merchant who ignores a nat's demand will find his boat sunk and his name struck from the ledgers. A farmer who disrespects the monsoon will find his fields barren and his children dreaming of drowning. In Ayokha, the rain is not weather---it is a creditor. And the God-King is behind on his payments.

\vspace{0.5em}

Never refuse a gift of betel leaf. It is an oath of hospitality. Refusing it is an insult that can end a bloodline.

\vspace{1em}

\noindent\textit{--- The God-King's Astrologer, who has not seen the sun in thirty years}
}

\vspace{1em}

\begin{almanac}
\textbf{What do people eat for breakfast?} Rice congee with dried shrimp---the rice is from the terraces, the shrimp is from the delta nets. The delta is where the river meets the sea. The sea is the God-King's mother. The God-King's mother is a nat. The nat is hungry.

\textbf{What does a child learn before they can walk?} To read the river. Ayokha children are taught the eight currents: Flood, Ebb, Eddy, Whirl, Still, and the ninth---the current that flows upstream. The upstream current is the monsoon's breath. The monsoon's breath is the God-King's voice. The God-King is silent.

\textbf{What is the first thing you notice about a stranger?} Their offering. Every traveler in Ayokha carries a gift for the nats---betel leaf, candle, a coin. A stranger whose offering is generous has been blessed. A stranger whose offering is small is testing the spirits. A stranger with no offering is a fool. The nats take fools first.

\textbf{What is the most common crime?} Nat-bribery---but bribery is not a crime, it is a negotiation. The nats set the tolls. The merchants pay the tolls. The nats spend the tolls on offerings. The offerings are burned. The ash falls on the fields. The fields are fertile. The circle is the economy.

\textbf{What does no one talk about but everyone knows?} The God-King is a puppet. The astrologer is the real ruler. The astrologer reads the stars. The stars say what the astrologer tells them to say. The astrologer has not read the stars in twenty years. He reads the nats' moods instead. The nats are fickle.

\textbf{Where do people go when they die?} To the river. The dead are cremated on the ghats. Their ashes are scattered on the water. The water carries them to the delta. The delta is the land of the nats. The nats decide whether the dead become ancestors or hungry ghosts. The nats are not generous.

\textbf{How do you apologize for a serious wrong?} You offer a nat---a spirit of the river, summoned by a medium. The nat hears your apology. If the nat accepts, the water is calm. If the nat refuses, the water rises. If the nat is silent, the water is still. Still water is the God-King's judgment. The God-King does not forgive.

\textbf{What is the most important festival?} The Welcoming of the Monsoon (summer). The God-King stands in the rain. The rain is the nats' blessing. The blessing is measured in inches. The inches determine the harvest. The harvest determines the taxes. The taxes determine the God-King's wealth. The God-King is not wealthy.

\textbf{What is the secret that could destroy a powerful person?} The God-King's daughter is not his daughter. She is the child of the astrologer and a nat---born from a bargain made during a failed monsoon. She can control the rain. She does not know. The astrologer knows. The astrologer is waiting for the right moment to tell her. The right moment is the next drought.
\end{almanac}

\newpage
\section{Ayokha --- ``The Monsoon Throne''}
\label{chap:ayokha}

\begin{quote}
  \textbf{God-King's Astrologer (Elite):}  
  ``The monsoon does not come because we pray. It comes because the stars align. I read the stars. I align them. Without me, the rains would fail, the rice would burn, and the kingdom would fall. I am not a servant of the God-King. I am his monsoon.''
\end{quote}

\begin{quote}
  \textbf{River Merchant (Commoner):}  
  ``My boat has more secrets than cargo. The lower deck is for silk. The upper deck is for lies. The tiller is for bribes. I have carried nine spies this year. The ninth was the God-King's own daughter.''
\end{quote}

\begin{tcolorbox}[colback=black!3,colframe=blue!60!black,title={Theme \& Atmosphere}]
Ayokha is the land of the Monsoon Throne, where the God-King rules by the rhythm of the rains. Rivers flood, rice terraces glow green, and the wharves are crowded with ships from Dhahara, Sihai, and distant isles. The Lethai-ar witness every contract; the spirit-mediums speak for the nats of the ford. Power here is measured in monsoon schedules, silk vouchers, and the favour of the God-King's astrologer. The ninth taboo is strong: never pour the ninth cup, never count the ninth wave, never speak the ninth name. The Hollow watches from the flooded fields.

\vspace{2mm}
\textbf{What you notice first:} The smell of wet earth and fermenting rice. The way every bridge has a spirit shrine at its foot. The constant sound of water — dripping, flowing, raining.

\textbf{The one rule that kills newcomers:} Never refuse a gift of betel leaf. It is an oath of hospitality. Refusing it is an insult that can end a bloodline.
\end{tcolorbox}

\subsection*{Regional Mood: The Weight of the Monsoon}

Power here is measured in monsoon predictions, silk bales, and the number of spirits placated. The God-King rules, but the astrologer sets the calendar, the Lethai-ar witness the contracts, and the spirit-mediums negotiate with the nats. A merchant who ignores a nat's demand will find his boat sunk. A farmer who disrespects the monsoon will find his fields barren.

\textbf{Starting Location:} A river market on stilts, where boats load silk and spices. A Lethai-ar Vigil holds a green lamp: ``The God-King's astrologer has predicted a late monsoon. The merchants are panicking. The rice will fail unless someone sacrifices to the river spirit. You are someone. Go.''

\subsection*{Prominent NPCs of Ayokha}
\label{sec:ayokha-npcs}

\begin{longtable}{|p{0.2\textwidth}|p{0.25\textwidth}|p{0.3\textwidth}|p{0.15\textwidth}|}
\hline
\textbf{Name} & \textbf{Role/Title} & \textbf{Motivation} & \textbf{Rumoured Quirk} \\
\hline
God-King's Astrologer & Royal astronomer & Predict the monsoon before the people starve & He has not seen the sun in thirty years. He works by candlelight. \\
\hline
Vigil Lady Suyin & Lethai-ar witness-master & Keep the river's ledger free of forgeries & She has cut nine witness cords. The ninth was her own. \\
\hline
River Merchant & Silk trader & Smuggle contraband under the God-King's tariffs & She has a hidden deck that only appears at low tide. \\
\hline
Spirit-Medium & Nat negotiator & Speak for the spirit of the ford & She carries a bowl of water; the spirit speaks through ripples. \\
\hline
The Drowned Nat (Ghost) & A spirit of a broken oath & Demand payment before the monsoon fails & It appears as a wet figure in torn robes. It never blinks. \\
\hline
\end{longtable}

\subsection*{Names of Ayokha}
\label{sec:ayokha-names}

Inspired by the monsoon kingdoms: soft, flowing consonants, honorifics (Lady, Master). Surnames denote a river, a wharf, or a spirit.

\begin{longtable}{|p{0.25\textwidth}|p{0.25\textwidth}|p{0.2\textwidth}|p{0.2\textwidth}|}
\hline
\textbf{First Names (Male)} & \textbf{First Names (Female)} & \textbf{Clan/Epithet} & \textbf{Monsoon-Name} \\
\hline
Hiro & Suyin & River-watch & \textit{the Uncut} \\
\hline
Tama & Mira & Silk-hand & \textit{Ninth Knot} \\
\hline
Kaito & Yuki & Lamp-bearer & \textit{Green Flame} \\
\hline
Riku & Hana & Oath-strand & \textit{the Witness} \\
\hline
Sora & Nami & Mud-dyer & \textit{the Unforgotten} \\
\hline
\end{longtable}

\paragraph*{(Market/Temple/Bridge)} A river market on stilts; a step-pyramid temple with monsoon carvings; a bridge that appears only once a century.

\section*{Spades — Places}
\label{sec:ayokha-places}
\begin{enumerate}
\setcounter{enumi}{1}
\item \textbf{River Market (Stilts)} --- A floating market where goods are passed from boat to boat.
  \hfill \textit{The prices are written in water. They change with the tide.}
\item \textbf{Step-Pyramid Temple} --- A stone pyramid rising from the riverbank. Carvings of the first monsoon.
  \hfill \textit{The carvings weep water during the rainy season. The water is holy.}
\item \textbf{Bridge of Three Echoes} --- A wooden bridge that appears only once a century.
  \hfill \textit{Those who cross it see three versions of themselves: past, present, and future.}
\item \textbf{Silk Village (Mulberry Groves)} --- A village where silk is spun and cords are woven.
  \hfill \textit{The mulberry trees were planted by the first God-King. Their leaves are sacred.}
\item \textbf{Astrologer's Observatory} --- A roofless tower with a bronze armillary sphere.
  \hfill \textit{The sphere has nine rings. The ninth ring is broken. The astrologer refuses to fix it.}
\item \textbf{River Court (Floating)} --- A barge that serves as a mobile courthouse.
  \hfill \textit{The barge is anchored at a different spot each day. The river is the judge.}
\item \textbf{Spirit-Medium's Hut} --- A small shelter on the riverbank, surrounded by offering bowls.
  \hfill \textit{The bowls are filled with rice, flowers, and small coins. The spirits take the offerings at night.}
\item \textbf{Well of the Drowned Nat} --- A stone well where a spirit was trapped.
  \hfill \textit{The well is covered by a stone slab. The slab has nine holes. The spirit breathes through them.}
\item \textbf{God-King's Palace (Floating)} --- A wooden palace built on pontoons, anchored in the river.
  \hfill \textit{The palace rises and falls with the flood. The God-King never touches dry land.}
\item[J] \textbf{The Ninth Terrace} --- A rice terrace that is flooded only once every nine years.
  \hfill \textit{The rice grown there is for the God-King alone. Eating it is treason.}
\item[Q] \textbf{Monsoon Bell Tower} --- A stone tower with a single brass bell. The bell is rung at the first rain.
  \hfill \textit{The bell's tone determines the harvest. A high tone means abundance. A low tone means famine.}
\item[K] \textbf{The First Bridge (Myth)} --- Where the first God-King crossed the river and claimed the throne.
  \hfill \textit{The bridge is invisible. Only those who have witnessed nine true oaths can see it.}
\item[A] \textbf{The Monsoon's Eye} --- A still pool in the astrologer's observatory. It reflects the sky.
  \hfill \textit{The pool shows the future of the monsoon. The astrologer is the only one who can read it.}
\end{enumerate}

\paragraph*{(Merchant/Medium/Vigil)} River merchant with a boat of secrets; spirit-medium with a bowl of water; Lethai-ar Vigil with a green lamp.

\section*{Hearts — People & Factions}
\label{sec:ayokha-people}
\begin{enumerate}
\setcounter{enumi}{1}
\item \textbf{River Merchant} --- A woman with a brass earring and a hidden deck. She never sails without a bribe.
  \hfill \textit{``The tariff is negotiable,'' she says. ``My silence is not.''}
\item \textbf{Spirit-Medium} --- A woman with a bowl of water. She watches the ripples.
  \hfill \textit{``The nat says the river is angry,'' she says. ``Someone broke a promise.''}
\item \textbf{Vigil Lady Suyin} --- Grey-robed, a green lamp at her hip. She never blinks when witnessing an oath.
  \hfill \textit{``The cord will tighten,'' she says. ``I do not need to watch. The river watches for me.''}
\item \textbf{God-King's Astrologer} --- A pale man in silk robes, surrounded by candles.
  \hfill \textit{``The monsoon will be late,'' he says. ``The rice will fail unless...'' He trails off.}
\item \textbf{Rice Farmer} --- A woman with mud-stained feet and a straw hat. She leaves offerings at the terrace edge.
  \hfill \textit{``The spirits are hungry,'' she says. ``Feed them, or they feed on you.''}
\item \textbf{Silk Weaver} --- A man with thread between his fingers. He weaves witness cords for the Lethai-ar.
  \hfill \textit{``The cord must be perfect,'' he says. ``A single flaw can break an oath.''}
\item \textbf{Monsoon Bell Ringer} --- A young woman who rings the bell at the first rain. She has never been late.
  \hfill \textit{``The rain waits for me,'' she says. ``I do not wait for the rain.''}
\item \textbf{Drowned Nat's Ghost} --- A wet figure in torn robes, standing by the well.
  \hfill \textit{It asks: ``Who sealed me here?'' The answer is a name. The name is the God-King's.}
\item \textbf{God-King (Offstage)} --- Never seen, rarely heard. His decrees are delivered by astrologer.
  \hfill \textit{Rumour says he is a spirit. Rumour says he is a child. Rumour says he is the river itself.}
\item[J] \textbf{The First God-King's Ghost} --- A crowned skeleton in the palace's lower deck.
  \hfill \textit{He asks: ``Is the monsoon still on time?'' If you say no, he weeps.}
\item[Q] \textbf{Monsoon Spirit (Kami)} --- A whirlwind with eyes, seen over the river during storms.
  \hfill \textit{It asks: ``Who rang the bell?'' The answer must be the ringer's name.}
\item[K] \textbf{The Ninth Terrace's Keeper} --- An old woman who tends the forbidden rice.
  \hfill \textit{She has never eaten the rice. She has never been tempted.}
\item[A] \textbf{Kuva's Shadow (Hollow)} --- A vast darkness that covers the river at midnight.
  \hfill \textit{It asks: ``Who poured the ninth cup?'' The answer is a name. The name is the speaker's.}
\end{enumerate}

\paragraph*{(Flood/Rumor/Offense)} The river floods — a bridge is swept away; a rival merchant spreads rumors of cursed silk; the nat of the well is offended — no one may draw water.

\section*{Clubs — Complications/Threats}
\label{sec:ayokha-complications}
\begin{enumerate}
\setcounter{enumi}{1}
\item \textbf{River Floods} --- The monsoon rises. The bridge is gone. The wharf is isolated.
  \hfill \textit{The flood is not natural. A river-spirit is angry. An oath was broken upstream.}
\item \textbf{Rumour of Cursed Silk} --- A merchant claims that a bale of silk is cursed. Trade stops.
  \hfill \textit{The curse is a lie. The merchant wants to drive down prices.}
\item \textbf{Well-Spirit Offended} --- A traveller threw trash into the sacred well. The spirit refuses water.
  \hfill \textit{The spirit demands an apology. The apology must be an oath. The oath must be witnessed.}
\item \textbf{Monsoon Late} --- The astrologer's prediction fails. The rains do not come.
  \hfill \textit{The astrologer is a fraud. Or the spirits are angry. Either way, the rice is dying.}
\item \textbf{Witness Cord Cut (False)} --- Someone cuts a witness cord that was tied in good faith.
  \hfill \textit{The cutter is the Vigil's own apprentice. She thought the oath was false. It was true.}
\item \textbf{Silk Village Poisoned} --- Someone poisons the mulberry trees. The silkworms die.
  \hfill \textit{The poisoner is a rival merchant. He wants to monopolise the silk trade.}
\item \textbf{River-Spirit Demands a Story} --- The nat appears at the ford. It asks for a broken oath.
  \hfill \textit{It will not leave until someone confesses. The confession must be true.}
\item \textbf{God-King's False Oath} --- The king swears a vow he does not intend to keep. The Vigil knows.
  \hfill \textit{If she exposes him, the king will exile the Lethai-ar. If she stays silent, the cord tightens.}
\item \textbf{Astrologer's Betrayal} --- The astrologer sells the monsoon schedule to a foreign power.
  \hfill \textit{The foreign power is Ashaan. They want to time an invasion with the floods.}
\item[J] \textbf{The Ninth Cup Poured (By Mistake)} --- A stranger pours the ninth cup at a river market. The Hollow comes.
  \hfill \textit{The Hollow asks: ``Who poured this cup?'' The stranger points at the party.}
\item[Q] \textbf{First Bridge Appears} --- The mythic bridge manifests. Crossing it grants a new name.
  \hfill \textit{The new name is a debt. The debt is to the first God-King. He will collect.}
\item[K] \textbf{Monsoon Bell Cracked} --- The bell is struck too hard. It cracks. The tone is wrong.
  \hfill \textit{The bell-ringer is a spy. She cracked it on purpose. The monsoon will fail.}
\item[A] \textbf{The Monsoon Fails} --- The rains do not come for nine weeks. The river dries. The Hollow walks.
  \hfill \textit{The Hollow asks: ``Who failed to ring the bell?'' The answer is a name.}
\end{enumerate}

\paragraph*{(Schedule/Voucher/Cord)} Monsoon schedule (predicts safe sailing days); silk voucher (trade goods at 10\% profit); Lethai-ar witness cord (enforce any oath).

\section*{Diamonds — Rewards/Leverage}
\label{sec:ayokha-rewards}
\begin{enumerate}
\setcounter{enumi}{1}
\item \textbf{Monsoon Schedule (Astrologer's Copy)} --- A scroll of safe sailing days for the next season.
  \hfill \textit{The schedule is written on dried palm leaf. The leaf will crumble after use.}
\item \textbf{Silk Voucher (Merchant's Seal)} --- A stamped document. Trade goods at 10\% profit in any Ayokha port.
  \hfill \textit{The seal is a fish. The fish glows when the voucher is genuine.}
\item \textbf{Lethai-ar Witness Cord} --- A silk cord dyed with river mud. Enforce any oath.
  \hfill \textit{The cord tightens when the oath is broken. It will not stop until blood is drawn.}
\item \textbf{God-King's Blessing (Scroll)} --- A royal decree. One national favour in Ayokha.
  \hfill \textit{The scroll is sealed with river clay. The clay is from the palace's foundation.}
\item \textbf{River Merchant's Hidden Deck Key} --- A brass key. Access any smuggler's hold on the river.
  \hfill \textit{The key is warm. It glows when a customs boat is near.}
\item \textbf{Spirit-Medium's Bowl} --- A clay bowl. Fill it with water to speak with any river-spirit once.
  \hfill \textit{The spirit will answer one question. The answer will be a ripple.}
\item \textbf{Vigil's Lamp (Green Flame)} --- A small oil lamp. Gain +1 die to social rolls in any Lethai-ar court.
  \hfill \textit{The flame flickers when someone lies. It also flickers when you lie.}
\item \textbf{Silk Weaver's Needle} --- A needle carved from a mulberry branch. Tie a witness cord without cutting the silk.
  \hfill \textit{The needle can also cut a cord without making a sound. Use once.}
\item \textbf{Monsoon Bell Tone (Recording)} --- A clay cylinder that captured the bell's ring. Play it to summon a monsoon spirit.
  \hfill \textit{The spirit will come. It will demand an offering. The offering must be a promise.}
\item[J] \textbf{First Bridge Plank} --- A piece of the mythic bridge. Cross any river as if on solid ground once.
  \hfill \textit{The plank crumbles after use. The first God-King will remember.}
\item[Q] \textbf{Ninth Terrace Rice (Grain)} --- A single grain from the forbidden terrace. Eat it to gain a vision of the future.
  \hfill \textit{The vision is always a flood. The flood is always coming.}
\item[K] \textbf{Astrologer's Sphere Ring (Broken)} --- The ninth ring from the armillary sphere. Touch it to predict the monsoon's arrival.
  \hfill \textit{The ring is cold. It warms when the rains are nine days away.}
\item[A] \textbf{The Monsoon's Eye (Water)} --- A vial of water from the astrologer's pool. Drink it to see the future of the harvest.
  \hfill \textit{The water tastes of salt. The vision is always a famine. The famine is not inevitable.}
\end{enumerate}

\section*{Quick use notes}
\label{sec:ayokha-quick-use}
\begin{itemize}
\item Draw until you have all four suits. Highest rank sets main Clock.
\item Diamonds are codified outcomes (schedules, vouchers, cords) that change position rather than call for a roll.
\item If any A appears, echo \textbf{monsoon-omens}—water that flows upstream, a lamp that burns green without oil, a cord that ties itself.
\end{itemize}

\subsection*{Additional Features}
\label{sec:ayokha-features}

\begin{itemize}
\item \textbf{Monsoon Clock [6]:} Tracks the approach of the rainy season. When it fills, the monsoon arrives — all river travel becomes Desperate for 2 legs. The clock ticks each day during the dry season. The astrologer can predict the exact day.
\item \textbf{River Court Justice:} Disputes on the river are settled by a floating court. The judge is a Lethai-ar Vigil. The verdict is written on a clay tablet and thrown into the water. If the tablet floats, the verdict is true. If it sinks, the verdict is false — and the judge is drowned.
\item \textbf{Betel Leaf Hospitality:} Offering a betel leaf is an oath of hospitality. Refusing is a grave insult. Accepting binds both parties to non-violence for the duration of the visit.
\end{itemize}

\subsection*{Lore Echoes (Cross-Expansion Hooks)}
\begin{itemize}
  \item \textbf{Monsoon Trade and Dhahara:} Ayokha's monsoon winds carry ships to the Brass Coast. A delay in Ayokha means famine in Dhaharan ports.
  \item \textbf{Witness Cords and the Sidhi (Way of Silk):} Sidhi freedmen carry Lethai-ar witness cords as proof of emancipation. A cord from Vigil Suyin is worth more than a freedman's seal.
  \item \textbf{River Spirits and the Ukwe (Sekogo):} The Ayokhan river is said to flow from the Ukwe's roots. A spirit-medium who visits the Ukwe can read the river's memory without a context key.
\end{itemize}

\subsection*{Ayokha Superstitions (burn for +1 die once)}
\begin{itemize}
  \item \textbf{Bread to the River:} Crumble a crust into the current. Ignore the first \emph{River Floods} penalty.
  \item \textbf{Knot the Cord:} Tie a single knot in a witness cord before using it. Cancel \emph{Witness Cord Cut} or take –1 die on your next oath.
  \item \textbf{Name to the Well:} Whisper a dead Vigil's name into the well. Gain +1 Effect when witnessing an oath, but mark \emph{Obligation} to the Monsoon Spirit.
\end{itemize}

\begin{tcolorbox}[colback=black!3,colframe=black!40!white,title={Patronage \& Power}]
In Ayokha, power flows through the monsoon, the river, and the weight of a witnessed oath. The God-King rules, but the astrologer sets the calendar, the Lethai-ar hold the ledgers, and the spirit-mediums speak for the nats.

\textbf{For the GM:}  
Use the \textbf{Monsoon Clock} as a visible campaign timer. Let \textbf{River Court Justice} be a high-stakes social encounter. Enforce \textbf{Betel Leaf Hospitality} as a binding ritual — breaking it invites the Hollow.
\end{tcolorbox}

\newpage

\chapter{Aelaerem of the Eastern Realms}
\label{chap:aelaerem-east}

\emph{The Hollow's Neighbors.}

\noindent
{\Large
\textit{My grandmother taught me that the Hollow is not evil. It is a forgetfulness. It forgets where the boundary lies, forgets that we are not food, forgets that it was once something else. So we remind it. With a thread, a bell, a bowl of salt. We are its memory.}

\vspace{0.5em}

I am a Thread-Binder. I tied a red thread around my wrist this morning. The Hollow walked past my door last night. It saw the thread. It kept walking. That is the bargain. I give it a boundary. It respects the boundary. Mostly. But last week it forgot the thread's color. I had to tie two.

\vspace{0.5em}

The Aelaerem of the eastern realms have adapted their hedge-magic to thresholds of a different kind: bamboo groves that whisper in the dark, temple stairs that count your heartbeats, and the silence after a storm bell that has no clapper. Here, the Hollow wears the face of the uncanny --- not hunger, but absence. It is the space between two bamboo stalks, the step that is always missing, the name you cannot quite remember.

They survive through ritual precision: a red thread tied at a crossroads with exactly eight knots, a bowl of rice left at dusk and never checked until dawn, a single bell rung once and then muffled with a cloth that has been washed in salt water. They whisper that the Hollow is the shadow of Kuva, cast when the first sun turned away, and that it wanders the world looking for the name it forgot.

Unlike western Aelaerem, who live in close-knit villages, the Eastern Aelaerem are often solitary or live in small family groups---yamabushi of the boundary, shrine hermits of the between. They are the ones who know where the veil is thinnest, and they are the ones paid to keep it closed. Their magic is humble: a knot to bind a promise, a thread to guide a lost ghost, a cup of sake to buy a night's silence.

\vspace{0.5em}

Never count the steps between your door and the crossroads. The Hollow counts with you. If you reach nine, it will be waiting at the ninth step. It will ask, ``Where are you going?'' Do not answer.

\vspace{1em}

\noindent\textit{--- Hedge-Keeper Mila, who has never slept indoors}
}

\vspace{1em}

\begin{tcolorbox}[colback=black!3,colframe=blue!60!black,title={Theme \& Atmosphere},breakable]
\textbf{What you notice first:} Red threads tied to every gate, every bridge, every well. The smell of burnt rice and old salt. The way the Aelaerem never whistle after dusk---the Hollow thinks you are calling it. And the way they never, ever count to nine.

\textbf{The one rule that kills newcomers:} Never count the steps between your door and the crossroads. The Hollow counts with you. If you reach nine, it will be waiting at the ninth step. It will ask, ``Where are you going?'' Do not answer.
\end{tcolorbox}

\subsection*{Regional Mood: The Weight of the Unspoken}

Power here is measured in red threads tied, offerings left, and the gaps between words. The Aelaerem have no leaders, no courts, no written laws. Each hedge-keeper knows her own stretch of boundary. She knows where the Hollow walks, when it is restless, and what gesture will turn it aside. Her authority is not given by any king --- it is earned by surviving the uncanny. The Aelaerem are not feared or respected; they are tolerated, because without them the Hollow would seep into every threshold. They are the neighbors of the thing that lives between the notes of a song.

\subsection*{Starting Location}

A crossroads at dusk, where a red thread is tied around a wooden post. An old Aelaerem woman is pouring a cup of sake onto the ground. She looks up: ``The Hollow walked last night. I left salt. It took the salt. It left a single red thread tied around my wrist. It moved. Sleep with that knot, and it will know where you are. Do you want to know where it went? It went toward your camp.''

\subsection*{The Almanac of Ordinary Days (Eastern Aelaerem)}
\begin{almanac}
\textbf{What do people eat for breakfast?} Rice porridge with pickled vegetables---the rice is from the paddies, the vegetables are salted in clay pots buried at the boundary. The salt is blessed by the hedge-keeper. The vegetables are always a little bitter. The bitterness is a reminder.

\textbf{What does a child learn before they can walk?} To tie the eight boundary knots. The ninth knot is never tied---it is left loose for the Hollow to find. If a child ties the ninth knot, the Hollow will visit that night and ask for a name. The child's name.

\textbf{What is the first thing you notice about a stranger?} Their shadow. A stranger whose shadow is too short has been touched by the Hollow. A stranger whose shadow moves independently is already lost. A stranger with no shadow is the Hollow itself, wearing a borrowed face. Do not invite them in.

\textbf{What is the most common crime?} Thread-theft---stealing a red thread from a boundary marker. The thread is not valuable, but cutting it breaks a promise. The Hollow becomes restless. The thief is marked: their shadow will shrink one inch each night until they return the thread.

\textbf{What does no one talk about but everyone knows?} The Hollow is not a single thing. It is a place---the space between heartbeats, the pause before a sneeze, the silence after a bell. The Aelaerem do not guard against a monster. They guard against the gaps.

\textbf{Where do people go when they die?} To the salt. The dead are cremated, and their ashes are mixed with salt and scattered along the boundary. The salt draws the Hollow's attention away from the living. The dead become the new boundary. They do not rest. They watch.

\textbf{How do you apologize for a serious wrong?} You leave a cup of sake and a lock of hair at the place where the wrong was done. The Hollow will witness. If the sake is gone by dawn and the hair is untouched, you are forgiven. If the hair is gone, the Hollow has taken your apology as a promise. You are now bound. Do not break that promise.

\textbf{What is the most important festival?} The Retting of the Threads (spring). All red threads on the boundary are taken down, washed in salt water, and re-tied. The village is silent for one night---no bells, no songs, no whispers. The Hollow listens. The Hollow remembers. If a thread is tied wrong, the Hollow will correct it. The correction is always a death.

\textbf{What is the secret that could destroy a powerful person?} The first red thread was not tied by an Aelaerem. It was tied by the Hollow itself, before it forgot what it was. The thread is still there, at the First Boundary. If it is untied, the Hollow will remember its true name. Then it will stop being hungry. It will become something else. No one knows what.
\end{almanac}

\newpage
\section{Aelaerem --- ``The Hollow's Neighbors''}
\label{chap:aelaerem}

\begin{quote}
  \textbf{Hedge-Keeper (Elite):}  
  ``My grandmother taught me that the Hollow is not evil. It is hungry. There is a difference. A hungry thing can be fed. Pour a cup, tie a thread, leave a bell on the doorstep. Do not feed it too much, or it will learn your name.''
\end{quote}

\begin{quote}
  \textbf{Thread-Binder (Commoner):}  
  ``I tied a red thread around my wrist this morning. The Hollow walked past my door last night. It saw the thread. It kept walking. That is the bargain. I give it a boundary. It respects the boundary. Mostly.''
\end{quote}

\begin{tcolorbox}[colback=black!3,colframe=blue!60!black,title={Theme \& Atmosphere}]
The Aelaerem of Nihori and Ayokha have adapted their hedge-magic to the eastern thresholds. Here the Hollow wears a different face — it is the space between bamboo stalks, the silence after the storm bell, the missing step on the temple stair. They survive through ritual appeasement: a red thread tied at crossroads, a bowl of rice left at dusk, a single bell rung and then muffled. They whisper that the Hollow is the shadow of Kuva, cast when the sky first turned from the sun. Unlike western Aelaerem, who live in close-knit villages, the Eastern Aelaerem are often solitary or live in small family groups. They are the ones who know where the boundary between the living and the Hollow is thinnest, and they are the ones paid to keep it closed. Their magic is humble: a knot to bind a promise, a thread to guide a lost ghost, a cup of sake to buy a night's silence.

\vspace{2mm}
\textbf{What you notice first:} Red threads tied to every gate, every bridge, every well. The smell of burnt rice and old sake. The way the Aelaerem never whistle after dusk — the Hollow thinks you are calling it.

\textbf{The one rule that kills newcomers:} Never count the steps between your door and the crossroads. The Hollow counts with you. If you reach nine, it will be waiting at the ninth step.
\end{tcolorbox}

\subsection*{Regional Mood: The Weight of the Boundary}

Power here is measured in red threads tied, offerings left, and the distance from the Hollow's path. The Aelaerem have no leaders, no courts, no written laws. Each hedge-keeper knows her own stretch of boundary. She knows where the Hollow walks, when it is hungry, and what offering will turn it aside. Her authority is not given by any king — it is earned by survival. The Aelaerem are not feared or respected; they are tolerated, because without them the Hollow would walk through every village. They are the neighbours of the thing that watches from under the hill.

\textbf{Starting Location:} A crossroads at dusk, where a red thread is tied around a wooden post. An old Aelaerem woman is pouring a cup of sake onto the ground. She looks up: ``The Hollow walked last night. I left butter. It took the butter. It left a single red thread tied around my wrist. It moved. Sleep with that knot, and it will know where you are. Do you want to know where it went? It went toward your camp.''

\subsection*{Prominent NPCs of the Aelaerem}
\label{sec:aelaerem-npcs}

\begin{longtable}{|p{0.2\textwidth}|p{0.25\textwidth}|p{0.3\textwidth}|p{0.15\textwidth}|}
\hline
\textbf{Name} & \textbf{Role/Title} & \textbf{Motivation} & \textbf{Rumoured Quirk} \\
\hline
Hedge-Keeper Mila & Warden of the Moss Road & Keep the Hollow from crossing into the lowlands & She has never slept indoors. She sleeps at the boundary. \\
\hline
Thread-Binder Yuki & Ties red threads at crossroads & Bind the Hollow's path before it reaches the village & She carries a spindle of red thread that never runs out. \\
\hline
Offering-Keeper Taro & Leaves rice and sake at dusk & Feed the Hollow so it does not feed on the living & He has not eaten rice in nine years. He gives it all away. \\
\hline
Bell-Muffler Rin & Rings and silences copper bells & Signal to the Hollow that the boundary is attended & She wears earplugs made of wax. She does not need to hear the bell. \\
\hline
The Oath-Keeper (Hollow) & A spirit that collects unpaid obligations & Demand the ninth cup, the ninth step, the ninth name & It appears as a shadow with too many fingers. It never speaks first. \\
\hline
\end{longtable}

\subsection*{Names of the Aelaerem}
\label{sec:aelaerem-names}

Inspired by the rural traditions of the eastern islands and coast: short, earthy, often with repeated syllables. Surnames denote a boundary, a thread, or a hedge.

\begin{longtable}{|p{0.25\textwidth}|p{0.25\textwidth}|p{0.2\textwidth}|p{0.2\textwidth}|}
\hline
\textbf{First Names (Male)} & \textbf{First Names (Female)} & \textbf{Clan/Epithet} & \textbf{Thread-Name} \\
\hline
Taro & Mila & Hedge-kin & \textit{the Uncrossed} \\
\hline
Ren & Yuki & Thread-born & \textit{Ninth Knot} \\
\hline
Kai & Rin & Bell-muffler & \textit{the Silent} \\
\hline
Sho & Nami & Offering-hand & \textit{Boundary-Keeper} \\
\hline
Jiro & Hana & Root-walker & \textit{the Still} \\
\hline
\end{longtable}

\paragraph*{(Crossroads/Threshold/Well)} A crossroads with a red thread tied around a post; a threshold with a missing ninth step; a well with a bowl of rice at its base.

\section*{Spades --- Places}
\label{sec:aelaerem-places}
\begin{enumerate}
\setcounter{enumi}{1}
\item \textbf{Crossroads of the Red Thread} \hfill \textit{[BOUNDARY] [OMEN]} \\
  A wooden post at a meeting of three paths, wrapped in red cord. \\
  \textit{The thread is retied each full moon. The Hollow cannot cross a fresh knot.}
\item \textbf{Threshold of the Missing Step} \hfill \textit{[THRESHOLD] [TRAP]} \\
  A doorway with eight steps leading up. The ninth is absent. \\
  \textit{The gap is a handspan wide. Step over it. Do not step in it. The Hollow is below.}
\item \textbf{Offering Well} \hfill \textit{[RITUAL] [DEBT]} \\
  A stone well where rice and sake are left at dusk. \\
  \textit{The offerings are gone by dawn. The Hollow takes them. Sometimes it leaves a thread.}
\item \textbf{Boundary Stone} \hfill \textit{[BOUNDARY] [COLD]} \\
  A moss-covered rock that marks where the Hollow's territory begins. \\
  \textit{The stone is warm on one side and cold on the other. The cold side is the Hollow's.}
\item \textbf{Hedge-Keeper's Hut} \hfill \textit{[REFUGE] [LIMINAL]} \\
  A small shelter built exactly on the boundary line. \\
  \textit{The door faces both ways. Inside, the floor is split — one half is safe, one half is not.}
\item \textbf{Thread-Binder's Stall} \hfill \textit{[TRADE] [WARD]} \\
  A cart where red thread is sold by the knot. \\
  \textit{The thread is dyed with berry juice and iron filings. The Hollow is repelled by iron.}
\item \textbf{Bell-Muffler's Shrine} \hfill \textit{[RITUAL] [SILENCE]} \\
  A small stone shrine with a copper bell on a rope. \\
  \textit{The bell is rung once at dusk, then muffled with a leather strap. The Hollow hears the ring but not the silence.}
\item \textbf{Rice Paddy of the Ninth Harvest} \hfill \textit{[HARVEST] [OMEN]} \\
  A field where one stalk is never cut. \\
  \textit{The uncut stalk is an offering. The Hollow takes one grain each night. When the stalk is bare, the Hollow comes.}
\item \textbf{Oath-Keeper's Path} \hfill \textit{[DEATH] [CURSE]} \\
  A straight line of bare earth through the grass, seen only at dawn. \\
  \textit{The path was walked by a Oath-Keeper. Follow it, and you will find what it was collecting.}
\item[J] \textbf{The Hollow's Door (Seasonal)} \hfill \textit{[THRESHOLD] [DANGER]} \\
  A crack in a hillside that appears only on the new moon. \\
  \textit{The crack is warm. It smells of old promises. Do not enter.}
\item[Q] \textbf{The Ninth Cup's Altar} \hfill \textit{[RITUAL] [TABOO]} \\
  A stone table with eight cups. The ninth spot is empty. \\
  \textit{The ninth cup is for the Hollow. If you pour it, you must drink with it.}
\item[K] \textbf{Thread Tree} \hfill \textit{[OMEN] [MEMORY]} \\
  An old oak covered in red threads, each tied by a different hand. \\
  \textit{The threads are promises to the Hollow. Untie one, and the Hollow will remember the promise as broken.}
\item[A] \textbf{The First Boundary (Myth)} \hfill \textit{[MYTH] [PERMANENT]} \\
  Where the first Aelaerem tied the first red thread. \\
  \textit{The thread is still there. It has never been retied. It is older than memory.}
\end{enumerate}

\paragraph*{(Keeper/Binder/Muffler)} Hedge-Keeper Mila with a red thread around her wrist; Thread-Binder Yuki with a spindle of red thread; Bell-Muffler Rin with earplugs of wax.

\section*{Hearts --- People \& Factions}
\label{sec:aelaerem-people}
\begin{enumerate}
\setcounter{enumi}{1}
\item \textbf{Hedge-Keeper} \hfill \textit{[AUTHORITY] [KNOWLEDGE]} \\
  Old, grey-haired, a red thread tied around her left wrist. She never enters a house. \\
  \textit{``The Hollow watches the thresholds,'' she says. ``I sleep where it cannot see.''}
\item \textbf{Thread-Binder} \hfill \textit{[CRAFT] [RITUAL]} \\
  A young woman with a spindle of red thread. She ties knots at crossroads. \\
  \textit{``A fresh knot each full moon,'' she says. ``The Hollow respects the new thread. For a while.''}
\item \textbf{Offering-Keeper} \hfill \textit{[SACRIFICE] [HUNGER]} \\
  A man with a bowl of rice and a flask of sake. He leaves them at dusk. \\
  \textit{``The Hollow is hungry,'' he says. ``Feed it, and it will eat the offering instead of you.''}
\item \textbf{Bell-Muffler} \hfill \textit{[RITUAL] [SILENCE]} \\
  A woman with wax earplugs. She rings a copper bell at dusk, then muffles it. \\
  \textit{``The Hollow hears the ring,'' she says. ``It knows we are watching. The silence is the boundary.''}
\item \textbf{Oath-Keeper (Visible)} \hfill \textit{[SPIRIT] [DEBT]} \\
  A tall shadow with too many fingers. It walks the path at midnight. \\
  \textit{It does not speak. It points. The direction is where a debt is owed.}
\item \textbf{Thread-Cutter (Exile)} \hfill \textit{[EXILE] [CURSE]} \\
  A man who cut a red thread. He is marked by the Hollow. \\
  \textit{His shadow is missing. The Hollow took it. He follows his shadow's memory.}
\item \textbf{Rice Farmer (Boundary-Knower)} \hfill \textit{[HARVEST] [KNOWLEDGE]} \\
  A woman who leaves one stalk uncut each harvest. \\
  \textit{``The Hollow takes one grain a night,'' she says. ``When the stalk is bare, I move.''}
\item \textbf{Hollow-Watcher Child} \hfill \textit{[OMEN] [PERCEPTION]} \\
  A boy who can see the Hollow when it walks. He points. \\
  \textit{He never speaks. He only points. The adults throw salt where he points.}
\item \textbf{Boundary Breaker (Desperate)} \hfill \textit{[DESPERATION] [DEBT]} \\
  A traveller who crossed the boundary without an offering. \\
  \textit{The Hollow followed him. He asks the Aelaerem for help. The price is a story.}
\item[J] \textbf{The First Hedge-Keeper (Ghost)} \hfill \textit{[GHOST] [MYTH]} \\
  A translucent figure tying a thread at the mythic boundary. \\
  \textit{She asks: ``Is the thread still there?'' If you say yes, she smiles. If you say no, she weeps.}
\item[Q] \textbf{The Hollow's Voice} \hfill \textit{[SPIRIT] [OMEN]} \\
  A whisper from the crack in the hillside. It speaks in counting rhymes. \\
  \textit{``One, two, three, four, five, six, seven, eight,'' it says. Then it stops. It never says nine.}
\item[K] \textbf{The Thread Tree's Keeper} \hfill \textit{[MEMORY] [AUTHORITY]} \\
  An old woman who tends the oak. She knows every thread's story. \\
  \textit{``This thread was tied by a mother who lost her child,'' she says. ``The Hollow did not take the child. It took the thread.''}
\item[A] \textbf{Kuva's Shadow (The Hollow Itself)} \hfill \textit{[SPIRIT] [CURSE]} \\
  A vast darkness that covers the boundary at midnight. \\
  \textit{It asks: ``Who tied the first thread?'' The answer is a name no one remembers.}
\end{enumerate}

\paragraph*{(Ashes/Cut/Step)} Ash from a burnt offering — the Hollow has been fed; a cut red thread — the boundary is broken; a missing step — the Hollow has crossed.

\section*{Clubs --- Complications/Threats}
\label{sec:aelaerem-complications}
\begin{enumerate}
\setcounter{enumi}{1}
\item \textbf{Cut Red Thread} \hfill \textit{[BOUNDARY] [VIOLATION]} \\
  Someone cut a boundary thread. The Hollow crosses. The village is at risk. \\
  \textit{The cutter is a stranger who did not know. Ignorance is not a defence. The Hollow demands a new thread.}
\item \textbf{Hollow Walks (Debt Called)} \hfill \textit{[DEBT] [SPIRIT]} \\
  A Oath-Keeper appears at the crossroads. It points at a party member. \\
  \textit{The member owes a debt. The debt is forgotten. The Hollow remembers. Pay in story or memory.}
\item \textbf{Missing Offering} \hfill \textit{[RITUAL] [HUNGER]} \\
  The offering bowl is empty at dawn. The Hollow took it and left a mark. \\
  \textit{The mark is a single red thread tied around the bowl. The Hollow wants more.}
\item \textbf{Bell Not Muffled} \hfill \textit{[RITUAL] [ALERT]} \\
  Someone rings a bell at dusk and does not silence it. The Hollow hears. \\
  \textit{The Hollow follows the sound. It will arrive at midnight. Muffle the bell before then.}
\item \textbf{Ninth Step Taken} \hfill \textit{[TABOO] [TRAP]} \\
  A traveller counts the steps to a threshold and reaches nine. The Hollow appears. \\
  \textit{The Hollow asks: ``Who counted?'' The traveller points at the party. The party must answer.}
\item \textbf{Ninth Cup Poured} \hfill \textit{[TABOO] [DEBT]} \\
  A stranger pours the ninth cup at a hedge-keeper's table. The Hollow drinks. \\
  \textit{The Hollow is now a guest. It will not leave until it is offered a story. The story must be true.}
\item \textbf{Thread Tree Burned} \hfill \textit{[DESTRUCTION] [CURSE]} \\
  Someone sets fire to the oak. The threads burn. The promises are broken. \\
  \textit{The Hollow is furious. It will walk every boundary. The Aelaerem cannot stop it alone.}
\item \textbf{Boundary Stone Moved} \hfill \textit{[BOUNDARY] [SHIFT]} \\
  A traveller moves the moss-covered rock. The boundary shifts. \\
  \textit{The Hollow now owns a new stretch of land. The village must negotiate. The negotiator is the party.}
\item \textbf{Oath-Keeper's Path Crossed} \hfill \textit{[CURSE] [DEBT]} \\
  The party camps on the straight line of bare earth. The Hollow comes. \\
  \textit{It asks: ``Why are you on my path?'' The answer must be a promise. Break it, and you become a Oath-Keeper.}
\item[J] \textbf{Hollow's Door Opens} \hfill \textit{[THRESHOLD] [DANGER]} \\
  The crack in the hillside widens. Cold air pours out. \\
  \textit{The Hollow invites the party inside. The invitation is a test. Refuse, and the Hollow is offended. Accept, and you may not return.}
\item[Q] \textbf{Thread-Cutter's Shadow Returns} \hfill \textit{[SPIRIT] [DEBT]} \\
  The exile's shadow appears at the boundary. It is looking for him. \\
  \textit{The shadow asks: ``Where is my body?'' The body is the exile. He is hiding. Do not tell.}
\item[K] \textbf{First Thread Untied} \hfill \textit{[MYTH] [CURSE]} \\
  Someone unties the mythic thread at the First Boundary. The Hollow is free. \\
  \textit{The Hollow spreads across the land. The Aelaerem have one day to re-tie the thread. The party must help.}
\item[A] \textbf{Kuva's Shadow Falls} \hfill \textit{[CURSE] [DARKNESS]} \\
  The Hollow covers the sun. Day becomes night. The boundary is invisible. \\
  \textit{The Hollow asks every person: ``What is your name?'' Answer, and it knows you. Refuse, and it takes your shadow.}
\end{enumerate}

\paragraph*{(Thread/Cup/Stone)} Red thread bundle (one safe crossing); Hollow cup (proof of a debt paid in memory); offering plate (right to camp in a dangerous spot).

\section*{Diamonds --- Rewards/Leverage}
\label{sec:aelaerem-rewards}
\begin{enumerate}
\setcounter{enumi}{1}
\item \textbf{Red Thread Bundle} \hfill \textit{[WARD] [ONCE]} \\
  A bundle of red-dyed thread. Tie it at a crossroads to mark a safe boundary for one night. \\
  \textit{The thread glows faintly at dusk. The Hollow will not cross.}
\item \textbf{Hollow Cup (Clay)} \hfill \textit{[PROOF] [DEBT]} \\
  A clay cup that has been used for a Hollow offering. Present it to prove a debt was paid. \\
  \textit{The cup is warm. It smells of old sake. The Hollow will recognise it.}
\item \textbf{Offering Plate (Wooden)} \hfill \textit{[REFUGE] [ONCE]} \\
  A wooden plate used for rice offerings. Camp anywhere on the boundary for one night. \\
  \textit{The plate is carved with eight circles. The ninth is missing. Place it on the ground to claim safety.}
\item \textbf{Thread-Binder's Spindle (Iron)} \hfill \textit{[WARD] [ONGOING]} \\
  A spindle of red thread that never runs out. Tie an unbreakable boundary. \\
  \textit{The thread is cold. It will not burn. The Hollow cannot cross it. Use once per session.}
\item \textbf{Bell-Muffler's Strap (Leather)} \hfill \textit{[SILENCE] [ONCE]} \\
  A leather strap for silencing bells. Muffle any bell to hide its sound from the Hollow. \\
  \textit{The strap is stained with wax. The Hollow cannot hear a bell muffled by it.}
\item \textbf{Boundary Stone Fragment} \hfill \textit{[SENSE] [ONGOING]} \\
  A chip of the moss-covered rock. Touch it to sense the nearest boundary line. \\
  \textit{The fragment is cold on one side. The cold side points to the Hollow.}
\item \textbf{Oath-Keeper's Finger (Bone)} \hfill \textit{[SUMMON] [ONCE]} \\
  A dried finger bone from a Oath-Keeper. Use it to call a Oath-Keeper once. \\
  \textit{The Oath-Keeper will collect one debt for you. Then it will ask for yours.}
\item \textbf{Hedge-Keeper's Salt Pouch} \hfill \textit{[BANISH] [ONCE]} \\
  A leather pouch of salt blessed by a hedge-keeper. Throw it to banish a Oath-Keeper for one scene. \\
  \textit{The salt is bitter. The Oath-Keeper will taste it. It will remember your face.}
\item \textbf{Ninth Cup's Altar Shard} \hfill \textit{[NEGOTIATION] [ONCE]} \\
  A piece of the stone altar. Pour a cup of sake on it to summon the Hollow for negotiation. \\
  \textit{The Hollow will come. It will listen. It will not attack — unless you lie.}
\item[J] \textbf{Thread Tree Branch} \hfill \textit{[WARD] [PERMANENT]} \\
  A branch from the oak, still covered in red threads. Plant it to create a new boundary. \\
  \textit{The branch will grow. The threads will remain. The Hollow will respect the new tree.}
\item[Q] \textbf{First Hedge-Keeper's Thread (Fragment)} \hfill \textit{[PERMANENT] [ONCE]} \\
  A piece of the mythic first thread. Touch it to any boundary to make it permanent. \\
  \textit{The thread is old. It crumbles after use. The boundary will last a century.}
\item[K] \textbf{The Hollow's Name (Whisper)} \hfill \textit{[COMMAND] [ONCE]} \\
  A single word: the name the Hollow had before it became hungry. \\
  \textit{Speak it, and the Hollow will stop for one minute. Then it will be angry. Use once.}
\item[A] \textbf{Kuva's Shadow (Glass Shard)} \hfill \textit{[FORGET] [ONCE]} \\
  A piece of green glass that contains a fragment of the Hollow's original form. \\
  \textit{Once, touch it to make the Hollow forget a debt. The shard will shatter. The Hollow will not forget who broke it.}
\end{enumerate}

\section*{Quick use notes}
\label{sec:aelaerem-quick-use}
\begin{itemize}
\item Draw until you have all four suits: Spade = place, Heart = actor, Club = pressure, Diamond = leverage. Highest rank sets the main Clock (2--5 $\rightarrow$ 4, 6--10 $\rightarrow$ 6, J/Q/K $\rightarrow$ 8, A $\rightarrow$ 10).
\item Diamonds are codified outcomes (threads, cups, plates) that change position rather than call for a roll.
\item If any A appears, echo \textbf{boundary-omens}—a red thread that ties itself, a cup that pours without a hand, a step that counts itself.
\end{itemize}

\subsection*{Additional Features (Cultural Mechanics)}
\label{sec:aelaerem-mechanics}

\begin{itemize}
\item \textbf{The Empty Cup:} Once per journey, when you are lost or pursued, you may pour a cup of sake or water onto the ground and whisper a name that is not yours. The Hollow will lead you away from danger — but you will forget one minor memory (the GM chooses a trivial fact; if you have no minor memories, advance the Hollow Attention clock [4]). The forgotten memory can be recovered by returning to the same spot and speaking your own name.
\item \textbf{Neighbor's Courtesy:} When camping, if you leave a small offering (rice, salt, a button) at the edge of the light, the night's travel hazard roll is made with Advantage (roll twice, take the less severe). If you forget three nights in a row, the Hollow begins to watch. Start a Hollow Attention clock [4].
\item \textbf{Kuva's Eavesdrop:} Once per night, if you sleep under open sky, you may ask one yes/no question about the next day's weather; the GM answers truthfully, but the question must be whispered to the wind. On a 1, the wind answers but the Hollow also hears.
\item \textbf{Red Thread Binding:} You may tie a red thread around an object or a threshold to ward against minor spirits (Cap 1) for one scene. If the thread is cut, the spirits notice — start a Spirit Attention clock [4].
\end{itemize}

\subsection*{Lore Echoes (Cross-Expansion Hooks)}
\begin{itemize}
  \item \textbf{Red Threads and the Lethai-ar (Eastern Realms):} Lethai-ar witness cords are sometimes dyed with the same berry juice as Aelaerem boundary threads. A witness cord tied with Aelaerem red thread is considered unbreakable.
  \item \textbf{The Hollow's Name and the Veiled Sisters (Ashaan):} The name of the Hollow is said to be a forgotten aspect of Ikasha — the part of her that did not fall asleep. Breath-less agents sometimes whisper the name to distract Oath-Keepers.
  \item \textbf{Boundary Stones and the Mountain Satrapies (Dhahara):} Aeler engineers have tried to replicate Aelaerem boundary stones using breath-bonds. The attempt failed; the mountain does not recognise iron.
  \item \textbf{The Ninth Cup and the Tulkani (Way of Silk):} Tulkani caravans never pour a ninth cup even when not in Aelaerem territory. The custom spread from the hedge-keepers. A Tulkani who pours a ninth cup is exiled.
  \item \textbf{The Hollow's Hunger and the Savage Heart (The Wilds):} In the deep wilderness, the Hollow is sometimes conflated with the Red Maw — a bear-spirit of endless appetite. Aelaerem hedge-keepers leave offerings of raw meat at the boundary during winter, when the Hollow is said to wear fur.
\end{itemize}

\subsection*{Aelaerem Superstitions (burn for +1 die once)}
\begin{itemize}
  \item \textbf{Bread to the Boundary:} Crumble a crust at the base of a boundary stone. Ignore the first \emph{Cut Red Thread} penalty.
  \item \textbf{Knot the Thread:} Tie a single knot in a red thread at a crossroads. Cancel \emph{Ninth Step Taken} or take –1 die on your next stealth roll (your choice).
  \item \textbf{Name to the Hollow:} Whisper a dead hedge-keeper's name into the empty cup. Gain +1 Effect when negotiating with the Hollow this scene, but mark \emph{Obligation} to Kuva.
\end{itemize}

\begin{tcolorbox}[colback=black!3,colframe=black!40!white,title={Patronage \& Power}]
Among the Aelaerem, power flows through threads, offerings, and the memory of the boundary. A Hedge-Keeper's authority rests on how long she has kept the Hollow at bay. There are no kings, no councils, no written laws — only the boundary, and the knowledge of where it lies.

\textbf{For the GM:}  
Power in Aelaerem society is measured in threads tied, offerings given, and boundaries maintained. Patronage often takes the form of red thread bundles, hollow cups, or offering plates — each tied to a specific crossroads, threshold, or debt. To emphasize the cultural mechanics:
\begin{itemize}
\item Use \textbf{The Empty Cup} as a desperate escape tool, but with a real cost (forgotten memories, Hollow Attention).
\item Enforce \textbf{Neighbor's Courtesy} — offerings matter. A party that forgets to leave offerings will attract the Hollow.
\item Let \textbf{Kuva's Eavesdrop} be a useful divination, but with a risk of Hollow attention on a roll of 1.
\item Track \textbf{Red Thread Binding} as a temporary ward; cut threads have consequences.
\end{itemize}
In Aelaerem society, the boundary forgets nothing — and the Hollow remembers every thread.
\end{tcolorbox}

% Index entries
\index{Aelaerem|see{Hollow's Neighbors}}
\index{Theme generator}
\index{Hollow's Neighbors generator}
\index{Places!Aelaerem}
\index{People!Aelaerem}
\index{Complications!Aelaerem}
\index{Rewards!Aelaerem}
\index{Red thread binding}
\index{Empty cup ritual}
\index{Hollow attention}
\index{Boundary stones}
\index{Oath-Keepers}
\index{Ninth cup taboo}
\index{Thread tree}
\index{First boundary}

\section*{Thematic SB Spend Table}
\label{sec:aelaerem-sb}

\subsection*{Minor Complications (1 SB)}
\begin{itemize}
\item \textbf{Exposure:} Your actions draw unwanted attention from \textbf{hedge-keepers or Hollow-watchers}.
\item \textbf{Noise:} Sounds of your actions alert nearby \textbf{thread-binders or offering-keepers}.
\item \textbf{Trace:} Evidence of your passage marks your route for \textbf{Oath-Keepers or the Hollow}.
\item \textbf{Delay:} A brief but meaningful setback costs you \textbf{time or favourable boundary position}.
\item \textbf{Supply Strain:} Mark +1 segment on a relevant \textbf{resource clock}.
\end{itemize}

\subsection*{Moderate Setbacks (2 SB)}
\begin{itemize}
\item \textbf{Alarm Raised:} \textbf{Local hedge-keeper or thread-binder} becomes aware and begins responding.
\item \textbf{Position Lost:} You lose advantageous ground/cover/stealth due to \textbf{fog or boundary shift}.
\item \textbf{Foe Appears:} A \textbf{Oath-Keeper or Hollow manifestation} arrives on scene.
\item \textbf{Gear Trouble:} A piece of equipment becomes \textbf{Compromised/Neglected}.
\item \textbf{Lock/Barrier:} A simple obstacle now requires a test to overcome.
\end{itemize}

\subsection*{Serious Trouble (3 SB)}
\begin{itemize}
\item \textbf{Reinforcements:} Additional \textbf{spirits, hollow watchers, or thread-cutters} arrive.
\item \textbf{Key Gear Breaks:} A crucial tool/weapon becomes temporarily unusable.
\item \textbf{Major Twist:} The situation fundamentally changes—\textbf{thread cut/offering taken/ninth step counted}.
\item \textbf{Rail Tick:} Advance a relevant campaign/front clock by 1 segment.
\item \textbf{Condition Applied:} Mark \textbf{Fatigue 1/Harm 1/Condition} appropriate to fiction.
\end{itemize}

\subsection*{Major Turns (4+ SB)}
\begin{itemize}
\item \textbf{Trap Springs:} A prepared danger activates with full effect.
\item \textbf{Authority Arrival:} \textbf{Hedge-Keeper Mila, the Hollow's Voice, or Kuva's Shadow} intervenes.
\item \textbf{Scene Shift:} The environment changes dramatically—\textbf{boundary shifts/thread tree burns/hollow's door opens}.
\item \textbf{Patron Omen:} Divine/arcane forces take notice—\textbf{omen appears/blessing lost/curse manifests}.
\item \textbf{Narrative Pivot:} The story takes an unexpected turn that reframes objectives.
\end{itemize}

\subsection*{Region-Specific SB Options}
\begin{itemize}
\item \textbf{Aelaerem (Boundary Law):} A red thread tightens without being pulled, a cup of sake steams, a boundary stone moves an inch.
\item \textbf{Aelaerem (Hollow Law):} A shadow detaches from its owner, a whisper says ``nine'', a step appears where there was none.
\item \textbf{Aelaerem (Offering Law):} Rice steams without heat, salt piles form into a hand, a bell rings once then falls silent.
\end{itemize}

\subsection*{Pursuit Ladder (Near / Mid / Far)}
\begin{itemize}
\item \textbf{Advance/Withdraw (Wits + Survival):} On a hit, shift one band. On a strong hit, also impose \emph{Red Thread}.
\item \textbf{Cut the Debt (Presence + Sway):} On a hit, freeze range for one exchange; if a hollow cup is in your possession, +1 die.
\item \textbf{Weather the Hollow (Resolve + Spirit):} On a hit in \emph{Hollow Walks}, you pick who is pointed at; on a miss, the GM does.
\end{itemize}

\subsection*{Boss Seeds (Thread \& Boundary)}
\begin{itemize}
\item \textbf{The Thread-Cutter (Nine Severed Promises)} — \emph{Phases}: Boundaries Broken $\rightarrow$ Orchestrated Hollow $\rightarrow$ Boundary Coup. \emph{Cracks}: re-tie a cut thread, expose a false offering, survive a night on the Oath-Keeper's path.
\item \textbf{The Unfed Hollow (Hunger of the Ninth Cup)} — \emph{Phases}: Offerings Ignored $\rightarrow$ Procession of Empty Bowls $\rightarrow$ Thread-Baptism of a Hedge-Keeper. \emph{Cracks}: relic offering plate, survivor testimony, the first thread re-tied.
\end{itemize}
\newpage

\chapter{Aelinnel of the Inland Sea}
\label{chap:aelinnel}

\emph{The Bell-Tenders.}

\noindent
{\Large
\textit{The Hollow is not a place. It is a miscount. The ninth step that is never there, the ninth cup that is never poured, the ninth name that is never spoken. We do not fight the Hollow. We leave one empty chair, one unlit lamp, one uncounted stone. That is the price of sitting at the table.}

\vspace{0.5em}

I asked an Aelinnel bell-tender how she kept the tide-spirits from flooding her village. She tapped her temple. ``I count them,'' she said. ``One, two, three, four, five, six, seven, eight. Then I stop. They know that if I ever count nine, I am not counting them.''

\vspace{0.5em}

The Aelinnel of the East do not hide from the spirits --- they coexist. Small spirits inhabit every rock, tree, and wave. The Aelinnel have learned to negotiate with them using copper bells, precise counting, and occasional offerings of pickled plums. They are the engineers of shrines, the builders of floating bridges, and the keepers of the bell-lines that keep tide-spirits from straying. Some say Kuva, the sky, taught them the first counting song, and that the Hollow is simply the space where the count fails.

Unlike their western kin, who lean toward the Gift of the Mind, the Eastern Aelinnel practice a folk mathematics that blends logic and superstition. They believe every spirit has a number-name --- a frequency that can be rung out of a bell, a sequence of steps that can be walked, a knot that can be tied. To offend a spirit is to miscount.

\vspace{0.5em}

Never count to nine aloud near a shrine. The spirits count with you. When you reach nine, they will ask: ``What comes after nine?'' The answer is always the Hollow.

\vspace{1em}

\noindent\textit{--- Bell-Warden Tumble, who writes letters on copper wax}
}

\vspace{1em}

\begin{tcolorbox}[colback=black!3,colframe=blue!60!black,title={Theme \& Atmosphere},breakable]
\textbf{What you notice first:} The sound of copper bells --- not ringing in sequence, but chiming in overlapping waves. The way every bridge has a missing ninth plank. The smell of pickled plums and salt.

\textbf{The one rule that kills newcomers:} Never count to nine aloud near a shrine. The spirits count with you. When you reach nine, they will ask: ``What comes after nine?'' The answer is always the Hollow.
\end{tcolorbox}

\subsection*{Regional Mood: The Weight of the Count}

Power here is measured in bells rung, tides predicted, and the precision of a threshold. The Aelinnel have no kings or councils; their authority rests on the bell-wardens who maintain the shrines and the tide-counters who read the waves. A village that loses its bell-warden is a village that the tide-spirits will flood. A bridge without its ninth plank missing is a bridge that the Hollow will cross. The Aelinnel do not worship the spirits; they transact with them. An offering is not a prayer --- it is a payment. A bell is not a musical instrument --- it is a contract.

\subsection*{Starting Location}

A floating shrine on the Inland Sea, where copper bells hang from every corner. A Bell-Warden rings a bell eight times, then stops: ``The tide-spirit is restless. It wants a ninth bell. I will not give it. You --- count the waves. Tell me how many break before the sun sets. If you miscount, the spirit will ask you for a story. Do not miscount.''

\subsection*{The Almanac of Ordinary Days (Aelinnel)}
\begin{almanac}
\textbf{What do people eat for breakfast?} Rice porridge with pickled plums --- the plums are fermented in salt brine and stored in copper-lidded barrels. The brine is never poured out; it is passed down through generations. The oldest batch is said to be nine generations old. No one has tasted it.

\textbf{What does a child learn before they can walk?} To count the eight safe numbers. Children are taught to count pebbles, footsteps, and bell-rings --- but always to stop at eight. The ninth is never spoken. A child who says ``nine'' aloud is sent to sit by the missing plank of the bridge until they forget they ever said it.

\textbf{What is the first thing you notice about a stranger?} Their fingers. Aelinnel tide-counters often have only eight fingers --- the ninth was a miscount they paid for. A stranger with nine fingers has never made a mistake. A stranger with seven has paid a heavy debt. A stranger with no fingers is either a spirit or someone who counted to nine and lost everything.

\textbf{What is the most common crime?} Bell-theft --- but not for the copper. A stolen bell can be rung at the wrong time, creating a miscount. The thief is not punished by law; they are given to the tide-spirit whose count was broken. The spirit decides the price.

\textbf{What does no one talk about but everyone knows?} The First Bell was not cast by an Aelinnel. It was found --- already ringing, already counting. No one knows who made it. No one knows why it rang. The shards are scattered across the Inland Sea. When the ninth shard is found, the First Bell will ring again. The Hollow will hear it.

\textbf{Where do people go when they die?} To the bell-line. The dead are not buried. Their ashes are mixed with copper dust and cast into new bells. The bells are hung on the bell-lines. When the tide-spirits ring them, the dead speak. Mostly they say: ``Count carefully.''

\textbf{How do you apologize for a serious wrong?} You offer a count --- a public recitation of the wrong, spoken in eights. ``I took the offering. I broke the bell. I counted to nine.'' Then you wait. If a spirit rings a bell eight times, you are forgiven. If it rings nine times, you are cursed. If no bell rings, the Hollow will visit you at midnight.

\textbf{What is the most important festival?} The Retuning of the Bells (autumn). Every bell on the Inland Sea is struck once, in sequence, from the highest tide to the lowest. The sound travels for miles. The spirits listen. If a bell is out of tune, the spirits will correct it by morning. The correction is always a death --- but the bell will ring true for another year.

\textbf{What is the secret that could destroy a powerful person?} The ninth bell was never meant to be cast. It was a mistake --- a foundry worker poured one bell too many. They tried to hide it, but it rang once on its own. The tone was the Hollow's name. The bell still exists. It is buried under a shrine on a small island. No one knows which island. The map was miscounted.
\end{almanac}

\newpage
\section{Aelinnel Regional Generator}
łabel{chap:aelinnel}

\textbf{Key Demographic: } Aelinnel (Gnomes)
\textbf{Captial: } Aevrossa

\subsection*{Quick Reference}
\textbf{Tagline:} \textit{Count to eight. Stop. The ninth is for the Hollow.}

\begin{quote}
  \textbf{Bell−Warden of the Seventh Tide (Elite):}  
  ``The Hollow is not a place. It is a miscount. The ninth step that is never there, the ninth cup that is never poured, the ninth name that is never spoken. We do not fight the Hollow. We leave one empty chair, one unlit lamp, one uncounted stone. That is the price of sitting at the table.''
\end{quote}

\begin{quote}
  \textbf{Tide−Counter (Commoner):}  
  ``I asked an Aelinnel bell−tender how she kept the tide−spirits from flooding her village. She tapped her temple. \textit{I count them,} she said. \textit{One, two, three, four, five, six, seven, eight. Then I stop. They know that if I ever count nine, I am not counting them.''}
\end{quote}

\textbf{Genre / Mood:} Folk mathematics, coastal spirit−bargaining, horror of the miscount.

\textbf{Starting Location:} A floating shrine on the Inland Sea, where copper bells hang from every corner. A Bell−Warden rings a bell eight times, then stops: ``The tide−spirit is restless. It wants a ninth bell. I will not give it. You −−− count the waves. Tell me how many break before the sun sets. If you miscount, the spirit will ask you for a story. Do not miscount.''

\vspace{2mm}
\begin{tcolorbox}[colback=black!3,colframe=blue!60!black,title={Theme & Atmosphere}]
The Aelinnel of the East do not hide from the spirits −−− they coexist. Small spirits inhabit every rock, tree, and wave. The Aelinnel have learned to negotiate with them using copper bells, precise counting, and occasional offerings of pickled plums. They are the engineers of shrines, the builders of floating bridges, and the keepers of the bell−lines that keep tide−spirits from straying. Some say Kuva, the sky, taught them the first counting song, and that the Hollow is simply the space where the count fails. Unlike their western kin, who lean toward the Gift of the Mind, the Eastern Aelinnel practice a folk mathematics that blends logic and superstition. They believe every spirit has a number−name −−− a frequency that can be rung out of a bell, a sequence of steps that can be walked, a knot that can be tied. To offend a spirit is to miscount.

\vspace{2mm}
\textbf{What you notice first:} The sound of copper bells −−− not ringing in sequence, but chiming in overlapping waves. The way every bridge has a missing ninth plank. The smell of pickled plums and salt.

\textbf{The one rule that kills newcomers:} Never count to nine aloud near a shrine. The spirits count with you. When you reach nine, they will ask: ``What comes after nine?'' The answer is always the Hollow.
\end{tcolorbox}

\subsection*{Regional Mood: The Weight of the Count}

Power here is measured in bells rung, tides predicted, and the precision of a threshold. The Aelinnel have no kings or councils; their authority rests on the bell−wardens who maintain the shrines and the tide−counters who read the waves. A village that loses its bell−warden is a village that the tide−spirits will flood. A bridge without its ninth plank missing is a bridge that the Hollow will cross. The Aelinnel do not worship the spirits; they transact with them. An offering is not a prayer −−− it is a payment. A bell is not a musical instrument −−− it is a contract.

\subsection*{Prominent NPCs of the Aelinnel}
łabel{sec:aelinnel−npcs}

\begin{longtable}{|p{0.2\textwidth}|p{0.25\textwidth}|p{0.3\textwidth}|p{0.15\textwidth}|}
ℎline
\textbf{Name} & \textbf{Role/Title} & \textbf{Motivation} & \textbf{Rumored Quirk} \\
ℎline
Bell−Warden Tumble & Keeper of the Seventh Tide & Keep the bell−lines ringing and the Hollow at bay & He writes letters on copper wax; the wax is his only memory. \\
ℎline
Tide−Counter Mira & Reader of the ninth wave & Predict the flood before it comes & She counts on her fingers but has only eight. The ninth is a phantom. \\
ℎline
Shrine−Builder Orin & Engineer of floating shrines & Build a bridge that never needs a missing plank & He uses no nails; the wood is joined by counting songs. \\
ℎline
Spirit−Negotiator Luma & Bargainer with tide−spirits & Trade pickled plums for safe passage & She has never lost a negotiation. The spirits fear her. \\
ℎline
The Miscount (Ghost) & A bell−warden who counted to nine & Warn others of the Hollow's attention & It appears as a flickering shadow, holding one copper bell that never rings. \\
ℎline
\end{longtable}

\subsection*{Names of the Aelinnel}
łabel{sec:aelinnel−names}

Inspired by the inland sea cultures: short, melodic, often with doubled consonants. Surnames denote a tide, a shrine, or a counting song.

\begin{longtable}{|p{0.25\textwidth}|p{0.25\textwidth}|p{0.2\textwidth}|p{0.2\textwidth}|}
ℎline
\textbf{First Names (Male)} & \textbf{First Names (Female)} & \textbf{Clan/Epithet} & \textbf{Bell−Name} \\
ℎline
Tumble & Mira & Seventh Tide & \textit{the Nine−Counter} \\
ℎline
Orin & Luma & Copper−kin & \textit{Bell−Born} \\
ℎline
Kellen & Nori & Wave−reader & \textit{the Unrung} \\
ℎline
Pip & Suki & Shrine−hand & \textit{Ninth Plank} \\
ℎline
Rinn & Tali & Counting−song & \textit{the Still Bell} \\
ℎline
\end{longtable}

\paragraph*{(Shrine/Bridge/Cove)} A floating shrine with copper bells; a bridge with a missing ninth plank; a spirit cove where offerings are left.

\section*{Spades −−− Places}
łabel{sec:aelinnel−places}
\begin{enumerate}
\setcounter{enumi}{1}
ℹtem \textbf{Floating Shrine (Copper Bells)} −−− A wooden platform on pontoons, hung with bells.
  ℎfill \textit{The bells ring when a spirit passes. The tone tells if the spirit is friendly.} \texttt{[SPIRIT]}
ℹtem \textbf{Bridge of the Missing Plank} −−− A wooden bridge with eight planks. The ninth is absent.
  ℎfill \textit{The gap is exactly one foot wide. Step over it. Do not step in it.} \texttt{[THRESHOLD]}
ℹtem \textbf{Spirit Cove (Offering Rocks)} −−− A sheltered beach where offerings are left at low tide.
  ℎfill \textit{The offerings are rice, pickled plums, and small copper coins. The spirits take them at night.} \texttt{[SPIRIT]}
ℹtem \textbf{Bell−Line Tower} −−− A stone tower with a rope of bells stretching across the channel.
  ℎfill \textit{The rope is made of braided hair and seaweed. The bells are tuned to the tide's frequency.} \texttt{[BELL]}
ℹtem \textbf{Counting School} −−− A wooden hut where children learn the eight−number system.
  ℎfill \textit{The teacher has only eight fingers. She was born that way. The children count on theirs.} \texttt{[SOCIAL]}
ℹtem \textbf{Tide−Watcher's Post} −−− A high platform on stilts, overlooking the sea.
  ℎfill \textit{The watcher marks each wave with a pebble. Nine pebbles mean the tide is turning.} \texttt{[WEATHER]}
ℹtem \textbf{Shrine of the Uncounted} −−− A small shrine with no bells, no offerings, no threshold.
  ℎfill \textit{It is for the spirits that have no number. They are the most dangerous.} \texttt{[DEATH]}
ℹtem \textbf{Plum Orchard (Pickled)} −−− A grove of plum trees, the fruit fermented in salt brine.
  ℎfill \textit{The pickling barrels are sealed with copper lids. The lids have eight holes.} \texttt{[TRADE]}
ℹtem \textbf{Copper Foundry} −−− A forge where bells are cast and tuned.
  ℎfill \textit{The foundry is silent. The workers use hand signals. Sound would disturb the spirits.} \texttt{[CRAFT]}
ℹtem[J] \textbf{The Ninth Bell (Forgotten)} −−− A bell that was cast but never hung. It sits in a field, overgrown.
  ℎfill \textit{It rings once a year, on the night of the ninth tide. No one knows why.} \texttt{[MYSTERY]}
ℹtem[Q] \textbf{Hollow's Gate (Submerged)} −−− A stone arch underwater, visible only at the lowest tide.
  ℎfill \textit{The arch is covered in counting marks. Eight marks, then a space. The space is the gate.} \texttt{[DEATH]}
ℹtem[K] \textbf{The Counting Stone} −−− A flat rock carved with eight grooves. Water flows through them.
  ℎfill \textit{The ninth groove is missing. The water falls into a hole. The hole is bottomless.} \texttt{[MAGIC]}
ℹtem[A] \textbf{The First Bell (Myth)} −−− A bell that rang once and shattered. Its shards are scattered across the Inland Sea.
  ℎfill \textit{Find all nine shards, and you can ring the bell that ended the first war with the Hollow.} \texttt{[LEGEND]}
\end{enumerate}

\paragraph*{(Warden/Counter/Builder)} Bell−Warden Tumble with a copper bell; Tide−Counter Mira with eight fingers; Shrine−Builder Orin with a counting song.

\section*{Hearts −−− People & Factions}
łabel{sec:aelinnel−people}
\begin{enumerate}
\setcounter{enumi}{1}
ℹtem \textbf{Bell−Warden} −−− Old, gray−haired, a copper bell on a leather thong. He never speaks above a whisper.
  ℎfill \textit{``The spirits hear everything,'' he says. ``A whisper is a negotiation. A shout is a declaration of war.''} \texttt{[SOCIAL]}
ℹtem \textbf{Tide−Counter} −−− A young woman with eight fingers. She counts waves by tapping her knee.
  ℎfill \textit{``The ninth wave is the spirit's breath,'' she says. ``Do not count it. It counts you.''} \texttt{[WEATHER]}
ℹtem \textbf{Shrine−Builder} −−− A man with calloused hands and a wooden mallet. He never uses nails.
  ℎfill \textit{``Nails are for coffins,'' he says. ``Joinery is for shrines. The spirits respect wood grain.''} \texttt{[CRAFT]}
ℹtem \textbf{Spirit−Negotiator} −−− A woman with a basket of pickled plums. She speaks to the air.
  ℎfill \textit{``I offer eight plums,'' she says. ``The ninth is for later. Do you accept?'' The air ripples.} \texttt{[SOCIAL]} \texttt{[SPIRIT]}
ℹtem \textbf{Copper−Smith (Silent)} −−− A dwarf−like figure who works in complete silence. He uses hand signs.
  ℎfill \textit{His hands are scarred from molten copper. He has not spoken in forty years.} \texttt{[CRAFT]}
ℹtem \textbf{Tide−Spirit (Visible)} −−− A shimmer of light over the water, shaped like a woman.
  ℎfill \textit{It asks: ``How many waves have you counted today?'' Answer with eight, or it will flood.} \texttt{[SPIRIT]} \texttt{[FEAR]}
ℹtem \textbf{Counting Teacher} −−− An old woman with a slate and chalk. She teaches children to stop at eight.
  ℎfill \textit{``What comes after eight?'' she asks. The children chorus: ``Nothing. We stop.''} \texttt{[SOCIAL]}
ℹtem \textbf{Hollow−Watcher} −−− A boy who sits at the edge of the bridge at dusk, watching the missing plank.
  ℎfill \textit{He can see the Hollow when it crosses. He whistles to warn the village.} \texttt{[DEATH]}
ℹtem \textbf{Offering−Keeper} −−− A woman who tends the spirit cove. She sweeps the beach at dawn.
  ℎfill \textit{``The offerings are gone,'' she says. ``The spirits were satisfied. We have one more day.''} \texttt{[SPIRIT]}
ℹtem[J] \textbf{The Miscount's Ghost} −−− A flickering shadow with a copper bell that never rings.
  ℎfill \textit{It follows travelers who count to nine. It whispers: ``Nine. Nine. What comes after nine?''} \texttt{[DEATH]} \texttt{[FEAR]}
ℹtem[Q] \textbf{First Bell's Shard (Keeper)} −−− A woman who holds one of the nine shards. She guards it with her life.
  ℎfill \textit{The shard is warm. It hums at low tide. It wants to be reunited.} \texttt{[MAGIC]}
ℹtem[K] \textbf{The Still Bell (Legend)} −−− A bell that hangs in a submerged shrine. It rings only when the Hollow walks.
  ℎfill \textit{No one has heard it ring in a century. The Hollow has been quiet. Too quiet.} \texttt{[LEGEND]}
ℹtem[A] \textbf{Kuva's Echo (Sky Spirit)} −−− A voice from the bell−line tower, speaking in counting songs.
  ℎfill \textit{``You have counted eight,'' it says. ``That is enough. The ninth is mine.''} \texttt{[PATRON]}
\end{enumerate}

\paragraph*{(Squall/Pirate/Whirlpool)} A sudden squall — all navigation rolls at disadvantage; a pirate attack — pay cargo or fight; a spirit−whirlpool — lose 1 Supply to escape.

\section*{Clubs −−− Complications/Threats}
łabel{sec:aelinnel−complications}
\begin{enumerate}
\setcounter{enumi}{1}
ℹtem \textbf{Sudden Squall} −−− A storm rises from calm waters. Navigation DV +2. \texttt{[WEATHER]}
  ℎfill \textit{The squall was summoned by an offended tide−spirit. Someone miscounted.}
ℹtem \textbf{Pirate Attack (Kalingan)} −−− Fast dhows appear from behind a reef. They demand cargo. \texttt{[COMBAT]}
  ℎfill \textit{The pirates are not human. They are spirit−touched. They count in nines.}
ℹtem \textbf{Spirit−Whirlpool} −−− A whirlpool opens beneath the boat. Lose 1 Supply to escape. \texttt{[DEATH]} \texttt{[TRAVEL]}
  ℎfill \textit{The whirlpool is a spirit's mouth. It is hungry. Feed it a memory.}
ℹtem \textbf{Bell−Line Breaks} −−− A rope of bells snaps. The tide−spirits are no longer bound. \texttt{[SPIRIT]} \texttt{[TOLL]}
  ℎfill \textit{The tide rises. The village will flood in nine hours. Repair the line.}
ℹtem \textbf{Miscount (Offended Spirit)} −−− A traveler counted to nine near a shrine. The spirit demands compensation. \texttt{[SOCIAL]} \texttt{[TOLL]}
  ℎfill \textit{The compensation is a story that ends with the number eight. The spirit will check.}
ℹtem \textbf{Missing Offering} −−− The spirits took the offering, but they left a mark. The mark is a warning. \texttt{[OMEN]}
  ℎfill \textit{The warning says: ``The Hollow is coming. Ring the ninth bell.'' But the ninth bell does not exist.}
ℹtem \textbf{Copper Foundry Fire} −−− The silent forge catches fire. The bells crack. The tones are wrong. \texttt{[DISASTER]}
  ℎfill \textit{The fire was set by a rival bell−warden. He wants to replace the old bells with his own.}
ℹtem \textbf{Bridge Plank Replaced (By Mistake)} −−− Someone nails a ninth plank onto the bridge. The Hollow crosses. \texttt{[DEATH]}
  ℎfill \textit{The Hollow asks: ``Who built this bridge?'' The builder is the party's ally.}
ℹtem \textbf{Counting School Poisoned} −−− The teacher's chalk is laced with a memory−erasing dust. \texttt{[MAGIC]} \texttt{[TOLL]}
  ℎfill \textit{The children forget to stop at eight. They count to nine. The Hollow comes for them.}
ℹtem[J] \textbf{The Ninth Bell Rings (Once)} −−− The forgotten bell in the field rings. Everyone hears it. \texttt{[OMEN]} \texttt{[DEATH]}
  ℎfill \textit{The Hollow's gate opens. Spirits flood the Inland Sea. The Aelinnel are terrified.}
ℹtem[Q] \textbf{Tide−Spirit Demands a Name} −−− A spirit rises from the water. It asks: ``What is the ninth number?'' \texttt{[RIDDLE]} \texttt{[DEATH]}
  ℎfill \textit{There is no ninth number. The correct answer is silence. The wrong answer is death.}
ℹtem[K] \textbf{First Bell Shard Stolen} −−− Someone steals one of the nine shards. The Hollow stirs. \texttt{[THEFT]} \texttt{[TOLL]}
  ℎfill \textit{The thief is a foreign merchant. He wants to sell the shard to the Veiled Sisters.}
ℹtem[A] \textbf{Kuva's Silence} −−− The counting song stops. The sky is empty. No birds, no wind, no waves. \texttt{[APOCALYPSE]}
  ℎfill \textit{Kuva has turned away. The Hollow is free. The Aelinnel are counting their last breaths.}
\end{enumerate}

\paragraph*{(Charter/Bead/Stone)} Bell−charter (right to ring a named shrine); tide−bead (safe passage on the Inland Sea for one crossing); counting stone (one reroll on a navigation test in fog).

\section*{Diamonds −−− Treasures & Discoveries}
łabel{sec:aelinnel−rewards}
\begin{enumerate}
\setcounter{enumi}{1}
ℹtem \textbf{Bell−Charter (Copper Token)} −−− A small copper disc. Right to ring any named shrine bell once. \texttt{[SOCIAL]}
  ℎfill \textit{The token is stamped with eight dots. The ninth dot is missing.}
ℹtem \textbf{Tide−Bead (Glass)} −−− A blue glass bead. Safe passage on the Inland Sea for one crossing. \texttt{[TRAVEL]}
  ℎfill \textit{The bead was blessed by a tide−spirit. It glows when a storm is coming.}
ℹtem \textbf{Counting Stone (Smooth Pebble)} −−− A pebble with eight grooves. Reroll a failed navigation test in fog. \texttt{[TRAVEL]}
  ℎfill \textit{The stone was worn by waves for a century. It knows the tide's schedule.}
ℹtem \textbf{Offering Box Key (Copper)} −−− A small key. Opens any Aelinnel offering box. \texttt{[STEALTH]}
  ℎfill \textit{The key is warm. It glows when an offering has been accepted.}
ℹtem \textbf{Shrine−Builder's Mallet} −−− A wooden mallet with no nails. Build a temporary shrine anywhere. \texttt{[CRAFT]}
  ℎfill \textit{The shrine lasts one night. The spirits will respect it.}
ℹtem \textbf{Spirit−Negotiator's Plum (Pickled)} −−− A pickled plum in a brine jar. Offer it to any tide−spirit. \texttt{[SOCIAL]}
  ℎfill \textit{The spirit will grant one favor. The favor is never large. It is always fair.}
ℹtem \textbf{Bell−Line Rope (Hair and Seaweed)} −−− A length of braided rope. Use it to bind a tide−spirit for one hour. \texttt{[MAGIC]}
  ℎfill \textit{The rope smells of salt. The spirit will remember the binding.}
ℹtem \textbf{Tide−Counter's Pebble Pouch} −−− A leather pouch of marked pebbles. Predict the tide for the next day. \texttt{[WEATHER]}
  ℎfill \textit{Count the pebbles. Eight mean high tide. Seven mean low. Nine mean run.}
ℹtem \textbf{Copper Bell (Untuned)} −−− A small bell that has not been tuned. Ring it to confuse a spirit. \texttt{[COMBAT]}
  ℎfill \textit{The spirit will hesitate for one exchange. The bell will crack after use.}
ℹtem[J] \textbf{Forgotten Bell's Tone (Recording)} −−− A clay cylinder that captured the ninth bell's ring. \texttt{[MAGIC]}
  ℎfill \textit{Play it to open the Hollow's gate. Use once. The Hollow will be angry.}
ℹtem[Q] \textbf{First Bell Shard (Green Glass)} −−− A shard of the mythic first bell. Touch it to any spirit to command obedience for one minute. \texttt{[MAGIC]}
  ℎfill \textit{The shard is warm. It hums. The spirit will obey — but it will remember.}
ℹtem[K] \textbf{Kuva's Counting Song (Flute)} −−− A bone flute that plays the eight−note counting song. \texttt{[SOCIAL]}
  ℎfill \textit{Play it to calm any spirit. The spirit will listen. It will not attack.}
ℹtem[A] \textbf{The Hollow's Silence (Bell)} −−− A copper bell that rings only for the Hollow. Ring it to banish an Oath−Keeper. \texttt{[DEATH]}
  ℎfill \textit{The bell will ring once. Then it will shatter. The Hollow will not forget.}
\end{enumerate}

\section*{Quick use notes}
łabel{sec:aelinnel−quick−use}
\begin{itemize}
ℹtem Draw until you have all four suits: Spade = place, Heart = actor, Club = pressure, Diamond = treasure. Highest rank sets the main Timer (2−−5 $→$ 4, 6−−10 $→$ 6, J/Q/K $→$ 8, A $→$ 10).
ℹtem Treasures (Diamonds) are codified outcomes (charters, beads, stones) that change position rather than call for a roll.
ℹtem If any A appears, echo \textbf{bell−omens} — a bell that rings without being struck, a missing plank that appears, a count that reaches nine and stops.
\end{itemize}

\subsection*{Additional Features (Cultural Mechanics)}
łabel{sec:aelinnel−mechanics}

\begin{itemize}
ℹtem \textbf{Bell−Tending:} Once per scene, when near a shrine or threshold, you may ring a copper bell. The next environmental hazard (tide, fog, landslide) is delayed by 1 segment of the relevant timer. The bell's tone also pacifies minor spirits (Cap 1) for the scene.
ℹtem \textbf{Fae Courtesy:} An Aelinnel never says ``no'' to a spirit. They say ``that road is closed for repairs'' or ``the season is not right.'' This phrasing cancels the first Story Beat generated by a spirit encounter. If forced to refuse outright, they mark 1 Fatigue (the spirit's disappointment is heavy).
ℹtem \textbf{Kuva's Count:} Once per journey, a character may roll Wits + Lore (DV 3) to recall a forgotten sky−omen; on success, the next travel leg ignores the first Club card drawn. On a miss, the omen is true, but the GM draws an additional Club for that leg.
ℹtem \textbf{The Ninth Step:} When you cross a threshold that has a missing ninth step (a stair, a bridge, a door), you may count the absence. Gain +1 die to resist the Hollow's attention for the scene, but mark 1 Obligation (the Hollow now knows your name).
ℹtem \textbf{Local Patrons (Syncretic Names):}
  \begin{itemize}
    ℹtem \textbf{The Unrung} — The Witness. \textit{A bell that never rings, yet remembers every promise.}
    ℹtem \textbf{The Counter} — Mykkiel. \textit{The divine arithmetician; a miscount is a sin.}
    ℹtem \textbf{The Ninth Wave} — Khemesh. \textit{The pressure beneath the tide; the hunger that waits after eight.}
    ℹtem \textbf{Kuva's Echo} — The Traveler. \textit{The sky−road and the counting song that guides lost ships.}
  \end{itemize}
\end{itemize}

\subsection*{Lore Echoes (Cross−Expansion Hooks)}
\begin{itemize}
  ℹtem \textbf{The Bitter Root:} Daughters of the Cord sometimes serve as bell−wardens, using the Hollow's attention to collect on promises. A renegade Sister hides in a floating shrine, trying to untie her cord from a tide−spirit's bargain.
  ℹtem \textbf{Malachai, the Chained Angel:} A silver coin that never spends appears in offering boxes. Take it, and the spirits will accept any gift — but the coin will reappear in the mouth of a drowned sailor.
  ℹtem \textbf{The Grey Wanderer's Grimoire:} Aelinnel psions are called ``Echo−Readers.'' They hear the count of waves in their dreams and can predict a miscount three tides before it happens.
  ℹtem \textbf{The Threshold Compendium:} The Ninth Bell is a Class III Procession – the sound of the bell walks the Inland Sea once a year, and those who hear it must answer a riddle or join the march.
  ℹtem \textbf{The Ninth:} In any Aelinnel scene, the ninth bell rung in an hour opens a threshold. The wise stop at eight. Those who ring the ninth see a gray figure standing on the water — waiting.
\end{itemize}

\subsection*{Regional Superstitions (burn for +1 die once)}
\begin{itemize}
  ℹtem \textbf{Bread to the Bell:} Crumble a crust at the base of a shrine bell. Ignore the first \emph{Bell−Line Breaks} penalty. \texttt{[SUPERSTITION]}
  ℹtem \textbf{Knot the Rope:} Tie a single knot in a bell−line rope. Cancel \emph{Miscount} or take –1 die on your next spirit negotiation (your choice). \texttt{[SUPERSTITION]}
  ℹtem \textbf{Name to the Wave:} Whisper a dead tide−counter's name into the sea. Gain +1 Effect when counting waves this scene, but mark \emph{Obligation} to Kuva. \texttt{[SUPERSTITION]}
\end{itemize}

\subsection*{Thematic SB Spend Table}
łabel{sec:aelinnel−sb}

\paragraph{Minor Complications (1 SB)}
\begin{itemize}
ℹtem \textbf{Exposure:} Your actions draw unwanted attention from \textbf{bell−wardens or tide−spirits}.
ℹtem \textbf{Noise:} Sounds of your actions alert nearby \textbf{shrine builders or offering−keepers}.
ℹtem \textbf{Trace:} Evidence of your passage marks your route for \textbf{spirits or the Hollow}.
ℹtem \textbf{Delay:} A brief but meaningful setback costs you \textbf{time or favorable tide}.
ℹtem \textbf{Supply Strain:} Mark +1 segment on a relevant \textbf{resource timer}.
\end{itemize}

\paragraph{Moderate Setbacks (2 SB)}
\begin{itemize}
ℹtem \textbf{Alarm Raised:} \textbf{Local Bell−Warden or tide−counter} becomes aware and begins responding.
ℹtem \textbf{Position Lost:} You lose advantageous ground/cover/stealth due to \textbf{fog or rising water}.
ℹtem \textbf{Foe Appears:} An \textbf{offended tide−spirit or Hollow agent} arrives on scene.
ℹtem \textbf{Gear Trouble:} A piece of equipment becomes \textbf{Compromised/Neglected}.
ℹtem \textbf{Lock/Barrier:} A simple obstacle now requires a test to overcome.
\end{itemize}

\paragraph{Serious Trouble (3 SB)}
\begin{itemize}
ℹtem \textbf{Reinforcements:} Additional \textbf{spirits, bell−wardens, or Hollow manifestations} arrive.
ℹtem \textbf{Key Gear Breaks:} A crucial tool/weapon becomes temporarily unusable.
ℹtem \textbf{Major Twist:} The situation fundamentally changes — \textbf{bell−line breaks/miscount called/ninth bell rings}.
ℹtem \textbf{Rail Tick:} Advance a relevant campaign/front timer by 1 segment.
ℹtem \textbf{Condition Applied:} Mark \textbf{Fatigue 1/Harm 1/Condition} appropriate to fiction.
\end{itemize}

\paragraph{Major Turns (4+ SB)}
\begin{itemize}
ℹtem \textbf{Trap Springs:} A prepared danger activates with full effect.
ℹtem \textbf{Authority Arrival:} \textbf{Bell−Warden Tumble, Kuva's Echo, or the Miscount's Ghost} intervenes.
ℹtem \textbf{Scene Shift:} The environment changes dramatically — \textbf{shrine floods/bell−line snaps/ninth plank appears}.
ℹtem \textbf{Patron Omen:} Divine/arcane forces take notice — \textbf{omen appears/blessing lost/curse manifests}.
ℹtem \textbf{Narrative Pivot:} The story takes an unexpected turn that reframes objectives.
\end{itemize}

\paragraph{Region−Specific SB Options}
\begin{itemize}
ℹtem \textbf{Aelinnel (Bell Law):} Bells ring in harmony, then stop; a missing plank glows; the count of waves becomes audible.
ℹtem \textbf{Aelinnel (Spirit Law):} A tide−spirit laughs, an offering steams, a copper bell tarnishes instantly.
ℹtem \textbf{Aelinnel (Hollow Law):} The ninth step appears, then vanishes; a whisper says ``count again''; a shadow crosses the bridge.
\end{itemize}

\subsection*{Pursuit Ladder (Near / Mid / Far)}
\begin{itemize}
ℹtem \textbf{Advance/Withdraw (Wits + Navigation):} On a hit, shift one band. On a strong hit, also impose \emph{Fog Bank}.
ℹtem \textbf{Cut the Bell (Presence + Lore):} On a hit, freeze range for one exchange; if a bell−charter is in your possession, +1 die.
ℹtem \textbf{Weather the Squall (Resolve + Survival):} On a hit in \emph{Sudden Squall}, you pick who keeps their footing; on a miss, the GM does.
\end{itemize}

\subsection*{Boss Seeds (Bell & Count)}
\begin{itemize}
ℹtem \textbf{The Miscount−King (Nine Errors)} — \emph{Phases}: Bell−Lines Severed $→$ Orchestrated Miscount $→$ Shrine Coup. \emph{Cracks}: re−tie a bell−line, expose a false count, survive a night at the forgotten bell.
ℹtem \textbf{The Drowned Counter (Spirit of the Ninth Tide)} — \emph{Phases}: Waves Miscounted $→$ Procession of Empty Shrines $→$ Bell−Baptism of a Warden. \emph{Cracks}: relic tide−bead, survivor testimony, the ninth bell refusing to ring.
\end{itemize}

\begin{tcolorbox}[colback=black!3,colframe=black!40!white,title={Patronage & Power}]
Among the Aelinnel, power flows through bells, counts, and the precision of a threshold. A Bell−Warden's authority rests on the number of shrines he tends; a Tide−Counter's on the accuracy of her predictions. There is no government, only the shared understanding that a miscount kills everyone.

\textbf{For the GM:}  
Power in Aelinnel society is measured in bells rung, offerings given, and miscounts avoided. Patronage often takes the form of bell−charters, tide−beads, or counting stones — each tied to a specific shrine, tide, or spirit. To emphasize the cultural mechanics:
\begin{itemize}
ℹtem Use \textbf{Bell−Tending} as a hazard mitigation tool. Let players delay environmental timers by ringing bells.
ℹtem Enforce \textbf{Fae Courtesy} — never say ``no'' to a spirit. Creative phrasing is rewarded.
ℹtem Let \textbf{Kuva's Count} be a journey−saving ability, but with a risk of additional complications.
ℹtem Track \textbf{The Ninth Step} as a source of Hollow Obligation. Each crossing of a threshold is a negotiation.
\end{itemize}
In Aelinnel society, the count forgets nothing — and the ninth bell is never rung.
\end{tcolorbox}

% Index entries
∈dex{Aelinnel|see{Bell−Tenders}}
∈dex{Theme generator}
∈dex{Bell−Tenders generator}
∈dex{Places!Aelinnel}
∈dex{People!Aelinnel}
∈dex{Complications!Aelinnel}
∈dex{Rewards!Aelinnel}
∈dex{Bell−tending}
∈dex{Counting songs}
∈dex{Kuva's count}
∈dex{Ninth step}
∈dex{Fae courtesy}
∈dex{Tide−spirits}
∈dex{Bell−lines}
∈dex{Forgotten bell}
∈dex{Local Patrons!Aelinnel}

\newpage

\chapter{Lethai-al of Nihori}
\label{chap:lethai-al}

\emph{The Storm-Touched Wanderers.}

\noindent
{\Large
\textit{We do not own the land. We walk it. The road is our shrine. The storm is our law. The sea does not forgive, but it does forget. If you can wait long enough, the waves will erase your debt. But we do not wait. We walk. From shrine to shrine, from storm to storm, we carry the names of the drowned. The road is a prayer. Every step is a syllable.}

\vspace{0.5em}

My grandmother drowned in the ninth wave. Her name became a tide-charm. I keep it in a pouch of oiled leather. When the storm grows angry, I whisper her name. The sea remembers her. It calms for a moment --- long enough to breathe. We are wanderers. The road is our home. The dead travel with us.

\vspace{0.5em}

The Lethai-al of Nihori once guarded a drowned kingdom --- an archipelago that sank when a forbidden name was spoken. Their blood still carries the salt of that ancient catastrophe. Now they are wanderers in the oldest sense: masterless, landless, bound to no lord but fiercely bound to each other and to the storm spirits that swirl above the waves. They walk from shrine to shrine, from port to port, following the processional roads that circle the island.

They do not farm. They do not build. They carry everything they own on their backs or in wooden carts. Their duels are not to the death --- they are to the first draw of blood, after which the loser must join the procession and carry the winner's burden for a season. Their courts are crossroads. Their temples are the road itself. The Eastern Ninth lives here: the Hollow does not hide under hills --- it breathes in the salt foam. The ninth wave is never the largest; it is the one that remembers your name.

\vspace{0.5em}

Never walk under a torii gate without making an offering. The gate is a threshold. The road on the other side is not the same road you left.

\vspace{1em}

\noindent\textit{--- Procession Elder Kenji, who has walked the circuit 99 times}
}

\vspace{1em}

\begin{tcolorbox}[colback=black!3,colframe=blue!60!black,title={Theme \& Atmosphere},breakable]
\textbf{What you notice first:} The smell of salt and incense. The way the Lethai-al never sit with their backs to the sea. The constant, low murmur of waves being predicted --- not by instruments, but by the tilt of a head, the twitch of a finger. The worn wooden plaques hung from their belts, each carved with a dead ancestor's name.

\textbf{The one rule that kills newcomers:} Never walk under a torii gate without making an offering. The gate is a threshold. The road on the other side is not the same road you left.
\end{tcolorbox}

\subsection*{Regional Mood: The Processional Road}

Power here is measured in shrines visited, storms survived, and the number of dead ancestors carried as tide-charms. The Lethai-al have no king, no courts, no written law. Their authority is personal and processional: a Procession Elder's word is law because she has walked the circuit more times than anyone else. Their guilds are not guilds --- they are processional bands, each with its own rhythm, its own counting song, its own patron shrine. A band that walks too fast offends the spirits; one that walks too slow loses the tide. The Lethai-al do not believe in ownership of land; they believe in passage rights, storm oaths, and the memory of the drowned. Every crossroads is a court. Every storm is a messenger. Every sunrise is a departure.

\subsection*{Starting Location}

A procession trail winding along the sea cliffs, marked by weathered torii gates and small stone shrines. A group of Lethai-al walk in single file, chanting a counting song. One carries a palanquin with a covered object --- a tide-charm shrine. The leader stops: ``You stand on the processional road. The drowned are watching from the waves. Do you walk with us? If so, you must carry nothing. The road provides. The road takes away.''

\subsection*{The Almanac of Ordinary Days (Lethai-al)}
\begin{almanac}
\textbf{What do people eat for breakfast?} Dried fish and rice balls --- the fish is smoked over driftwood fires, the rice is pressed into shapes that fit in a pouch. The food is always cold. The Lethai-al do not stop to cook. The road does not wait.

\textbf{What does a child learn before they can walk?} To count the nine waves. Children are taught to watch the sea and count: one, two, three, four, five, six, seven, eight --- then stop. The ninth wave is the spirit's wave. It counts you. If you count it, you owe it.

\textbf{What is the first thing you notice about a stranger?} Their feet. A Lethai-al who has walked the circuit for many years has calloused soles cracked like dry earth. A stranger with soft feet has not walked far. A stranger with no feet is a ghost who lost the road.

\textbf{What is the most common crime?} Charm-theft --- stealing a tide-charm from a sleeping traveler. The thief gains the dead ancestor's favor for one night. The ancestor wakes. The ancestor is angry. The punishment is to carry the victim's burden for the rest of the circuit.

\textbf{What does no one talk about but everyone knows?} The drowned kingdom did not sink by accident. The forbidden name was spoken by a Lethai-al --- a woman who lost her child to the sea. She wanted the sea to remember. It remembered too well. Her name is not spoken. Her tide-charm hangs in the Sunken Gate, still wet.

\textbf{Where do people go when they die?} Into the tide-charms. The dead are not buried. Their names are carved onto shells and placed in oiled leather pouches. The living carry them. The dead walk with the living. When the pouch is lost, the dead become the Hollow.

\textbf{How do you apologize for a serious wrong?} You walk a crossroads circuit --- nine times around a crossroads shrine, counting your steps in eights. On the ninth circuit, you must speak the wrong aloud. If the wind carries your voice away, you are forgiven. If the wind throws it back, you are not. If the wind is silent, the Hollow is listening.

\textbf{What is the most important festival?} The Procession of the Ninth Wave (autumn). All Lethai-al bands walk to the shore at low tide. When the ninth wave crashes, they wade into the foam and whisper the names of those who drowned that year. The sea answers with a single wave. The wave is always the same height. The height is the number of names.

\textbf{What is the secret that could destroy a powerful person?} The Ninth Wave Herald is not a child. She is the first drowned --- the woman who spoke the forbidden name. She did not die. She became the ninth wave. She returns every year to point at the horizon. She is pointing at the day the sea will rise again. No one knows which day.
\end{almanac}

\newpage
\section{Lethai-al --- ``The Storm-Touched Wanderers of Nihori''}
\label{chap:lethai-al}

\begin{quote}
  \textbf{Storm-Reader (Elite):}  
  ``The sea does not forgive, but it does forget. If you can wait long enough, the waves will erase your debt. But I do not wait. I walk into the rain and let the storm decide. The wind has a voice. You just have to forget your name long enough to hear it.''
\end{quote}

\begin{quote}
  \textbf{Tide-Charm Keeper (Commoner):}  
  ``My grandmother drowned in the ninth wave. Her name became a tide-charm. I keep it in a pouch of oiled leather. When the storm grows angry, I whisper her name. The sea remembers her. It calms for a moment — long enough to breathe.''
\end{quote}

\begin{tcolorbox}[colback=black!3,colframe=blue!60!black,title={Theme \& Atmosphere}]
The Lethai-al of Nihori once guarded a drowned kingdom — an archipelago that sank when a volcanic name was spoken. Their blood still carries the salt of that ancient catastrophe. Now they are wanderers in the oldest sense: masterless, bound to no lord, but fiercely bound to each other and to the storm spirits that swirl above the waves. They dwell in hidden mountain groves, emerging to duel winds, pacify storm-spirits, and occasionally lift burdens from over-taxed merchants. Their strength is endurance rather than brute force. A Lethai-al can walk three days without sleep, hold breath beneath a tidal wave, and survive a fall that would shatter mortal bones. They say this is the blessing of their drowned ancestors, who learned to live where the sea did not want them. The Eastern Ninth lives here: the Hollow does not hide under hills — it breathes in the salt foam. The ninth wave is never the largest; it is the one that remembers.

\vspace{2mm}
\textbf{What you notice first:} The smell of salt and wet stone. The way the Lethai-al never sit with their backs to the sea. The constant, low murmur of waves being predicted — not by instruments, but by the tilt of a head, the twitch of a finger.

\textbf{The one rule that kills newcomers:} Never refuse a cup of sake before a duel. The first cup is courtesy, the second is negotiation, the third is fate. Refuse, and the duel becomes a death oath.
\end{tcolorbox}

\subsection*{Regional Mood: The Weight of the Ninth Wave}

Power here is measured in storms ridden, duels witnessed, and the number of drowned ancestors remembered. The Lethai-al have no king, no courts, no written law. Their authority is personal: a Storm-Reader's word is law because she has survived more typhoons than anyone else. Their duels are not to the death — they are to the first draw of blood, after which the loser must offer a gift. The gift is usually a story, a promise, or a tide-charm. The Lethai-al do not believe in ownership of land; they believe in passage rights, storm oaths, and the memory of the drowned.

\textbf{Starting Location:} A bamboo grove on a cliff overlooking the sea, where a duel is about to begin. Two Lethai-al face each other, cups of sake in hand. A Storm-Reader stands between them: ``The dispute is over a tide-charm. One says it was stolen. The other says it was inherited. The sea will not judge. You will. Watch the duel. The first blood tells the truth.''

\subsection*{Prominent NPCs of the Lethai-al}
\label{sec:lethai-al-npcs}

\begin{longtable}{|p{0.2\textwidth}|p{0.25\textwidth}|p{0.3\textwidth}|p{0.15\textwidth}|}
\hline
\textbf{Name} & \textbf{Role/Title} & \textbf{Motivation} & \textbf{Rumoured Quirk} \\
\hline
Storm-Reader Kenji & Elder of the cliff grove & Read the monsoon before it kills the fishing fleet & He has not slept in nine days; the storm keeps him awake. \\
\hline
Tide-Charm Keeper Yuki & Guardian of drowned names & Prevent the Hollow from taking her grandmother's charm & She carries an oiled pouch with nine charms. The ninth is empty. \\
\hline
Duelist Ronin (Unnamed) & Masterless warrior & Find a worthy opponent before the next full moon & He has never lost a duel. He has also never drawn first blood. \\
\hline
Shrine Maiden Hana & Storm-spirit negotiator & Pacify the spirit of the ninth wave & She rings a bell that can only be heard underwater. \\
\hline
The Drowned Captain (Ghost) & A Lethai-al who died before fulfilling an oath & Ask the living to complete his voyage & He appears on foggy nights, holding a skeletal oar. \\
\hline
\end{longtable}

\subsection*{Names of the Lethai-al}
\label{sec:lethai-al-names}

Inspired by the northern isles and coastal traditions: short, with vowel endings and occasional glottal stops. Surnames denote a tide, a wave, or a drowned ancestor.

\begin{longtable}{|p{0.25\textwidth}|p{0.25\textwidth}|p{0.2\textwidth}|p{0.2\textwidth}|}
\hline
\textbf{First Names (Male)} & \textbf{First Names (Female)} & \textbf{Clan/Epithet} & \textbf{Tide-Name} \\
\hline
Kenji & Yuki & Wave-walker & \textit{Ninth Breath} \\
\hline
Ronan & Hana & Storm-kin & \textit{Salt-Tongue} \\
\hline
Kael & Mira & Drowned-line & \textit{Unforgotten} \\
\hline
Takeshi & Sora & Foam-born & \textit{Tide-Charm} \\
\hline
Jiro & Nami & Ancestor-salt & \textit{the Unshadowed} \\
\hline
\end{longtable}

\paragraph*{(Grove/Cliff/Shrine)} A bamboo grove where duels are fought at dawn; a cliff shrine with ropes for climbing during high tide; a storm shrine on a peak, bells ringing in the wind.

\section*{Spades --- Places}
\label{sec:lethai-al-places}
\begin{enumerate}
\setcounter{enumi}{1}
\item \textbf{Bamboo Grove (Duel Ground)} --- A clearing where the stalks are scarred by blade cuts.
  \hfill \textit{The bamboo bleeds sap when a duel is dishonourable. The sap is red.}
\item \textbf{Cliff Shrine (Storm-Spirit)} --- A wooden platform on a sea cliff, with ropes for climbing.
  \hfill \textit{The ropes are braided with human hair. The hair is from sailors who survived shipwrecks.}
\item \textbf{Storm Shrine (Peak)} --- A stone circle with bells that ring in the wind.
  \hfill \textit{Each bell has a different tone. The tones predict the weather. The shrine maiden reads them.}
\item \textbf{Drowned Palace (Ruins)} --- Half-submerged columns visible at low tide.
  \hfill \textit{The columns are carved with names. The names are the drowned kingdom's royalty.}
\item \textbf{Tide-Charm Vault (Cave)} --- A sea cave where Lethai-al store their ancestors' charms.
  \hfill \textit{The charms are bones, shells, and pieces of shipwreck. Each hums with a different memory.}
\item \textbf{Foam-Ford (Crossing)} --- A shallow channel that can only be crossed at the exact moment of low tide.
  \hfill \textit{The crossing lasts nine breaths. Count carefully.}
\item \textbf{Volcano Shrine (Dormant)} --- A crater where the name that sank the kingdom was spoken.
  \hfill \textit{The crater is warm. Speaking any name aloud here causes the ground to tremble.}
\item \textbf{Sake House of the Ninth Cup} --- A tavern where duels are negotiated before they are fought.
  \hfill \textit{The ninth cup is never poured. It is for the dead.}
\item \textbf{Ancestor's Salt Pan} --- A flat where seawater is evaporated to make salt.
  \hfill \textit{The salt is mixed with ash from funeral pyres. It tastes of grief.}
\item[J] \textbf{The Sunken Gate} --- An underwater archway that leads to the drowned kingdom.
  \hfill \textit{It can only be reached by holding your breath for nine heartbeats. The gate is guarded.}
\item[Q] \textbf{Storm-Reader's Lookout} --- A wooden tower on the highest cliff.
  \hfill \textit{From here, the Storm-Reader can see three storms coming. He has never been wrong.}
\item[K] \textbf{The Ninth Wave's Rock} --- A black stone that the ninth wave strikes every full moon.
  \hfill \textit{The wave leaves behind a single white shell. The shell contains a memory.}
\item[A] \textbf{The Drowned Capital (Myth)} --- A sunken city said to rise once every century.
  \hfill \textit{Those who enter can speak to the dead. The price is a breath.}
\end{enumerate}

\paragraph*{(Reader/Keeper/Maiden)} Storm-Reader Kenji with a bone flute; Tide-Charm Keeper Yuki with an oiled pouch; Shrine Maiden Hana with a copper bell.

\section*{Hearts --- People \& Factions}
\label{sec:lethai-al-people}
\begin{enumerate}
\setcounter{enumi}{1}
\item \textbf{Storm-Reader} --- Grey-haired, eyes the colour of storm clouds. He can feel the pressure change in his bones.
  \hfill \textit{``The storm will arrive in nine hours,'' he says. ``Bring your offerings.''}
\item \textbf{Tide-Charm Keeper} --- A young woman with a leather pouch at her hip. She never swims.
  \hfill \textit{``My grandmother is in this pouch,'' she says. ``She cannot get wet.''}
\item \textbf{Duelist (Ronin)} --- A man with no master, a scarred blade, and a cup of sake.
  \hfill \textit{He offers a cup to every opponent. ``Drink,'' he says. ``The duel will be fair.''}
\item \textbf{Shrine Maiden} --- A woman with a copper bell on a rope. She rings it at dawn.
  \hfill \textit{The bell's tone tells the fishing fleet whether to sail. It has never been wrong.}
\item \textbf{Wave-Walker} --- A woman who can walk on water during a storm. She does not know how.
  \hfill \textit{``The storm carries me,'' she says. ``I just don't fight it.''}
\item \textbf{Sake Brewer} --- A fat man with a wooden ladle. He brews sake from rice and seawater.
  \hfill \textit{``The salt makes it strong,'' he says. ``The sea makes it honest.''}
\item \textbf{Drowned Scholar (Ghost)} --- A translucent figure who haunts the Sunken Gate.
  \hfill \textit{He asks: ``Has anyone spoken the name?'' He will not say which name.}
\item \textbf{Storm-Spirit (Kami)} --- A whirlwind with eyes. It speaks in the roar of the wind.
  \hfill \textit{It demands a gift before it will leave. The gift must be a promise.}
\item \textbf{Tide-Hunter} --- A woman who hunts sea monsters. She has a harpoon made from a ship's keel.
  \hfill \textit{She has never killed a monster. She has ridden nine of them.}
\item[J] \textbf{The Drowned Captain's Ghost} --- A skeleton in a captain's coat, carrying an oar.
  \hfill \textit{He asks: ``Will you complete my voyage?'' The voyage is to the sunken capital.}
\item[Q] \textbf{Volcano Priest} --- A man who tends the dormant crater. He speaks only in whispers.
  \hfill \textit{``The name that sank the kingdom is still here,'' he says. ``Do not ask what it is.''}
\item[K] \textbf{The Ninth Wave's Herald} --- A child who appears before every storm. She points at the horizon.
  \hfill \textit{She never speaks. She only points. The direction tells where the storm will hit.}
\item[A] \textbf{The First Drowned (Legend)} --- The first Lethai-al who died in the sinking. She sits at the bottom of the sea.
  \hfill \textit{If you bring her a memory, she will teach you to breathe water.}
\end{enumerate}

\paragraph*{(Duel/Pirate/Ashfall)} A duel challenge is issued — refusal loses face; a pirate flotilla claims the strait — pay toll or fight; ashfall from the volcano — all actions Desperate.

\section*{Clubs --- Complications/Threats}
\label{sec:lethai-al-complications}
\begin{enumerate}
\setcounter{enumi}{1}
\item \textbf{Duel Challenge (Issued)} --- A Lethai-al challenges a party member to a duel. Refusal loses face.
  \hfill \textit{The duel is to first blood. The loser owes a story. The story must be true.}
\item \textbf{Pirate Flotilla} --- Kalingan pirates block the strait. They demand a toll of one tide-charm.
  \hfill \textit{The charm must be from a drowned ancestor. Refuse, and they attack.}
\item \textbf{Ashfall (Volcano)} --- The dormant crater emits ash. Visibility drops. All actions are Desperate.
  \hfill \textit{The ash is warm. It smells of old names. Speaking a name makes the ash thicken.}
\item \textbf{Storm-Spirit Offended} --- A traveller insulted the wind. The storm-spirit demands compensation.
  \hfill \textit{The compensation is a promise. The promise must be about the sea. Break it, and the spirit drowns you.}
\item \textbf{Tide-Charm Stolen} --- Someone steals a pouch of charms. The ancestors are angry.
  \hfill \textit{The thief is a foreign merchant. The charms are in a cargo hold. The ship sails at dawn.}
\item \textbf{Drowned Capital Rising} --- The sunken city emerges at low tide. Spirits walk the streets.
  \hfill \textit{They ask for news of the surface. Lies will anger them. Truths will sadden them.}
\item \textbf{Volcano Name Whispered} --- Someone speaks the forbidden name near the crater. The ground shakes.
  \hfill \textit{The speaker is a child. The child did not know. The volcano remembers.}
\item \textbf{Ninth Wave Warning} --- The full moon is tomorrow. The ninth wave will be larger than any in a century.
  \hfill \textit{The fishing fleet cannot be recalled in time. Someone must warn them.}
\item \textbf{Sake House Brawl} --- A dispute over a duel's fairness turns violent. The ninth cup is poured.
  \hfill \textit{The ninth cup is for the dead. Someone has just died. The brawl is now a blood feud.}
\item[J] \textbf{Sunken Gate Opens} --- The underwater archway becomes passable. Something swims out.
  \hfill \textit{It is a drowned spirit. It wants its name back. It will trade a secret for it.}
\item[Q] \textbf{Storm-Reader's Silence} --- Kenji stops reading the storm. The fishermen do not know where to sail.
  \hfill \textit{He is grieving. His son died at sea. He will not speak until the body is found.}
\item[K] \textbf{Duel to the Death (Betrayal)} --- A duel that was meant to be to first blood becomes a killing.
  \hfill \textit{The killer is a spy from Ashaan. The duel was an assassination. The party is blamed.}
\item[A] \textbf{The Ninth Wave Takes} --- The massive wave crashes over the cliff shrine. The bells are silent.
  \hfill \textit{The shrine maiden is gone. The storm-spirit is free. The Hollow walks the shore.}
\end{enumerate}

\paragraph*{(Charm/Arrow/Bead)} Storm-charm (one weather bypass); broken arrow (proof of oath); ancestor shard (a piece of the drowned kingdom's crystal).

\section*{Diamonds --- Rewards/Leverage}
\label{sec:lethai-al-rewards}
\begin{enumerate}
\setcounter{enumi}{1}
\item \textbf{Storm-Charm} --- A bone carved with wave patterns. Use it to ignore one weather penalty on a journey leg.
  \hfill \textit{The bone is from a sailor who died in a storm. His spirit protects you.}
\item \textbf{Broken Arrow (Oath-Proof)} --- A snapped arrow shaft. Present it to prove an oath was fulfilled.
  \hfill \textit{The arrow was broken by a duelist who lost fairly. It carries his honour.}
\item \textbf{Ancestor Shard (Crystal)} --- A piece of crystal from the drowned kingdom. Once, reroll a failed swimming or endurance test.
  \hfill \textit{The shard is cold. It warms when an ancestor is near.}
\item \textbf{Tide-Charm (Leather Pouch)} --- A small pouch containing a drowned ancestor's name on a shell.
  \hfill \textit{Whisper the name to calm a storm for one hour. The ancestor will remember the favour.}
\item \textbf{Duelist's Sake Cup} --- A ceramic cup used in a fair duel. Present it to demand a duel to first blood.
  \hfill \textit{The cup is chipped. The chip is from a blade that drew first blood.}
\item \textbf{Shrine Maiden's Bell} --- A copper bell. Ring it to pacify a storm-spirit for one scene.
  \hfill \textit{The bell's tone is unique. The spirit will recognise it.}
\item \textbf{Volcano Priest's Whisper} --- A single word: the name that is forbidden. Speak it to cause a tremor.
  \hfill \textit{The word burns the tongue. Use it once. The volcano will remember.}
\item \textbf{Wave-Walker's Sandal} --- A single sandal woven from sea grass. Wear it to walk on calm water for one minute.
  \hfill \textit{The sandal will not get wet. Your foot will.}
\item \textbf{Drowned Captain's Oar} --- A skeletal oar that weighs nothing. Use it to row any boat against the current.
  \hfill \textit{The captain's ghost will row with you. He will ask for a story.}
\item[J] \textbf{Ninth Wave's Shell} --- A white shell found on the Ninth Wave's Rock. It contains a memory.
  \hfill \textit{Break it to experience the memory. The memory is always a death.}
\item[Q] \textbf{Storm-Reader's Flute} --- A bone flute. Play it to predict the weather for the next 24 hours.
  \hfill \textit{The tune is sad. It always predicts a storm. It is always right.}
\item[K] \textbf{Sunken Gate Key (Coral)} --- A piece of coral shaped like a key. Opens the Sunken Gate.
  \hfill \textit{The key is alive. It grows. It will only fit the gate once.}
\item[A] \textbf{First Drowned's Breath (Glass Vial)} --- A vial of seawater from the bottom of the sunken capital.
  \hfill \textit{Drink it to breathe water for one hour. The taste is salt and sorrow.}
\end{enumerate}

\section*{Quick use notes}
\label{sec:lethai-al-quick-use}
\begin{itemize}
\item Draw until you have all four suits: Spade = place, Heart = actor, Club = pressure, Diamond = leverage. Highest rank sets the main Clock (2–5 → 4, 6–10 → 6, J/Q/K → 8, A → 10).
\item Diamonds are codified outcomes (charms, arrows, shards) that change position rather than call for a roll.
\item If any A appears, echo \textbf{sea-omens}—a wave that breaks twice, a bell that rings underwater, a name that echoes from the foam.
\end{itemize}

\subsection*{Additional Features (Cultural Mechanics)}
\label{sec:lethai-al-mechanics}

\begin{itemize}
\item \textbf{Storm-Reading:} Once per journey, a character with Lethai-al training may predict the weather for the next 24 hours with perfect accuracy (Wits + Survival vs. DV 3). On success, choose one leg to ignore weather penalties. On a miss, the storm finds you anyway, and the GM gains 1 Story Beat to spend on a thunderclap or flash flood.
\item \textbf{Broken Oath Compensation:} When a Lethai-al breaks a promise (even a minor one), they must immediately offer compensation — a gift, a service, or a truth. If they refuse, they gain the condition \emph{Wave-Walker's Shame} (-1 die to all social rolls until they make amends). A Lethai-al who accepts a geas may gain +1 die to actions that serve it, but breaking it costs 2 Harm (Spirit).
\item \textbf{Tidal Resilience:} Once per scene, when a character would take Harm from a water-based or storm-based effect, they may reduce it by 1 level (to a minimum of 0). Mark 1 Fatigue.
\item \textbf{The Eastern Ninth (Hollow in Salt Foam):} When a Lethai-al drowns, their name becomes a tide-charm. Living kin may whisper it to calm a storm, but the Hollow will collect the debt at the next full moon. The Hollow's attention ticks whenever a tide-charm is used.
\end{itemize}

\subsection*{Lore Echoes (Cross-Expansion Hooks)}
\begin{itemize}
  \item \textbf{Tide-Charms and the Tulkani (Way of Silk):} Tulkani story-tellers sometimes carry Lethai-al tide-charms as proof of a witnessed drowning. A charm can be traded for a story that the Tulkani have not heard.
  \item \textbf{Storm-Spirits and the Veiled Sisters (Ashaan):} The Veiled Sisters have tried to bind storm-spirits for use in their fleet. The Lethai-al have foiled nine attempts. A Storm-Reader may ask the party to sabotage a Sister's ritual.
  \item \textbf{The Drowned Capital and the Kalrantha Amulet (Dhahara):} The sunken city is said to contain a temple that holds a fragment of the demon Varash. The amulet's creator visited the capital before it sank.
  \item \textbf{The Ninth Wave and the Hollow (Eastern Realms):} The Hollow's Debt-Walkers ride the ninth wave. Lethai-al who survive the ninth wave are said to owe the Hollow nothing. Aelaerem hedge-keepers leave offerings of salt at the high tide line to appease the wave.
\end{itemize}

\subsection*{Lethai-al Superstitions (burn for +1 die once)}
\begin{itemize}
  \item \textbf{Bread to the Foam:} Crumble a crust at the water's edge. Ignore the first \emph{Pirate Flotilla} penalty.
  \item \textbf{Knot the Rope:} Tie a single knot in a shrine climbing rope. Cancel \emph{Storm-Spirit Offended} or take –1 die on your next sailing roll (your choice).
  \item \textbf{Name to the Wave:} Whisper a drowned ancestor's name into the ninth wave. Gain +1 Effect when navigating a storm this scene, but mark \emph{Obligation} to the Hollow.
\end{itemize}

\begin{tcolorbox}[colback=black!3,colframe=black!40!white,title={Patronage \& Power}]
Among the Lethai-al, power flows through storms, duels, and the memory of the drowned. A Storm-Reader's word is law because she alone can read the wind. A Tide-Charm Keeper holds the names of the dead; she can calm a storm or call one. There are no kings, no courts, no written laws — only the sea, and the sea does not negotiate.

\textbf{For the GM:}  
Power in Lethai-al society is measured in storms survived, duels witnessed, and tide-charms kept. Patronage often takes the form of storm-charms, broken arrows, or ancestor shards — each tied to a specific ancestor, duel, or storm. To emphasize the cultural mechanics:
\begin{itemize}
\item Use \textbf{Storm-Reading} as a journey tool. Let players predict weather and avoid hazards, but make failure costly.
\item Enforce \textbf{Broken Oath Compensation} — a Lethai-al PC who breaks a promise must make amends immediately, or suffer social penalties.
\item Let \textbf{Tidal Resilience} be a lifesaver in shipwrecks, tidal waves, and storm scenes.
\item Track the \textbf{Hollow's Attention} when tide-charms are used. The Debt-Walkers may appear at the worst moment.
\end{itemize}
In Lethai-al society, the sea forgets nothing — and the ninth wave remembers every name.
\end{tcolorbox}

% Index entries
\index{Lethai-al|see{Storm-Touched Wanderers}}
\index{Theme generator}
\index{Storm-Touched Wanderers generator}
\index{Places!Lethai-al}
\index{People!Lethai-al}
\index{Complications!Lethai-al}
\index{Rewards!Lethai-al}
\index{Storm-reading}
\index{Tide-charms}
\index{Drowned kingdom}
\index{Ninth wave}
\index{Hollow in salt foam}
\index{Duel customs}
\index{Storm-spirits}
\index{Sunken capital}
\index{Saint Adar!Lethai-al}
\index{Saint Hervor!Lethai-al}
\index{Tide-Washed Covenant!Lethai-al}

\section*{Thematic SB Spend Table}
\label{sec:lethai-al-sb}

\subsection*{Minor Complications (1 SB)}
\begin{itemize}
\item \textbf{Exposure:} Your actions draw unwanted attention from \textbf{storm-spirits or pirate scouts}.
\item \textbf{Noise:} Sounds of your actions alert nearby \textbf{shrine maidens or duelists}.
\item \textbf{Trace:} Evidence of your passage marks your route for \textbf{tide-hunters or the Hollow}.
\item \textbf{Delay:} A brief but meaningful setback costs you \textbf{time or favourable tide}.
\item \textbf{Supply Strain:} Mark +1 segment on a relevant \textbf{resource clock}.
\end{itemize}

\subsection*{Moderate Setbacks (2 SB)}
\begin{itemize}
\item \textbf{Alarm Raised:} \textbf{Local Storm-Reader or shrine maiden} becomes aware and begins responding.
\item \textbf{Position Lost:} You lose advantageous ground/cover/stealth due to \textbf{ashfall or rising tide}.
\item \textbf{Foe Appears:} A \textbf{pirate crew or a storm-spirit} arrives on scene.
\item \textbf{Gear Trouble:} A piece of equipment becomes \textbf{Compromised/Neglected}.
\item \textbf{Lock/Barrier:} A simple obstacle now requires a test to overcome.
\end{itemize}

\subsection*{Serious Trouble (3 SB)}
\begin{itemize}
\item \textbf{Reinforcements:} Additional \textbf{pirates, storm-spirits, or volcanic ash clouds} arrive.
\item \textbf{Key Gear Breaks:} A crucial tool/weapon becomes temporarily unusable.
\item \textbf{Major Twist:} The situation fundamentally changes—\textbf{duel becomes death/storm shifts/sunken gate opens}.
\item \textbf{Rail Tick:} Advance a relevant campaign/front clock by 1 segment.
\item \textbf{Condition Applied:} Mark \textbf{Fatigue 1/Harm 1/Condition} appropriate to fiction.
\end{itemize}

\subsection*{Major Turns (4+ SB)}
\begin{itemize}
\item \textbf{Trap Springs:} A prepared danger activates with full effect.
\item \textbf{Authority Arrival:} \textbf{Storm-Reader Kenji, Drowned Captain, or Ninth Wave Herald} intervenes.
\item \textbf{Scene Shift:} The environment changes dramatically—\textbf{volcano erupts/storm breaks/sunken capital rises}.
\item \textbf{Patron Omen:} Divine/arcane forces take notice—\textbf{omen appears/blessing lost/curse manifests}.
\item \textbf{Narrative Pivot:} The story takes an unexpected turn that reframes objectives.
\end{itemize}

\subsection*{Region-Specific SB Options}
\begin{itemize}
\item \textbf{Lethai-al (Storm Law):} Lightning strikes the same place twice, rain falls upward, a bell rings without being struck.
\item \textbf{Lethai-al (Duel Law):} Blades glow faintly, spilled sake steams, the first drop of blood evaporates.
\item \textbf{Lethai-al (Drowned Law):} A name appears in wet sand, a wave carries a message, a drowned face looks up from still water.
\end{itemize}

\subsection*{Pursuit Ladder (Near / Mid / Far)}
\begin{itemize}
\item \textbf{Advance/Withdraw (Wits + Sail):} On a hit, shift one band. On a strong hit, also impose \emph{Storm Cover}.
\item \textbf{Cut the Duel (Presence + Sway):} On a hit, freeze range for one exchange; if a duelist's sake cup is in your possession, +1 die.
\item \textbf{Weather the Storm (Resolve + Survival):} On a hit in \emph{Storm-Spirit Offended}, you pick who is struck by lightning; on a miss, the GM does.
\end{itemize}

\subsection*{Boss Seeds (Wave \& Name)}
\begin{itemize}
\item \textbf{The Name-Eater (Nine Drownings)} — \emph{Phases}: Tide-Charms Stolen $\rightarrow$ Orchestrated Storm $\rightarrow$ Drowned Coup. \emph{Cracks}: recover a stolen charm, expose a false ancestor, survive the ninth wave.
\item \textbf{The Unforgotten Storm (Spirit of the Ninth Wave)} — \emph{Phases}: Waves Rising $\rightarrow$ Procession of Empty Shrines $\rightarrow$ Foam-Baptism of a Storm-Reader. \emph{Cracks}: relic storm-charm, survivor testimony, the ninth wave breaking on a calm sea.
\end{itemize}

\newpage

\chapter{Lethai-ar of Ayokha}
\label{chap:lethai-ar-ayokha}

\emph{The River-Bound Oath-Keepers.}

\noindent
{\Large
\textit{A witness cord is not a leash. It is a mirror. If you speak truly, you see yourself. If you lie, you see what you could have been. I have cut two cords in my life. I still dream of the faces in the silk. The river remembers every word spoken on its banks. We are merely the scribes.}

\vspace{0.5em}

My lamp has burned for nine years without being refilled. The flame is green. It flickers when a merchant lies about the weight of his silk. The river does not need oil. It needs witnesses. I sit on this wharf and watch. That is my offering.

\vspace{0.5em}

The Lethai-ar of Ayokha are the living bridges between mortals and the river spirits that launder the world's promises. Their ancestors were river-folk that ascended to the realm of the fair folk, and they now walk the middle path, a step between bodies of water and the world of men. They keep lanterns that never go out, record every contract spoken on the monsoon wharves, and bind their vows with silken cords dyed in river mud. Unlike their northern kin who swore to distant deities, the Lethai-ar serve \textbf{Benzaiten} (the river of words) or \textbf{Inari} (the fox's ledger). Their vows are physical cords, braided from their own hair and dyed with river mud. When a vow is fulfilled, the cord is burned. When it is broken, the cord tightens around the oath-breaker's wrist until blood is drawn. The Ayokhan taboo against the ninth cup is enforced by the Lethai-ar: the ninth cup is for the river spirits, and pouring it is an invitation to the Hollow.

\vspace{0.5em}

Never pour the ninth cup of tea or sake in a Lethai-ar's presence. The ninth cup is for the witness. Pour it, and you have invited the river to judge you. The river is not merciful.

\vspace{1em}

\noindent\textit{--- Vigil Lady Suyin, who has cut nine witness cords --- the ninth was her own}
}

\vspace{1em}

\begin{tcolorbox}[colback=black!3,colframe=blue!60!black,title={Theme \& Atmosphere},breakable]
\textbf{What you notice first:} The smell of river water and oiled silk. The way every Lethai-ar carries a small oil lamp, even at noon. The constant murmur of contracts being witnessed --- a merchant's promise, a captain's oath, a lover's vow.

\textbf{The one rule that kills newcomers:} Never pour the ninth cup of tea or sake in a Lethai-ar's presence. The ninth cup is for the witness. Pour it, and you have invited the river to judge you. The river is not merciful.
\end{tcolorbox}

\subsection*{Regional Mood: The Weight of the Witness Cord}

Power here is measured in lanterns kept lit, cords tied, and the number of oaths witnessed. The Lethai-ar have no army, no treasury, no territory. They have the river, and the river remembers everything. Every contract in Ayokha requires a Lethai-ar witness to be legally binding. The God-King's astrologers may set the calendar, but the Lethai-ar hold the ledger. Their neutrality is absolute: they will witness a vow between enemies as readily as between allies. The price of a witness cord is not silver --- it is a promise to keep the river clean, to light a lamp at dusk, to speak no word that cannot be unsaid.

\subsection*{Starting Location}

A monsoon wharf on the great river, where silk barges are being loaded. A Vigil in gray robes holds a lantern of green flame: ``The merchant and the captain have made a vow. The captain says the silk is eight bales. The merchant says it is nine. The witness cord is tied. It will tighten around the liar's wrist. I do not need to know who is lying. The cord will tell me. But you --- you are not bound. Watch. Learn. And do not pour the ninth cup.''

\subsection*{The Almanac of Ordinary Days (Lethai-ar of Ayokha)}
\begin{almanac}
\textbf{What do people eat for breakfast?} Rice porridge with river fish --- the fish are caught in the sacred currents, and the bones are returned to the water as an offering. The porridge is salted with evaporated river brine. It tastes of memory.

\textbf{What does a child learn before they can walk?} To tie the eight witness knots. Each knot represents a clause of a fair contract. The ninth knot is never tied --- it is left for the river to judge. A child who ties the ninth knot is taken to the step-pyramid and made to watch the green flame until they can recite every oath their family has ever sworn.

\textbf{What is the first thing you notice about a stranger?} Their lamp. A Lethai-ar's lamp is a public record of their witness. A stranger whose lamp burns green has kept every oath. A stranger whose lamp flickers has witnessed a broken vow. A stranger whose lamp is dark has been cut from the Vigil. Do not speak to them. The river may be listening.

\textbf{What is the most common crime?} Cord-forging --- creating a false witness cord that mimics the dye of a respected Vigil. The forger is not executed. Their hands are bound with a cord that tightens each time they speak. If they never lie, the cord will eventually loosen. If they lie once, it will cut.

\textbf{What does no one talk about but everyone knows?} The First Cord was not woven by a Lethai-ar. It was found --- already knotted, already witnessing --- floating down the river on a broken barge. No one knows who tied it. No one knows what oath it witnessed. The cord is kept in a lead-lined box at the bottom of the Well of Broken Oaths. The box is warm. The cord is still waiting.

\textbf{Where do people go when they die?} To the green flame. The dead are cremated, and their ashes are mixed with lamp oil. The oil is poured into the Vigil lanterns. The dead become the light. The light witnesses. The light does not lie.

\textbf{How do you apologize for a serious wrong?} You offer a new cord --- woven from your own hair, dyed with your own blood. The cord is tied between you and the wronged party. If the cord tightens, you are not forgiven. If it loosens, you are. If it breaks, the river will decide your fate. The river is patient.

\textbf{What is the most important festival?} The Reckoning of the Cords (spring). All witness cords that have been cut or broken are taken from the Well of Broken Oaths and burned in a barge on the river. The smoke rises. The river spirits watch. If a cord does not burn, the oath it witnessed is still unpaid. The debtor has one more year. If the same cord fails to burn twice, the Hollow comes.

\textbf{What is the secret that could destroy a powerful person?} The God-King's own witness cord is a forgery. He swore a vow to protect the Lethai-ar nine years ago. His cord was cut the next day. The Vigil knows. They have not told anyone because the cord that cut his was their own. The ninth cord that Lady Suyin cut was the king``s. She has been waiting for him to break his oath ever since. He will. They always do.
\end{almanac}

\newpage
\section{Lethai−ar (Ayokha)}
łabel{chap:lethai−ar−ayokha−revised}

\textbf{Key Demographics: } Lethai−ar
\textbf{Capital City: } Nakhon Phanom

\subsection*{Quick Reference}
\textbf{Tagline:} \textit{The cord remembers. The river witnesses. The flame judges.}

\begin{quote}
  \textbf{Vigil Lady Suyin (Elite):}  
  ``A witness cord is not a leash. It is a mirror. If you speak truly, you see yourself. If you lie, you see what you could have been. I have cut two cords in my life. I still dream of the faces in the silk. The river remembers every word spoken on its banks. We are merely the scribes.''
\end{quote}

\begin{quote}
  \textbf{Wharf−Witness (Commoner):}  
  ``My lamp has burned for nine years without being refilled. The flame is green. It flickers when a merchant lies about the weight of his silk. The river does not need oil. It needs witnesses. I sit on this wharf and watch. That is my offering.''
\end{quote}

\textbf{Genre / Mood:} Riverine legal drama, oath−magic, communal memory, spirit diplomacy.

\textbf{Starting Location:} A monsoon wharf on the great river, where silk barges are being loaded. A Vigil in gray robes holds a lantern of green flame: ``The merchant and the captain have made a vow. The captain says the silk is eight bales. The merchant says it is nine. The witness cord is tied. It will tighten around the liar's wrist. I do not need to know who is lying. The cord will tell me. But you −−− you are not bound. Watch. Learn. And do not pour the ninth cup.''

\vspace{2mm}
\begin{tcolorbox}[colback=black!3,colframe=blue!60!black,title={Theme & Atmosphere}]
The Lethai−ar of Ayokha are the living bridges between mortals and the river spirits that carry the world's promises. Their ancestors were river−folk that ascended to the realm of the fair folk, and they now walk the middle path, a step between bodies of water and the world of men. They keep lanterns that never go out, record every contract spoken on the monsoon wharves, and bind their vows with silken cords dyed in river mud. Unlike their northern kin who swore to distant deities, the Lethai−ar serve \textbf{Benzaiten} (aligned with \textit{The Witness}) or \textbf{Inari} (aligned with \textit{The Traveler}). Their vows are physical cords, braided from their own hair and dyed with river mud. When a vow is fulfilled, the cord is burned. When it is broken, the cord tightens around the oath−breaker's wrist until blood is drawn (Harm 1). The Ayokhan taboo against the ninth cup is enforced by the Lethai−ar: the ninth cup is for the river spirits, and pouring it is an invitation to the Hollow.

\vspace{2mm}
\textbf{What you notice first:} The smell of river water and oiled silk. The way every Lethai−ar carries a small oil lamp, even at noon. The constant murmur of contracts being witnessed −−− a merchant's promise, a captain's oath, a lover's vow.

\textbf{The one rule that kills newcomers:} Never pour the ninth cup of tea or sake in a Lethai−ar's presence. The ninth cup is for the witness. Pour it, and you have invited the river to judge you. The river is not merciful.
\end{tcolorbox}

\subsection*{Regional Mood: The Weight of the Witness Cord}

Power here is measured in lanterns kept lit, cords tied, and the number of oaths witnessed. The Lethai−ar have no army, no treasury, no territory. They have the river, and the river remembers everything. Every contract in Ayokha requires a Lethai−ar witness to be legally binding. The God−King's astrologers may set the calendar, but the Lethai−ar keep the chronicle. Their neutrality is absolute: they will witness a vow between enemies as readily as between allies. The price of a witness cord is not silver −−− it is a promise to keep the river clean, to light a lamp at dusk, to speak no word that cannot be unsaid.

\subsection*{Prominent NPCs of the Lethai−ar (Ayokha)}
łabel{sec:lethai−ar−npcs−revised}

\begin{longtable}{|p{0.2\textwidth}|p{0.25\textwidth}|p{0.3\textwidth}|p{0.15\textwidth}|}
ℎline
\textbf{Name} & \textbf{Role/Title} & \textbf{Motivation} & \textbf{Rumored Quirk} \\
ℎline
Vigil Lady Suyin & Master of the central wharf & Keep the river's memory free of falsehoods & She has cut nine witness cords. The ninth was her own. \\
ℎline
Cord−Maker Tama & Weaver of witness silk & Create a cord that cannot lie & She dyes her silk in river mud that glows under moonlight. \\
ℎline
Lamp−Keeper Hiro & Guardian of the undying flame & Ensure the green lamp never goes out & He has not slept in nine years. The lamp keeps him awake. \\
ℎline
River−Spirit (Nat) & Benzaiten's messenger & Collect stories from broken oaths & It appears as a woman with wet hair and riverweed. \\
ℎline
The Oath−Broken (Exile) & A Lethai−ar who cut a false cord & Earn back her place in the Vigil & She carries a cord with nine knots. The ninth is cut. \\
ℎline
\end{longtable}

\subsection*{Names of the Lethai−ar (Ayokha)}
łabel{sec:lethai−ar−names−revised}

Inspired by the river kingdoms and monsoon coasts: soft, flowing consonants, honorifics (Lady, Master). Surnames denote a river, a wharf, or a cord.

\begin{longtable}{|p{0.25\textwidth}|p{0.25\textwidth}|p{0.2\textwidth}|p{0.2\textwidth}|}
ℎline
\textbf{First Names (Male)} & \textbf{First Names (Female)} & \textbf{Clan/Epithet} & \textbf{Cord−Name} \\
ℎline
Hiro & Suyin & River−watch & \textit{the Uncut} \\
ℎline
Tama & Mira & Silk−hand & \textit{Ninth Knot} \\
ℎline
Kaito & Yuki & Lamp−bearer & \textit{Green Flame} \\
ℎline
Riku & Hana & Oath−strand & \textit{the Witness} \\
ℎline
Sora & Nami & Mud−dyer & \textit{the Unforgotten} \\
ℎline
\end{longtable}

\paragraph*{(Wharf/Temple/Cord−House)} A monsoon wharf with silk barges and green lanterns; a step−pyramid temple where the river spirit is honored; a cord−maker's house with hanging strands of dyed silk.

\section*{Spades −−− Places}
łabel{sec:lethai−ar−places−revised}
\begin{enumerate}
\setcounter{enumi}{1}
ℹtem \textbf{Monsoon Wharf (Central)} −−− A stone pier where silk is loaded and oaths are witnessed.
  ℎfill \textit{The pier is carved with the names of every oath kept. The names glow when a vow is fulfilled.} [\textbf{TRADE}] [\textbf{SOCIAL}]
ℹtem \textbf{Step−Pyramid Temple (Benzaiten)} −−− A stone pyramid rising from the riverbank. The top is a pool.
  ℎfill \textit{The pool reflects not your face but the faces of those you have wronged.} [\textbf{SPIRIT}] [\textbf{TRUTH}]
ℹtem \textbf{Cord−Maker's House} −−− A wooden building with silk strands hanging from the rafters.
  ℎfill \textit{The strands are dyed in nine shades of brown. The ninth shade is made from grave mud.} [\textbf{CRAFT}] [\textbf{MAGIC}]
ℹtem \textbf{Lamp−Keeper's Shrine} −−− A small stone hut with a single green flame.
  ℎfill \textit{The flame has not gone out in three centuries. The keeper sleeps on a bed of ashes.} [\textbf{RITUAL}] [\textbf{COMMUNITY}]
ℹtem \textbf{River Court (Floating)} −−− A barge that serves as a mobile courthouse.
  ℎfill \textit{The barge is anchored at a different spot each day. The river is the judge.} [\textbf{JUSTICE}] [\textbf{SOCIAL}]
ℹtem \textbf{Vigil's Archive} −−− A cave behind a waterfall, filled with witness cords.
  ℎfill \textit{Each cord is labeled with a date and an oath. The oldest cord is nine centuries old.} [\textbf{MEMORY}] [\textbf{LEGEND}]
ℹtem \textbf{Silk Village (Mulberry Groves)} −−− A village where silk is spun and cords are woven.
  ℎfill \textit{The mulberry trees were planted by the first Vigil. Their leaves are sacred.} [\textbf{COMMUNITY}] [\textbf{TRADE}]
ℹtem \textbf{Bridge of the Ninth Clause} −−− A wooden bridge where contracts are signed.
  ℎfill \textit{The ninth clause is always the same: ``This oath is witnessed by the river.''} [\textbf{THRESHOLD}] [\textbf{LAW}]
ℹtem \textbf{Well of Broken Oaths} −−− A dry well where cut cords are thrown.
  ℎfill \textit{The well is bottomless. The cords burn as they fall. The smoke smells of betrayal.} [\textbf{DEATH}] [\textbf{MEMORY}]
ℹtem[J] \textbf{The Oath−Broken's Hut} −−− A small shelter on the river's far bank, away from the wharves.
  ℎfill \textit{The exile lives here. She tends a single lamp. The flame is blue, not green.} [\textbf{SOCIAL}] [\textbf{MYSTERY}]
ℹtem[Q] \textbf{River−Spirit's Pool} −−− A still pool behind the temple where Benzaiten's messenger appears.
  ℎfill \textit{The pool is warm. It whispers. The whispers are the stories of broken oaths.} [\textbf{SPIRIT}] [\textbf{TRUTH}]
ℹtem[K] \textbf{The God−King's Witness Chamber} −−− A hall where the king's own oaths are recorded.
  ℎfill \textit{The walls are covered in witness cords. The king has broken nine. The cords are still tight.} [\textbf{AUTHORITY}] [\textbf{JUSTICE}]
ℹtem[A] \textbf{The First Cord (Myth)} −−− A silk strand said to be woven from the river's first current.
  ℎfill \textit{Whoever holds it can witness any oath anywhere. The price is a year of silence.} [\textbf{LEGEND}] [\textbf{TRUTH}]
\end{enumerate}

\paragraph*{(Vigil/Cord−Maker/Spirit)} Vigil Lady Suyin with a green lamp; Cord−Maker Tama with silk−dyed fingers; the River−Spirit with wet hair and riverweed.

\section*{Hearts −−− People & Factions}
łabel{sec:lethai−ar−people−revised}
\begin{enumerate}
\setcounter{enumi}{1}
ℹtem \textbf{Vigil (Lady Suyin)} −−− Gray−robed, a green lamp at her hip. She never blinks when witnessing an oath.
  ℎfill \textit{``The cord will tighten,'' she says. ``I do not need to watch. The river watches for me.''} [\textbf{JUSTICE}] [\textbf{SOCIAL}]
ℹtem \textbf{Cord−Maker (Tama)} −−− A woman with hands stained brown from river mud. She weaves silk and hair.
  ℎfill \textit{``Each cord is a promise,'' she says. ``The color tells the weight. White for truth. Red for blood. Black for death.''} [\textbf{CRAFT}] [\textbf{TRUTH}]
ℹtem \textbf{Lamp−Keeper (Hiro)} −−− An old man with ash on his face. He tends the green flame.
  ℎfill \textit{``If the flame goes out,'' he says, ``the river will forget every oath. I cannot sleep.''} [\textbf{RITUAL}] [\textbf{COMMUNITY}]
ℹtem \textbf{River−Spirit (Benzaiten's Messenger)} −−− A woman in wet robes, her hair dripping riverweed.
  ℎfill \textit{She speaks only in the voices of those who have broken oaths. ``Remember your promise,'' she whispers.} [\textbf{SPIRIT}] [\textbf{TRUTH}]
ℹtem \textbf{Merchant of the Ninth Bale} −−− A silk trader who always weighs his bales short by one.
  ℎfill \textit{He has never been caught. His witness cord is tied by a Vigil he has bribed.} [\textbf{TRADE}] [\textbf{SOCIAL}]
ℹtem \textbf{Captain of the Monsoon Ship} −−− A woman with a brass earring and a scarred hand.
  ℎfill \textit{She has never broken an oath. Her witness cord has nine knots. All are tight.} [\textbf{TRAVEL}] [\textbf{HONOR}]
ℹtem \textbf{Oath−Broken Exile} −−− A woman with a cut cord on her wrist. She sits apart from the wharf.
  ℎfill \textit{``I cut a false cord,'' she says. ``I thought the merchant was honest. He was not. Now I wait.''} [\textbf{COMMUNITY}] [\textbf{REDEMPTION}]
ℹtem \textbf{God−King's Astrologer} −−− A man who reads the monsoon and advises on oath timing.
  ℎfill \textit{``Never swear an oath during an eclipse,'' he says. ``The river cannot see you.''} [\textbf{POLITICS}] [\textbf{WEATHER}]
ℹtem \textbf{Lamp−Child (Apprentice)} −−− A young girl learning to tend the green flame.
  ℎfill \textit{She has never spoken. The lamp is her voice. The flame flickers when she is happy.} [\textbf{RITUAL}] [\textbf{INNOCENCE}]
ℹtem[J] \textbf{The Cut Cord's Ghost} −−− A translucent strand that floats through the archive at night.
  ℎfill \textit{It follows the oath−breaker. It whispers the words of the broken promise.} [\textbf{DEATH}] [\textbf{FEAR}]
ℹtem[Q] \textbf{River−Spirit's Rival (Inari's Fox)} −−− A fox with nine tails, seen at the cord−maker's house.
  ℎfill \textit{It offers a deal: a cord that cannot break, in exchange for a memory. The memory must be sweet.} [\textbf{MAGIC}] [\textbf{TRICKERY}]
ℹtem[K] \textbf{The First Vigil (Legend)} −−− A skeleton in gray robes, sitting at the bottom of the Well of Broken Oaths.
  ℎfill \textit{She holds a cord with nine knots. The ninth knot is loose. Untie it, and she will speak.} [\textbf{LEGEND}] [\textbf{DEATH}]
ℹtem[A] \textbf{Benzaiten (The River of Words)} −−− A voice heard from the step−pyramid's pool.
  ℎfill \textit{``Every word is a current,'' it says. ``Speak gently, or you will drown.''} [\textbf{PATRON}] [\textbf{TRUTH}]
\end{enumerate}

\paragraph*{(Flood/Rumor/Curse)} The river floods −−− a bridge is swept away; a rival merchant spreads rumors of cursed silk; the nat of the well is offended −−− no one may draw water until appeased.

\section*{Clubs −−− Complications/Threats}
łabel{sec:lethai−ar−complications−revised}
\begin{enumerate}
\setcounter{enumi}{1}
ℹtem \textbf{River Floods (Bridge Swept)} −−− The monsoon rises. The bridge is gone. The wharf is isolated. [\textbf{WEATHER}] [\textbf{TRAVEL}]
  ℎfill \textit{The flood is not natural. A river−spirit is angry. An oath was broken upstream.}
ℹtem \textbf{Rumor of Cursed Silk} −−− A merchant claims that a bale of silk is cursed. Trade stops. [\textbf{SOCIAL}] [\textbf{TRADE}]
  ℎfill \textit{The curse is a lie. The merchant wants to drive down prices. He will buy cheap.}
ℹtem \textbf{Well−Spirit Offended} −−− A traveler threw trash into the sacred well. The spirit refuses water. [\textbf{SPIRIT}] [\textbf{DEBT}]
  ℎfill \textit{The spirit demands an apology. The apology must be an oath. The oath must be witnessed.}
ℹtem \textbf{Witness Cord Cut (False)} −−− Someone cuts a witness cord that was tied in good faith. [\textbf{JUSTICE}] [\textbf{SOCIAL}]
  ℎfill \textit{The cutter is the Vigil's own apprentice. She thought the oath was false. It was true.}
ℹtem \textbf{Green Lamp Gutters} −−− The eternal flame flickers. The Lamp−Keeper is dying. [\textbf{RITUAL}] [\textbf{COMMUNITY}]
  ℎfill \textit{He has one day to train a successor. The successor is the party's least likely candidate.}
ℹtem \textbf{Cord−Maker's Silk Poisoned} −−− Someone dyes the silk with a solvent that weakens it. [\textbf{CRAFT}] [\textbf{TRAP}]
  ℎfill \textit{The cords will break under pressure. The poisoner is a rival cord−maker.}
ℹtem \textbf{River−Spirit Demands a Story} −−− Benzaiten's messenger appears. She asks for a broken oath. [\textbf{SPIRIT}] [\textbf{TRUTH}]
  ℎfill \textit{She will not leave until someone confesses. The confession must be true.}
ℹtem \textbf{God−King's False Oath} −−− The king swears a vow he does not intend to keep. The Vigil knows. [\textbf{POLITICS}] [\textbf{JUSTICE}]
  ℎfill \textit{If she exposes him, the king will exile the Lethai−ar. If she stays silent, the cord tightens on her wrist.}
ℹtem \textbf{Oath−Broken's Redemption Attempt} −−− The exile tries to rejoin the Vigil. She offers a true cord. [\textbf{COMMUNITY}] [\textbf{SOCIAL}]
  ℎfill \textit{The cord is her own hair. Cutting it would kill her. The Vigil must decide.}
ℹtem[J] \textbf{The Ninth Cup Poured (By Mistake)} −−− A stranger pours the ninth cup at a tea house. The Hollow comes. [\textbf{DEATH}] [\textbf{FEAR}]
  ℎfill \textit{The Hollow asks: ``Who poured this cup?'' The stranger points at the party.}
ℹtem[Q] \textbf{First Cord Found (Forgery)} −−− Someone claims to have found the mythic First Cord. It is a fake. [\textbf{LEGEND}] [\textbf{TRAP}]
  ℎfill \textit{The forger is the God−King's astrologer. He wants to control all oaths.}
ℹtem[K] \textbf{River Changes Course} −−− The river shifts overnight. The wharf is now on dry land. [\textbf{WEATHER}] [\textbf{DISASTER}]
  ℎfill \textit{The river was offended by a broken oath. It will return when the oath is fulfilled.}
ℹtem[A] \textbf{Benzaiten Speaks (Judgment)} −−− The river itself speaks from the step−pyramid pool. [\textbf{PATRON}] [\textbf{DEATH}]
  ℎfill \textit{``Nine oaths have been broken this season. The ninth was yours. Keep your word or drown.''}
\end{enumerate}

\paragraph*{(Lamp/Cord/Mirror)} Vigil lantern (social edge); silk witness cord (enforce any oath); monsoon mirror (reflects hidden intentions, one use).

\section*{Diamonds −−− Rewards/Leverage}
łabel{sec:lethai−ar−rewards−revised}
\begin{enumerate}
\setcounter{enumi}{1}
ℹtem \textbf{Vigil Lantern} −−− A small oil lamp with a green flame. Gain +1 die to social rolls in any Lethai−ar court. [\textbf{SOCIAL}]
  ℎfill \textit{The flame flickers when someone lies. It also flickers when you lie.}
ℹtem \textbf{Silk Witness Cord (Crimson)} −−− A braided cord dyed with river mud. Use it to enforce any oath. [\textbf{SOCIAL}] [\textbf{JUSTICE}]
  ℎfill \textit{The cord tightens when the oath is broken. It will not stop until blood is drawn or the promise is kept.}
ℹtem \textbf{Monsoon Mirror} −−− A polished brass disc. Reflects hidden intentions for one scene. [\textbf{TRUTH}] [\textbf{SOCIAL}]
  ℎfill \textit{Hold it before a speaker. If they lie, the mirror fogs. The fog smells of river water.}
ℹtem \textbf{Cord−Maker's Needle (Bone)} −−− A needle carved from a river−fish's spine. Tie a cord without cutting the silk. [\textbf{CRAFT}]
  ℎfill \textit{The needle can also cut a cord without making a sound. Use once.}
ℹtem \textbf{Lamp−Keeper's Ash} −−− A pouch of ash from the green flame's hearth. Throw it to calm a river−spirit. [\textbf{SPIRIT}]
  ℎfill \textit{The ash is warm. It glows faintly. The spirit will recognize it.}
ℹtem \textbf{River−Spirit's Tear (Vial)} −−− A vial of water from the step−pyramid pool. Drink it to recall a forgotten oath. [\textbf{MEMORY}] [\textbf{TRUTH}]
  ℎfill \textit{The water tastes of iron. The memory is painful. The pain is the price.}
ℹtem \textbf{Oath−Broken's Cord (Cut)} −−− A severed witness cord. Present it to demand a second chance from any Vigil. [\textbf{SOCIAL}] [\textbf{REDEMPTION}]
  ℎfill \textit{The cord will re−tie itself if the demand is just. The sound is a whisper.}
ℹtem \textbf{God−King's Witness Seal} −−− A golden seal. Commands a Lethai−ar to witness any oath, regardless of price. [\textbf{AUTHORITY}] [\textbf{SOCIAL}]
  ℎfill \textit{The seal is heavy. It was stolen from the king's chamber. He will want it back.}
ℹtem \textbf{Bridge of the Ninth Clause Map} −−− A chart showing every bridge where contracts can be witnessed. [\textbf{TRAVEL}] [\textbf{LAW}]
  ℎfill \textit{The map is written on silk. The silk is a witness cord. Do not cut it.}
ℹtem[J] \textbf{First Cord Fragment (Silk Strand)} −−− A piece of the mythic First Cord. Touch it to witness any oath without a cord. [\textbf{LEGEND}] [\textbf{JUSTICE}]
  ℎfill \textit{The fragment burns after use. The oath is still binding. The river will remember.}
ℹtem[Q] \textbf{Benzaiten's Whisper (Shell)} −−− A river shell that contains the voice of the river−spirit. [\textbf{TRUTH}] [\textbf{SPIRIT}]
  ℎfill \textit{Once, ask the river one question. The answer will be true. The river will remember the favor.}
ℹtem[K] \textbf{The Cut Cord's Ghost (Bottle)} −−− A ghost of a broken oath sealed in a glass jar. [\textbf{DEATH}] [\textbf{SOCIAL}]
  ℎfill \textit{Release it to haunt an oath−breaker. The ghost will whisper the broken promise until the oath is fulfilled.}
ℹtem[A] \textbf{The River's Pardon} −−− A scroll signed by Benzaiten herself. Voids any oath, any witness cord. [\textbf{LEGEND}] [\textbf{JUSTICE}]
  ℎfill \textit{The scroll is wet. It will never dry. The river will ask for a new oath in exchange.}
\end{enumerate}

\section*{Quick use notes}
łabel{sec:lethai−ar−quick−use−revised}
\begin{itemize}
ℹtem Draw until you have all four suits: Spade = place, Heart = actor, Club = pressure, Diamond = leverage. Highest rank sets the main Timer (2−−5 $→$ 4, 6−−10 $→$ 6, J/Q/K $→$ 8, A $→$ 10).
ℹtem Diamonds are codified outcomes (lanterns, cords, mirrors) that change Position rather than call for a roll.
ℹtem If any A appears, echo \textbf{river−omens}−−−water that flows upstream, a lamp that burns green without oil, a cord that ties itself.
\end{itemize}

\subsection*{Additional Features (Cultural Mechanics)}
łabel{sec:lethai−ar−mechanics−revised}

\begin{itemize}
ℹtem \textbf{Witness Clause:} Once per session, when you are present for an oath or contract, you may declare yourself a witness. That oath gains +1 DV to break, and if broken, the oath−breaker must grant you a personal favor (in addition to any other penalty). The witness cord glows faintly for the duration. \textit{Mechanical Note: Player gains 1 Boon if they witness a truthful oath.}
ℹtem \textbf{Lantern of Attendance:} You carry a small oil lamp (the Vigil lantern). As long as it burns, you cannot be surprised in social situations (the flame flickers when deception begins). If the lamp goes out, you lose this benefit for the scene. Refilling the lamp requires a minor offering (a prayer to Benzaiten, a coin in a river). \textit{Mechanical Note: Gain +1 Position on social rolls involving deception detection.}
ℹtem \textbf{Silk Token:} Once per arc, you may offer a witness cord as a token to a faction or NPC. While they hold it, you gain +1 die to all social rolls against them, but they may also call on you to witness one of their own oaths (refusing costs the token and 1 Fatigue). \textit{Mechanical Note: Treat as a String asset.}
ℹtem \textbf{The Ninth Cup Taboo:} In Ayokha, the ninth cup of tea or sake is never poured. Pouring it invites the Hollow to witness the conversation. The Lethai−ar enforce this taboo strictly; a character who pours the ninth cup gains the Hollow's Attention (advance the timer by 1 segment) and must succeed on a Spirit + Resolve test (DV 3) or forget the next minute of conversation. \textit{Mechanical Note: Banks 1 Story Beat for the GM.}
ℹtem \textbf{Local Patrons (Syncretic Names):}
  \begin{itemize}
    ℹtem \textbf{Benzaiten} −−− Aligned with \textit{The Witness}. \textit{The river of words; every oath flows to her.}
    ℹtem \textbf{Inari} −−− Aligned with \textit{The Traveler}. \textit{The fox's shadow; contracts are paths, and paths must be witnessed.}
    ℹtem \textbf{The Green Flame} −−− Aligned with \textit{The Sealed Gate}. \textit{The lamp that never goes out; the threshold between truth and lie.}
    ℹtem \textbf{The Ninth Current} −−− Aligned with \textit{Khemesh}. \textit{The depth beneath the words; the weight of unspoken promises.}
  \end{itemize}
\end{itemize}

\subsection*{Lore Echoes (Cross−Expansion Hooks)}
\begin{itemize}
  ℹtem \textbf{Witness Cords and the Covenant (Mykkiel):} The Covenant recognizes Lethai−ar witness cords as legally binding. A Vigil may be called to testify in a Covenant court; her cord is admissible as evidence.
  ℹtem \textbf{Silk Cords and the Sidhi (Way of Silk):} Sidhi freedmen sometimes carry Lethai−ar witness cords as proof of emancipation. A cord from a respected Vigil is worth more than a freedman's seal.
  ℹtem \textbf{The River Spirit and the Ukwe (Sekogo):} Benzaiten's river is said to flow from the Ukwe's roots. A Lethai−ar who visits the Ukwe can read the water's memory without a context key.
  ℹtem \textbf{The Ninth Cup and the Hollow (Eastern Realms):} The Hollow's Oath−Keepers are attracted to the ninth cup. Aelaerem hedge−keepers never pour a ninth cup in a Lethai−ar's presence; they pour the ninth cup into the ground instead.
\end{itemize}

\subsection*{Lethai−ar Superstitions (burn for +1 die once)}
\begin{itemize}
  ℹtem \textbf{Bread to the River:} Crumble a crust into the current. Ignore the first \emph{River Floods} penalty. [\textbf{SUPERSTITION}]
  ℹtem \textbf{Knot the Cord:} Tie a single knot in a witness cord before using it. Cancel \emph{Witness Cord Cut} or take −1 die on your next oath (your choice). [\textbf{SUPERSTITION}]
  ℹtem \textbf{Name to the Lamp:} Whisper a dead Vigil's name into the green flame. Gain +1 Effect when witnessing an oath this scene, but mark \emph{Obligation} to Benzaiten. [\textbf{SUPERSTITION}]
\end{itemize}

\subsection*{Thematic SB Spend Table}
łabel{sec:lethai−ar−sb−revised}

\paragraph{Minor Complications (1 SB)}
\begin{itemize}
ℹtem \textbf{Exposure:} Your actions draw unwanted attention from \textbf{Vigils or river−spirits}.
ℹtem \textbf{Noise:} Sounds of your actions alert nearby \textbf{lamp−keepers or cord−makers}.
ℹtem \textbf{Trace:} Evidence of your passage marks your route for \textbf{oath−breakers or the Hollow}.
ℹtem \textbf{Delay:} A brief but meaningful setback costs you \textbf{time or standing with the Vigil}.
ℹtem \textbf{Supply Strain:} Mark +1 segment on a relevant \textbf{resource timer}.
\end{itemize}

\paragraph{Moderate Setbacks (2 SB)}
\begin{itemize}
ℹtem \textbf{Alarm Raised:} \textbf{Local Vigil or wharf−master} becomes aware and begins responding.
ℹtem \textbf{Position Lost:} You lose advantageous ground/cover/stealth due to \textbf{rising river or lantern failure}.
ℹtem \textbf{Foe Appears:} An \textbf{oath−breaker's ghost or a rival cord−maker} arrives on scene.
ℹtem \textbf{Gear Trouble:} A piece of equipment becomes \textbf{Compromised/Neglected}.
ℹtem \textbf{Lock/Barrier:} A simple obstacle now requires a test to overcome.
\end{itemize}

\paragraph{Serious Trouble (3 SB)}
\begin{itemize}
ℹtem \textbf{Reinforcements:} Additional \textbf{Vigils, river−spirits, or Hollow agents} arrive.
ℹtem \textbf{Key Gear Breaks:} A crucial tool/weapon becomes temporarily unusable.
ℹtem \textbf{Major Twist:} The situation fundamentally changes−−−\textbf{cord cut/lamp out/river changes course}.
ℹtem \textbf{Rail Tick:} Advance a relevant campaign/front timer by 1 segment.
ℹtem \textbf{Condition Applied:} Mark \textbf{Fatigue 1/Harm 1/Condition} appropriate to fiction.
\end{itemize}

\paragraph{Major Turns (4+ SB)}
\begin{itemize}
ℹtem \textbf{Trap Springs:} A prepared danger activates with full effect.
ℹtem \textbf{Authority Arrival:} \textbf{Vigil Lady Suyin, River−Spirit, or Benzaiten's voice} intervenes.
ℹtem \textbf{Scene Shift:} The environment changes dramatically−−−\textbf{wharf floods/temple cracks/cords tighten everywhere}.
ℹtem \textbf{Patron Omen:} Divine/arcane forces take notice−−−\textbf{omen appears/blessing lost/curse manifests}.
ℹtem \textbf{Narrative Pivot:} The story takes an unexpected turn that reframes objectives.
\end{itemize}

\paragraph{Region−Specific SB Options}
\begin{itemize}
ℹtem \textbf{Lethai−ar (Witness Law):} Green flames flicker in unison, cords glow, a river's current reverses.
ℹtem \textbf{Lethai−ar (Cord Law):} A witness cord tightens without cause, a knot unties itself, a cut cord re−knots.
ℹtem \textbf{Lethai−ar (River Law):} Water levels rise without rain, a fish speaks a name, the ninth cup pours itself.
\end{itemize}

\subsection*{Pursuit Ladder (Near / Mid / Far)}
\begin{itemize}
ℹtem \textbf{Advance/Withdraw (Wits + Stealth):} On a hit, shift one band. On a strong hit, also impose \emph{Lantern Glare}.
ℹtem \textbf{Cut the Cord (Presence + Sway):} On a hit, freeze range for one exchange; if a witness cord is in your possession, +1 die.
ℹtem \textbf{Weather the Flood (Body + Athletics):} On a hit in \emph{River Floods}, you pick who reaches high ground; on a miss, the GM does.
\end{itemize}

\subsection*{Boss Seeds (Cord & River)}
\begin{itemize}
ℹtem \textbf{The Cord−Cutter (Nine Broken Oaths)} −−− \emph{Phases}: Witness Cords Severed $→$ Orchestrated Perjury $→$ Vigil Coup. \emph{Cracks}: re−tie a cut cord, expose a false Vigil, survive a night in the Well of Broken Oaths.
ℹtem \textbf{The River−Dryer (Benzaiten's Wrath)} −−− \emph{Phases}: Water Levels Dropping $→$ Procession of Empty Wharves $→$ Flood−Baptism of a Merchant. \emph{Cracks}: relic river water, survivor testimony, the green flame burning without oil.
\end{itemize}

\begin{tcolorbox}[colback=black!3,colframe=black!40!white,title={Patronage & Power}]
Among the Lethai−ar of Ayokha, power flows through witness cords, lantern flames, and the memory of the river. A Vigil's authority rests on the number of oaths she has witnessed and never broken. The lamp is the symbol of that authority; a Vigil without a lamp is not a Vigil.

\textbf{For the GM:}  
Power in Lethai−ar society is measured in cords tied, lamps kept lit, and oaths witnessed. Patronage often takes the form of lanterns, witness cords, or mirrors −−− each tied to a specific Vigil, wharf, or contract. To emphasize the cultural mechanics:
\begin{itemize}
ℹtem Use the \textbf{Witness Clause} to create social dilemmas: should the party witness an oath that might be broken?
ℹtem Enforce the \textbf{Lantern of Attendance} −−− a party member who carries a Vigil lantern gains social advantages but must maintain the flame.
ℹtem Let \textbf{Silk Token} be a form of communal currency: a witness cord given to a faction creates a mutual obligation.
ℹtem Track the \textbf{Ninth Cup Taboo} as a source of Hollow Attention. A single mistake can summon an Oath−Keeper.
\end{itemize}
In Lethai−ar society, the river forgets nothing −−− and the witness cord remembers every word.
\end{tcolorbox}

% Index entries
∈dex{Lethai−ar (Ayokha)|see{River−Bound Oath−Keepers}}
∈dex{Theme generator}
∈dex{River−Bound Oath−Keepers generator}
∈dex{Places!Lethai−ar}
∈dex{People!Lethai−ar}
∈dex{Complications!Lethai−ar}
∈dex{Rewards!Lethai−ar}
∈dex{Witness cords}
∈dex{Vigil lantern}
∈dex{Silk token}
∈dex{Ninth cup taboo}
∈dex{Benzaiten}
∈dex{River spirits}
∈dex{Monsoon wharves}
∈dex{Oath−Broken exile}
∈dex{Local Patrons!Lethai−ar}
\newpage

\chapter{Lethai-thora of Sihai}
\label{chap:lethai-thora}

\emph{The Silent Ledgers.}

\noindent
{\Large
\textit{We do not burn the dangerous scrolls. We bury them in plain sight, on shelves where no one thinks to look. The Emperor's Censor once stood in our archive and asked for a tax ledger. We gave him one. He did not notice the shelf behind him held the true history of his own grandfather's death. Silence is not emptiness. It is a door that opens only inward.}

\vspace{0.5em}

My grandmother tied this scroll with a silk cord. The cord has nine knots. The ninth knot is a secret she never told me. When I die, I will tie a tenth. Someone will untie it in a century. They will learn why the emperor's concubine wept on the night of the lantern festival. Some secrets are too heavy for one lifetime.

\vspace{0.5em}

The Lethai-thora were once courtiers of a celestial fey realm that spanned the Valewood and the eastern hills. When the old fey courts fell --- when the Alder King withdrew and the Hazel Queen sealed her gates --- many Lethai-thora fled east. They crossed the mountains with nothing but armfuls of scrolls and the memory of a civilisation that had forgotten how to die.

In the fertile river valleys of what would become Sihai, they found scattered agrarian tribes fighting over water and silence. The Lethai-thora did not conquer. They taught. They taught the tribes to write contracts instead of drawing blood, to measure grain instead of stealing it, to record debts instead of nursing grudges. Over generations, the tribes became a kingdom, the kingdom became an empire, and the empire became Sihai. The Lethai-thora became its memory.

Now they dwell in hidden pavilions, preserving ''forgotten ledgers`` that record what the Mandate of Heaven prefers to forget. Their knowledge is a burden: each secret they preserve slowly erases a piece of themselves. A Lethai-thora who guards a forbidden truth may find that she can no longer recall the color of her mother's dress --- a small price, she says, for keeping the ledger whole. They are no longer fey. They are no longer fully human. They are something in between: scribes of a story that has no ending, archivists of a past that never quite happened.

\vspace{0.5em}

Never touch a scroll without announcing your lineage. Even a false lineage is better than silence. The scrolls remember who opened them, and they will whisper your name to the Hollow.

\vspace{1em}

\noindent\textit{--- Keeper Lian, who has not spoken aloud in twenty years --- she writes on water}
}

\vspace{1em}

\begin{tcolorbox}[colback=black!3,colframe=blue!60!black,title={Theme \& Atmosphere},breakable]
\textbf{What you notice first:} The smell of old parchment and incense. The way scrolls are stacked with their spines facing the wall --- so the words cannot see who reads them. The silence when a secret is spoken; even the insects stop chirping.

\textbf{The one rule that kills newcomers:} Never touch a scroll without announcing your lineage. Even a false lineage is better than silence. The scrolls remember who opened them, and they will whisper your name to the Hollow.
\end{tcolorbox}

\subsection*{Regional Mood: The Weight of the Forgotten}

Power here is measured in scrolls sealed, secrets kept, and the length of a silence. The Lethai-thora have no king, no army, no treasury. They have archives, and the archives are their fortress. The Emperor's Censorate fears them because they cannot audit a secret that has no ink. The Mandate of Heaven is written in public; the Lethai-thora keep the unwritten ledger --- the one that records why dynasties fall. They are not rebels; they are witnesses. And a witness who never speaks is the most dangerous kind.

\subsection*{Starting Location}

A hidden pavilion behind a waterfall, where scrolls are stacked to the ceiling and the air is dry as bone. A Keeper of the Seventh Gutter looks up from a brass lamp: ''The Emperor's Quiet Chancellor has sent a messenger. He wants a scroll that does not exist. He wants the truth about his mother's death. The truth is written in water. It will vanish if we speak it. But you --- you are not bound by our silence. Will you carry a memory for us?``

\subsection*{The Almanac of Ordinary Days (Lethai-thora)}
\begin{almanac}
\textbf{What do people eat for breakfast?} Rice porridge with dried fish --- the fish are smoked over fires fed with old scrolls. The smoke smells of forgotten contracts. The porridge is always lukewarm. The archives are never truly warm.

\textbf{What does a child learn before they can walk?} To read without moving their lips. A Lethai-thora child who whispers a secret aloud is sent to sit in the Chamber of the Ninth Scroll until they learn that silence is a shield. Some children never come back. They become scribes of the empty shelves.

\textbf{What is the first thing you notice about a stranger?} Their ink-stained fingers. A Lethai-thora has ink under every fingernail. A stranger whose fingers are clean has never copied a forbidden text. A stranger whose fingers are stained black has copied too many. A stranger whose fingers are stained red has copied a death warrant --- their own.

\textbf{What is the most common crime?} Context key theft --- stealing the carved stone that unlocks a scroll's meaning. The thief is not punished by law. They are given a blank scroll and told to write their own confession. The scroll will fill itself if the thief lies. The ink is their own blood.

\textbf{What does no one talk about but everyone knows?} The Ninth Scroll is not a scroll. It is a memory so dangerous that nine different keepers each sealed a fragment of it in their own minds. When the ninth keeper dies, the memory will be lost forever --- or will escape.

\textbf{Where do people go when they die?} To the Well of Forgotten Names. The dead are not buried. Their bodies are placed in the well, and the well swallows them. The water turns black. The black is the ink of their unwritten secrets. The well is bottomless. It is never full.

\textbf{How do you apologize for a serious wrong?} You offer a silence --- a period of total stillness, no breath, no heartbeat, no word. The silence is timed by a water timer. If you hold the silence for nine drops, you are forgiven. If you break it, you must speak the wrong aloud in the Chamber of the Ninth Scroll. The scroll will remember. The scroll does not forgive.

\textbf{What is the most important festival?} The Unrolling of the Ledgers (winter). Once a year, the keepers unroll a single forbidden scroll and read it aloud to each other --- but they read it in a whisper, and only the walls hear. The walls remember. The walls do not judge. The festival ends with all nine keepers drinking ink from a silver cup. The ink is their own memory. The next year, they will have forgotten.

\textbf{What is the secret that could destroy a powerful person?} The emperor's true lineage is recorded on a scroll that no living keeper has opened. The scroll is sealed with a cord tied by the first keeper. The cord has nine knots. The ninth knot is a lie --- the emperor's grandfather was not of the blood. The empire was founded by a concubine's son. The scroll is warm. It is kept in the First Pavilion. The First Pavilion is invisible. Only those who have forgotten nine names can see it.
\end{almanac}

\newpage
\section{Lethai−thora −−− ``The Silent Ledgers of Sihai''}
łabel{chap:lethai−thora}

\textbf{Key Demographics: } Lethai−thoraa
\textbf{Capital City: } Linshu

\subsection*{Quick Reference}
\textbf{Tagline:} \textit{We do not burn the dangerous scrolls. We bury them on shelves where no one thinks to look. Silence is not emptiness. It is a door that opens only inward.}

\begin{quote}
  \textbf{\textit{Keeper of the Seventh Gutter (Elite):}}  
  \textit{``We do not burn the dangerous scrolls. We bury them in plain sight, on shelves where no one thinks to look. The Emperor's Censor once stood in our archive and asked for a tax ledger. We gave him one. He did not notice the shelf behind him held the true history of his own grandfather's death. Silence is not emptiness. It is a door that opens only inward.''}
\end{quote}

\begin{quote}
  \textbf{\textit{Scroll−Binder (Commoner):}}  
  \textit{``My grandmother tied this scroll with a silk cord. The cord has nine knots. The ninth knot is a secret she never told me. When I die, I will tie a tenth. Someone will untie it in a century. They will learn why the emperor's concubine wept on the night of the lantern festival. Some secrets are too heavy for one lifetime.''}
\end{quote}

\textbf{Genre / Mood:} Bureaucratic melancholy, hidden history, silence as power.

\textbf{Starting Location:} A hidden pavilion behind a waterfall, where scrolls are stacked to the ceiling and the air is dry as bone. A Keeper of the Seventh Gutter looks up from a brass lamp: \textit{``The Emperor's Quiet Chancellor has sent a messenger. He wants a scroll that does not exist. He wants the truth about his mother's death. The truth is written in water. It will vanish if we speak it. But you — you are not bound by our silence. Will you carry a memory for us?''}

\vspace{2mm}
\begin{tcolorbox}[colback=black!3,colframe=blue!60!black,title={Theme & Atmosphere}]
The Lethai−thora were once courtiers of a celestial fey realm that spanned the Valewood and the eastern hills. When the old fey courts fell — when the Alder King withdrew and the Hazel Queen sealed her gates — many Lethai−thora fled east. They crossed the mountains with nothing but armfuls of scrolls and the memory of a civilisation that had forgotten how to die.

In the fertile river valleys of what would become Sihai, they found scattered agrarian tribes fighting over water and silence. The Lethai−thora did not conquer. They taught. They taught the tribes to write contracts instead of drawing blood, to measure grain instead of stealing it, to record debts instead of nursing grudges. Over generations, the tribes became a kingdom, the kingdom became an empire, and the empire became Sihai. The Lethai−thora became its memory.

Now they dwell in hidden pavilions, preserving \textit{forgotten ledgers} that record what the Mandate of Heaven prefers to forget. Their knowledge is a burden: each secret they preserve slowly erases a piece of themselves. An Lethai−thora who guards a forbidden truth may find that she can no longer recall the colour of her mother's dress — a small price, she says, for keeping the ledger whole. They are no longer fey. They are no longer fully human. They are something in between: scribes of a story that has no ending, archivists of a past that never quite happened.

\vspace{2mm}
\textbf{What you notice first:} The smell of old parchment and incense. The way scrolls are stacked with their spines facing the wall — so the words cannot see who reads them. The silence when a secret is spoken; even the insects stop chirping.

\textbf{The one rule that kills newcomers:} Never touch a scroll without announcing your lineage. Even a false lineage is better than silence. The scrolls remember who opened them, and they will whisper your name to the Hollow.
\end{tcolorbox}

\subsection*{Regional Mood: The Weight of the Forgotten}

Power here is measured in scrolls sealed, secrets kept, and the length of a silence. The Lethai−thora have no king, no army, no treasury. They have archives, and the archives are their fortress. The Emperor's Censorate fears them because they cannot audit a secret that has no ink. The Mandate of Heaven is written in public; the Lethai−thora keep the unwritten ledger — the one that records why dynasties fall. They are not rebels; they are witnesses. And a witness who never speaks is the most dangerous kind.

\subsection*{The Fledging of the Empire (Origin Story)}
łabel{sec:lethai−thora−origin}

When the Valewood's fey courts crumbled, a caravan of Lethai−thora refugees crossed the eastern passes. Their scrolls were water−stained, their children hungry, their memories already fraying. They found a patchwork of warring clans along the Sihon River — farmers who fought over silt, herders who stole each other's sheep, chieftains who could not agree on a name for the season.

The Lethai−thora did not offer swords. They offered ink.

They taught the clans to write their disputes on clay, to seal them with witnesses, to let the record speak instead of the spear. A generation later, the first emperor rose — a man who had learned his letters from a Lethai−thora tutor. He did not forget. He built his capital around the Lethai−thora archive, and he decreed that no law could be written without a copy placed in their care.

The Lethai−thora have been part of Sihai ever since — embedded so deeply that they can no longer imagine leaving, even if they wanted to. They are the empire's memory, but also its conscience, its secret‑keeper, and its silent judge. Some among them resent this entanglement; they dream of returning to the Valewood, of hearing the old songs again. But the Valewood has changed. The fey courts do not answer. And the Lethai−thora have become something they never intended: the architects of a mortal empire.

\subsection*{Prominent NPCs of the Lethai−thora}
łabel{sec:lethai−thora−npcs}

\begin{longtable}{|p{0.2\textwidth}|p{0.25\textwidth}|p{0.3\textwidth}|p{0.15\textwidth}|}
ℎline
\textbf{Name} & \textbf{Role/Title} & \textbf{Motivation} & \textbf{Rumoured Quirk} \\
ℎline
Keeper Lian & Master of the Seventh Gutter & Preserve the true history of the emperor's grandfather & She has not spoken aloud in twenty years. She writes on water. \\
ℎline
Scribe Wei & Scroll−binder & Tie a secret so tightly it can never be untied & He has forgotten his own name. He signs his work with a thumbprint. \\
ℎline
Ink−Maker Hana & Forger of context keys & Create inks that reveal different truths in different seasons & She has a jar of ink made from crushed moonstones. It only writes during eclipses. \\
ℎline
The Quiet Chancellor & Imperial censor (secret ally) & Find the scroll that proves the emperor's claim is false & He speaks in couplets; the second couplet is always a bribe. \\
ℎline
The Ghost of the Ninth Scroll & A memory that was never written & A secret that killed its keeper before it could be recorded & It appears as a blank scroll that weeps ink. The ink is the secret. \\
ℎline
\end{longtable}

\subsection*{Names of the Lethai−thora}
łabel{sec:lethai−thora−names}

Inspired by the central kingdom traditions: soft consonants, vowel endings, honorifics. Surnames denote an archive, a gutter, or a lineage.

\begin{longtable}{|p{0.25\textwidth}|p{0.25\textwidth}|p{0.2\textwidth}|p{0.2\textwidth}|}
ℎline
\textbf{First Names (Male)} & \textbf{First Names (Female)} & \textbf{Clan/House} & \textbf{Epithets} \\
ℎline
Lian & Hana & Seventh Gutter & \textit{the Silent} \\
ℎline
Wei & Mira & Forbidden Archive & \textit{the Unwitnessed} \\
ℎline
Chen & Nuan & Ink−Thread & \textit{Ninth Scroll} \\
ℎline
Jin & Yue & Water−Pavilion & \textit{the Forgetful} \\
ℎline
Liam & Sora & Root−Memory & \textit{the Context−Keeper} \\
ℎline
\end{longtable}

\paragraph*{(Archive/Chamber/Shrine)} A hidden pavilion with scrolls stacked to the ceiling; a brass lamp that never goes out; a shrine to the forgotten gods, covered in dust.

\section*{\spadecard{J} Places}
łabel{sec:lethai−thora−places}
\begin{enumerate}
\setcounter{enumi}{1}
ℹtem \textbf{Hidden Pavilion (Waterfall)} −−− An archive behind a curtain of falling water.
  ℎfill \textit{The water muffles sound. Secrets spoken here cannot be overheard — but the water remembers.} \texttt{[SECRECY]}
ℹtem \textbf{Chamber of the Ninth Scroll} −−− A locked room with a single blank scroll on a pedestal.
  ℎfill \textit{The scroll is not blank. It is written in disappearing ink. The reagent is moonlight.} \texttt{[MYSTERY]}
ℹtem \textbf{Well of Forgotten Names} −−− A stone well where the Lethai−thora throw scrolls that are too dangerous to keep.
  ℎfill \textit{The well is bottomless. The scrolls burn as they fall. The smoke forms faces.} \texttt{[DEATH]}
ℹtem \textbf{Ink−Maker's Studio} −−− A cave with grinding stones and jars of pigment.
  ℎfill \textit{The ink is made from crushed beetles, iron ore, and tears. The tears are the ink−maker's own.} \texttt{[CRAFT]}
ℹtem \textbf{Silent Scriptorium} −−− A hall where scribes copy scrolls without speaking.
  ℎfill \textit{They write with their left hands. Their right hands are for swearing oaths.} \texttt{[LABOR]}
ℹtem \textbf{Context Key Archive} −−− A vault of stones, each carved with a season and a place.
  ℎfill \textit{A scroll cannot be read without the correct context key. The keys are scattered across Sihai.} \texttt{[TRAVEL]}
ℹtem \textbf{Emperor's Censorate (Monitoring Station)} −−− A building near the archive, staffed by imperial spies.
  ℎfill \textit{They pretend not to know about the hidden pavilion. The Lethai−thora pretend not to notice them.} \texttt{[AUTHORITY]}
ℹtem \textbf{Bridge of Unspoken Names} −−− A stone arch where the Lethai−thora whisper secrets to the river.
  ℎfill \textit{The river carries the secrets to the sea. The sea does not tell.} \texttt{[RITUAL]}
ℹtem \textbf{Shrine of the Silent Gods} −−− A ruined temple where the old fey beings were once worshipped.
  ℎfill \textit{The gods are not dead. They are waiting. The Lethai−thora leave offerings of ink.} \texttt{[LEGEND]}
ℹtem[J] \textbf{The Forgotten Ledger's Vault} −−− A subterranean room with a single iron door.
  ℎfill \textit{The door has nine locks. The keys are held by nine different keepers. The ninth keeper is dead.} \texttt{[SECRECY]}
ℹtem[Q] \textbf{Memory−Thread Chamber} −−− A room where Lethai−thora weave their hair into cords.
  ℎfill \textit{Each thread contains a memory. Cutting a thread releases the memory — and kills the keeper.} \texttt{[DEATH]}
ℹtem[K] \textbf{The Emperor's Secret Archive (Hidden)} −−− A duplicate archive built by a paranoid emperor.
  ℎfill \textit{It contains the same scrolls as the Lethai−thora's, plus the emperor's own confessions.} \texttt{[AUTHORITY]}
ℹtem[A] \textbf{The First Pavilion (Myth)} −−− Where the first Lethai−thora hid the first forbidden scroll.
  ℎfill \textit{It is said to be invisible. Only those who have forgotten nine names can see it.} \texttt{[LEGEND]}
\end{enumerate}

\paragraph*{(Keeper/Scribe/Chancellor)} Keeper Lian with a brass lamp; Scribe Wei with ink−stained fingers; the Quiet Chancellor with a sealed letter.

\section*{ℎeartcard{Q} People & Factions}
łabel{sec:lethai−thora−people}
\begin{enumerate}
\setcounter{enumi}{1}
ℹtem \textbf{Keeper of the Seventh Gutter} −−− Old, silent, writing on water with a reed brush.
  ℎfill \textit{``The words vanish,''} she says. \textit{``That is how I know they are true.''} \texttt{[JUSTICE]}
ℹtem \textbf{Scroll−Binder (Wei)} −−− A man with no shadow. He lost it to a scroll that ate his reflection.
  ℎfill \textit{He ties silk cords around each scroll. The knots are the secrets.} \texttt{[CRAFT]}
ℹtem \textbf{Ink−Maker (Hana)} −−− A woman with stained hands and a jar of moonstone ink.
  ℎfill \textit{``This ink only writes during an eclipse,''} she says. \textit{``The next eclipse is in nine days.''} \texttt{[MAGIC]}
ℹtem \textbf{Quiet Chancellor} −−− An imperial censor in grey robes. He speaks in couplets.
  ℎfill \textit{The first couplet is a greeting. The second is a threat. The third is a bribe.} \texttt{[POLITICS]}
ℹtem \textbf{Memory−Weaver} −−− A woman with a loom made of human hair. She weaves memories into cloth.
  ℎfill \textit{The cloth is soft. It whispers the memory when you touch it.} \texttt{[ART]}
ℹtem \textbf{Imperial Spy (Undercover)} −−− A young man posing as a scribe. He has a hidden seal.
  ℎfill \textit{He is looking for a scroll that names the emperor's true father. He will pay.} \texttt{[DECEPTION]}
ℹtem \textbf{Ghost of the Ninth Scroll} −−− A blank scroll that floats through the archive at midnight.
  ℎfill \textit{It weeps ink. The ink forms words. The words are a secret that killed its keeper.} \texttt{[DEATH]}
ℹtem \textbf{Forbidden Scholar} −−− A monk who was excommunicated for reading a scroll he should not have.
  ℎfill \textit{His eyes are gone. He reads by touch. The scrolls burn his fingers.} \texttt{[TRAGEDY]}
ℹtem \textbf{Silent Guard} −−− A warrior who took a vow of silence to protect the archive.
  ℎfill \textit{He communicates by writing on his own skin. His skin is full of scars.} \texttt{[COMBAT]}
ℹtem[J] \textbf{The First Keeper (Ghost)} −−− A translucent figure in scholar's robes, seen at the well.
  ℎfill \textit{She asks: ``What secret do you wish to forget?'' The answer becomes a scroll.} \texttt{[MYSTERY]}
ℹtem[Q] \textbf{The Censorate's Interrogator} −−− A man with a brass scale. He weighs truth and lies.
  ℎfill \textit{His scale is balanced. He has never found a truth heavy enough.} \texttt{[AUTHORITY]}
ℹtem[K] \textbf{The Emperor's Dreamer} −−− A woman who dreams the emperor's nightmares. She knows what he fears.
  ℎfill \textit{She is kept in a locked room. The Lethai−thora want to free her.} \texttt{[MAGIC]}
ℹtem[A] \textbf{The Forgotten God's Echo} −−− A voice from the ruined shrine. It speaks in a language no one understands.
  ℎfill \textit{The Lethai−thora say it is the last fey being. It is asking for a memory.} \texttt{[LEGEND]}
\end{enumerate}

\paragraph*{(Scroll/Forgery/Witness)} A scroll is found to be a forgery — the party is blamed; a context key is missing — the scroll cannot be read; the Quiet Chancellor demands a witness to a secret.

\section*{\clubcard{K} Complications/Threats}
łabel{sec:lethai−thora−complications}
\begin{enumerate}
\setcounter{enumi}{1}
ℹtem \textbf{Forged Scroll} −−− A tax ledger is discovered to be a fake. The forgery is perfect. \texttt{[DECEPTION]}
  ℎfill \textit{The forger is a Lethai−thora. The scroll is not a forgery — it is the original. The empire's copy is wrong.}
ℹtem \textbf{Missing Context Key} −−− A scroll cannot be read without a specific stone. The stone is lost. \texttt{[TRAVEL]}
  ℎfill \textit{The stone was stolen by the Censorate. They want to destroy the scroll.}
ℹtem \textbf{Witness Demanded} −−− The Quiet Chancellor asks the party to witness a secret exchange. \texttt{[SOCIAL]}
  ℎfill \textit{If they agree, they become bound by the Lethai−thora's silence. If they refuse, they are exiled from the archive.}
ℹtem \textbf{Archive Fire (Accident)} −−− A lamp falls. Scrolls burn. The smoke forms names. \texttt{[DISASTER]}
  ℎfill \textit{The names are those of the keepers. One is yours.}
ℹtem \textbf{Censorate Raid} −−− Imperial inspectors arrive with a warrant. They have a list of forbidden scrolls. \texttt{[AUTHORITY]}
  ℎfill \textit{The list is accurate. Someone inside the archive is a spy.}
ℹtem \textbf{Memory−Thread Cut} −−− A keeper's memory−thread is severed. She forgets everything. \texttt{[DEATH]}
  ℎfill \textit{The thread was cut by an assassin. The assassin wants a secret only she knew.}
ℹtem \textbf{Silence Broken} −−− A keeper speaks aloud for the first time in decades. The Hollow hears. \texttt{[DEATH]}
  ℎfill \textit{The Hollow sends an Oath−Keeper. The Oath−Keeper asks: ``What did you say?''}
ℹtem \textbf{Ink of Truth (Forced Confession)} −−− Someone poisons the ink well. Writing with it forces you to confess. \texttt{[MAGIC]}
  ℎfill \textit{The poison was placed by the Censorate. They want to know where the forbidden scroll is.}
ℹtem \textbf{Emperor's Death (Succession Crisis)} −−− The emperor dies. The Mandate wavers. The archive is vulnerable. \texttt{[POLITICS]}
  ℎfill \textit{The new emperor will burn the scrolls. The Lethai−thora must hide them — or die.}
ℹtem[J] \textbf{Ghost of the Ninth Scroll Awakens} −−− The blank scroll fills with ink. The ink is a name. \texttt{[DEATH]}
  ℎfill \textit{The name is the next person to die. The name is yours.}
ℹtem[Q] \textbf{The First Pavilion Appears} −−− The mythic archive manifests. It is invisible to most. The party can see it. \texttt{[MYSTERY]}
  ℎfill \textit{They are now witnesses to a secret that has been hidden for a thousand years. The Hollow will want it.}
ℹtem[K] \textbf{Censorate's Interrogation} −−− The party is arrested and brought before the brass scale. \texttt{[AUTHORITY]}
  ℎfill \textit{The interrogator asks one question: ``What did the Lethai−thora tell you?'' The scale will know if you lie.}
ℹtem[A] \textbf{The Forgotten God Wakes} −−− The voice from the shrine becomes a scream. The ground shakes. \texttt{[LEGEND]}
  ℎfill \textit{The god wants its memory back. The memory is held in a scroll. The scroll is the Ninth Scroll.}
\end{enumerate}

\paragraph*{(Scroll/Stone/Cord)} Forbidden scroll (one truth no one wants spoken); context key stone (unlocks one scroll's meaning); memory−thread (a cord containing a secret).

\section*{⋄card{A} Rewards/Leverage}
łabel{sec:lethai−thora−rewards}
\begin{enumerate}
\setcounter{enumi}{1}
ℹtem \textbf{Forbidden Scroll} −−− A scroll sealed with silk cord. It contains one truth no one wants spoken. \texttt{[TRUTH]}
  ℎfill \textit{Read it aloud, and the secret becomes known. The speaker will be hunted.}
ℹtem \textbf{Context Key Stone} −−− A carved stone. Unlocks the meaning of one specific scroll. \texttt{[TRAVEL]}
  ℎfill \textit{The stone is warm. It glows when held near the correct scroll.}
ℹtem \textbf{Memory−Thread (Silk Cord)} −−− A cord woven from a keeper's hair. Cut it to release a memory. \texttt{[MEMORY]}
  ℎfill \textit{The memory is a secret. The keeper will forget it. The cutter will remember it.}
ℹtem \textbf{Ink−Stone (Moonstone)} −−− A jar of ink made from crushed moonstones. Writes only during eclipses. \texttt{[MAGIC]}
  ℎfill \textit{The ink reveals truths that cannot be spoken. Use it to write a secret that will last.}
ℹtem \textbf{Silence Token (Brass)} −−− A small brass charm. Present it to demand a moment of silence in any court. \texttt{[SOCIAL]}
  ℎfill \textit{During the silence, no one can lie. The truth will echo.}
ℹtem \textbf{Keeper's Lamp} −−− A brass lamp that never goes out. The flame flickers when a lie is spoken nearby. \texttt{[JUSTICE]}
  ℎfill \textit{The lamp also flickers when a secret is hidden in the room.}
ℹtem \textbf{Quiet Chancellor's Seal} −−− A jade seal. Press it to any document to grant it imperial immunity from censorship. \texttt{[AUTHORITY]}
  ℎfill \textit{The seal is a forgery. The Chancellor made it himself. It still works.}
ℹtem \textbf{Scroll of Forgotten Names} −−− A scroll that lists the names of everyone who has been erased from history. \texttt{[TRUTH]}
  ℎfill \textit{Speak a name, and that person's deeds become known. The Hollow will remember them too.}
ℹtem \textbf{Ink−Maker's Grinding Stone} −−− A stone bowl. Grind any pigment into ink that writes the truth. \texttt{[CRAFT]}
  ℎfill \textit{The truth cannot be erased. The ink will burn through lies.}
ℹtem[J] \textbf{Ghost's Tear (Vial)} −−− A vial of ink wept by the Ghost of the Ninth Scroll. \texttt{[DEATH]}
  ℎfill \textit{Write a name with it, and that person will forget one secret. The secret is your choice.}
ℹtem[Q] \textbf{First Pavilion's Key (Invisible)} −−− An invisible key that opens the mythic archive. \texttt{[LEGEND]}
  ℎfill \textit{You cannot see the key. You can feel it. It is cold.}
ℹtem[K] \textbf{Emperor's Nightmare (Scroll)} −−− A scroll containing the emperor's most feared secret. \texttt{[TRUTH]}
  ℎfill \textit{Read it aloud to make the emperor confess — or to drive him mad.}
ℹtem[A] \textbf{Forgotten God's Memory (Glass Shard)} −−− A piece of green glass that hums. It contains a fragment of a fey god's memory. \texttt{[LEGEND]}
  ℎfill \textit{Once, touch it to learn a secret of the old world. The secret will cost a memory of equal weight.}
\end{enumerate}

\section*{Quick use notes}
łabel{sec:lethai−thora−quick−use}
\begin{itemize}
ℹtem Draw until you have all four suits: Spade = place, Heart = actor, Club = pressure, Diamond = leverage. Highest rank sets the main Timer (2–5 \(→\) 4, 6–10 \(→\) 6, J/Q/K \(→\) 8, A \(→\) 10).
ℹtem Diamonds are codified outcomes (scrolls, stones, threads) that change position rather than call for a roll.
ℹtem If any A appears, echo \textbf{ink−omens}—scrolls that rustle in still air, a lamp that burns green, a forgotten name whispered from a shelf.
\end{itemize}

\subsection*{Additional Features (Cultural Mechanics)}
łabel{sec:lethai−thora−mechanics}

\begin{itemize}
ℹtem \textbf{Silence as Sanctuary:} Once per scene, when asked a direct question, a character may spend 1 Boon to answer with silence. Treat Position as Dominant for resisting interrogation — but the GM gains 1 Story Beat (the silence has weight). In Lethai−thora halls, a silence is worth more than a thousand words.
ℹtem \textbf{Context Keys:} To read an old text or recall a forbidden history, you need three keys: place, season, and witness. Without all three, roll at disadvantage. A Lethai−thora may create a temporary Context Key by spending 1 Boon and describing a small ritual (lighting incense, touching a stone, speaking a name).
ℹtem \textbf{The Keeper's Toll:} When using a Lethai−thora Rite or consulting a hidden archive, mark 1 Fatigue. The knowledge costs something — a moment of warmth, a pleasant memory, the scent of rain. Some keepers have forgotten how to smile.
ℹtem \textbf{Local Patrons (Syncretic Names):}
  \begin{itemize}
    ℹtem \textbf{The Unwritten} — The Witness. \textit{A secret that never needs to be spoken.}
    ℹtem \textbf{The Quiet Scale} — Mykkiel. \textit{The law that weighs silence as heavily as testimony.}
    ℹtem \textbf{The Ink Father} — Maelstraeus. \textit{The debt written in disappearing ink; the interest that accrues in forgetting.}
    ℹtem \textbf{The Echo of the First Pavilion} — The Traveler. \textit{The path that leads only inward.}
  \end{itemize}
\end{itemize}

\subsection*{Lore Echoes (Cross−Expansion Hooks)}
\begin{itemize}
  ℹtem \textbf{The Ninth Scroll and the Kalrantha Amulet (Dhahara):} The Ninth Scroll is said to describe the creation of the amulet — and the demon sealed within. The Veiled Sisters have been searching for it for centuries. The Lethai−thora have hidden it in plain sight.
  ℹtem \textbf{Silence Tokens and the Covenant (Mykkiel):} The Covenant recognises a Lethai−thora silence token as a valid request for a closed hearing. A silence token can pause any court proceeding for the duration of a single breath.
  ℹtem \textbf{Context Keys and the Ukwe (Sekogo):} The Ukwe's roots are a natural context key — they remember the season, the place, and the witness. A Lethai−thora who consults the Ukwe does not need to spend a Boon.
  ℹtem \textbf{Memory−Threads and the Hollow (Eastern Realms):} The Hollow feeds on forgotten memories. A Lethai−thora's memory−thread is a feast. The Oath−Keepers will cut a thread if given the chance.
  ℹtem \textbf{The Fledging of the Empire (Sihai):} The Lethai−thora's oldest scrolls still bear the water stains of the Valewood crossing. A keeper who unrolls one of these scrolls may glimpse a memory of the old fey courts — but the glimpse costs a year of her life.
\end{itemize}

\subsection*{Lethai−thora Superstitions (burn for +1 die once)}
\begin{itemize}
  ℹtem \textbf{Bread to the Archive:} Crumble a crust at the threshold. Ignore the first \textit{Archive Fire} penalty. \texttt{[SUPERSTITION]}
  ℹtem \textbf{Knot the Scroll:} Tie a single knot in a silk cord. Cancel \textit{Forged Scroll} or take –1 die on your next forgery detection roll (your choice). \texttt{[SUPERSTITION]}
  ℹtem \textbf{Name to the Well:} Whisper a dead keeper's name into the well. Gain +1 Effect when recalling a forgotten fact this scene, but mark \textit{Obligation} to the Lethai−thora. \texttt{[SUPERSTITION]}
\end{itemize}

\subsection*{Thematic SB Spend Table}
łabel{sec:lethai−thora−sb}

\paragraph{Minor Complications (1 SB)}
\begin{itemize}
ℹtem \textbf{Exposure:} Your actions draw unwanted attention from \textbf{Censorate spies or imperial guards}.
ℹtem \textbf{Noise:} Sounds of your actions alert nearby \textbf{scribes or keepers}.
ℹtem \textbf{Trace:} Evidence of your passage marks your route for \textbf{memory−hunters or the Hollow}.
ℹtem \textbf{Delay:} A brief but meaningful setback costs you \textbf{time or access to an archive}.
ℹtem \textbf{Supply Strain:} Mark +1 segment on a relevant \textbf{resource timer}.
\end{itemize}

\paragraph{Moderate Setbacks (2 SB)}
\begin{itemize}
ℹtem \textbf{Alarm Raised:} \textbf{Local keeper or Censorate agent} becomes aware and begins responding.
ℹtem \textbf{Position Lost:} You lose advantageous ground/cover/stealth due to \textbf{flickering lamp or silent alarm}.
ℹtem \textbf{Foe Appears:} A \textbf{Censorate interrogator or a Oath−Keeper} arrives on scene.
ℹtem \textbf{Gear Trouble:} A piece of equipment becomes \textbf{Compromised/Neglected}.
ℹtem \textbf{Lock/Barrier:} A simple obstacle now requires a test to overcome.
\end{itemize}

\paragraph{Serious Trouble (3 SB)}
\begin{itemize}
ℹtem \textbf{Reinforcements:} Additional \textbf{Censorate agents, imperial guards, or Hollow spirits} arrive.
ℹtem \textbf{Key Gear Breaks:} A crucial tool/weapon becomes temporarily unusable.
ℹtem \textbf{Major Twist:} The situation fundamentally changes—\textbf{scroll burns/thread cut/pavilion appears}.
ℹtem \textbf{Rail Tick:} Advance a relevant campaign/front timer by 1 segment.
ℹtem \textbf{Condition Applied:} Mark \textbf{Fatigue 1/Harm 1/Condition} appropriate to fiction.
\end{itemize}

\paragraph{Major Turns (4+ SB)}
\begin{itemize}
ℹtem \textbf{Trap Springs:} A prepared danger activates with full effect.
ℹtem \textbf{Authority Arrival:} \textbf{High Keeper, Quiet Chancellor, or Ghost of the Ninth Scroll} intervenes.
ℹtem \textbf{Scene Shift:} The environment changes dramatically—\textbf{archive floods/well overflows/first pavilion materialises}.
ℹtem \textbf{Patron Omen:} Divine/arcane forces take notice—\textbf{omen appears/blessing lost/curse manifests}.
ℹtem \textbf{Narrative Pivot:} The story takes an unexpected turn that reframes objectives.
\end{itemize}

\paragraph{Region−Specific SB Options}
\begin{itemize}
ℹtem \textbf{Lethai−thora (Silence Law):} Lamps flicker without breeze, scrolls shuffle their own order, a whisper is heard from a sealed jar.
ℹtem \textbf{Lethai−thora (Memory Law):} A keeper forgets your face, a memory−thread glows, you recall a fact you never learned.
ℹtem \textbf{Lethai−thora (Context Law):} The season changes inexplicably, a stone context key warms, a witness appears who was not there.
\end{itemize}

\subsection*{Pursuit Ladder (Near / Mid / Far)}
\begin{itemize}
ℹtem \textbf{Advance/Withdraw (Wits + Stealth):} On a hit, shift one band. On a strong hit, also impose \textit{Silent Step}.
ℹtem \textbf{Cut the Memory (Presence + Sway):} On a hit, freeze range for one exchange; if a memory−thread is in your possession, +1 die.
ℹtem \textbf{Weather the Interrogation (Resolve + Spirit):} On a hit in \textit{Censorate's Interrogation}, you pick who is questioned; on a miss, the GM does.
\end{itemize}

\subsection*{Boss Seeds (Ink & Silence)}
\begin{itemize}
ℹtem \textbf{The Scroll−Eater (Nine Forgeries)} — \textit{Phases}: Scrolls Burned \(→\) Orchestrated Erasure \(→\) Archive Coup. \textit{Cracks}: recover a true scroll, expose a false context key, survive a night in the well.
ℹtem \textbf{The Forgotten God (Memory Eater)} — \textit{Phases}: Voices from the Shrine \(→\) Procession of Blank Scrolls \(→\) Ink−Baptism of a Keeper. \textit{Cracks}: relic memory−thread, survivor testimony, the first pavilion opening its doors.
\end{itemize}

\begin{tcolorbox}[colback=black!3,colframe=black!40!white,title={Patronage & Power}]
Among the Lethai−thora, power flows through silence, ink, and the weight of a forgotten secret. A Keeper's authority rests on how many secrets she has buried and how many she has forgotten. The archive is the fortress; the silence is the shield. Yet they are also entangled with the empire they helped create — the Censorate watches them, the Emperor needs them, and the common people fear them. The Lethai−thora can no longer leave Sihai, even if they wished. Their fate is bound to the Mandate of Heaven.

\textbf{For the GM:}  
Power in Lethai−thora society is measured in scrolls sealed, context keys held, and the length of a pause. Patronage often takes the form of forbidden scrolls, context stones, or memory−threads — each tied to a specific secret, archive, or keeper. To emphasise the entangled relationship with Sihai:
\begin{itemize}
ℹtem Use the \textbf{Silence as Sanctuary} mechanic to create tense social scenes where a pause is a weapon.
ℹtem Require \textbf{Context Keys} to unlock ancient lore — the party must gather place, season, and witness before a scroll can be read.
ℹtem Let the \textbf{Keeper's Toll} be a meaningful cost: fatigue, lost memories, or forgotten faces.
ℹtem Remind players that the Lethai−thora are not outsiders; they are the empire's conscience. A Lethai−thora NPC may hold a secret about a PC's family, or about the true cause of a war, or about a crime the emperor himself committed.
\end{itemize}
In Lethai−thora society, the ink forgets nothing — and the silence remembers everything.
\end{tcolorbox}

% Index entries
∈dex{Lethai−thora|see{Silent Ledgers}}
∈dex{Theme generator}
∈dex{Silent Ledgers generator}
∈dex{Places!Lethai−thora}
∈dex{People!Lethai−thora}
∈dex{Complications!Lethai−thora}
∈dex{Rewards!Lethai−thora}
∈dex{Forbidden scrolls}
∈dex{Context keys}
∈dex{Memory−threads}
∈dex{Silence as sanctuary}
∈dex{Keeper's toll}
∈dex{Quiet Chancellor}
∈dex{Censorate}
∈dex{First Pavilion}
∈dex{Fledging of the Empire}
∈dex{Local Patrons!Lethai−thora}
\newpage